% Grands mouvements de l'atmosphère et de l'hydrosphère — SVT 1ère TI — Cours signé OnBuch+
% Programme officiel MINESEC. Compiler avec : tectonic NCT07-mouvements-atmosphere-hydrosphere.tex
\newcommand{\DOCMATIERE}{SVT}
\newcommand{\DOCNIVEAU}{1ère TI}
\newcommand{\DOCMODULE}{Module 3 : L'Éducation à l'Environnement et au Développement Durable}
\newcommand{\DOCLECON}{NCT07}
\newcommand{\DOCTITRE}{Grands mouvements de l'atmosphère et de l'hydrosphère}
% OnBuch+ — préambule « cours signés » ENRICHI (compilé avec tectonic / XeLaTeX)
% Système visuel « Pop » d'OnBuch+ : fond crème (jamais blanc), bordures encre
% épaisses, ombres « dures » décalées, titres Archivo Black en capitales.
% Version MATHÉMATIQUES : graphes (pgfplots/tikz), boîtes pédagogiques
% (définition, propriété, méthode, attention, le savais-tu, exercice résolu,
% exercice, TP, bilan), page de garde et signature OnBuch+ ancrée en bas.
%
% Macros attendues dans main.tex AVANT \input{preamble} :
%   \DOCMATIERE  \DOCNIVEAU  \DOCMODULE  \DOCLECON (numéro)  \DOCTITRE
\documentclass[11pt]{article}

\usepackage{fontspec}
\usepackage[french]{babel}
\usepackage[a4paper,margin=2cm,top=3.4cm,bottom=2.8cm,headheight=1.4cm,footskip=1.3cm]{geometry}
\usepackage{amsmath,amssymb,amsfonts}
\usepackage{mathtools} % \xmapsto, \coloneqq... souvent utilisés par les modèles
\usepackage{cancel} % simplifications barrées (bilans de forces)
% \abs{z} (module, valeur absolue) et \norm{u} (norme), utilisés par les modèles
\providecommand{\abs}{}\let\abs\relax\DeclarePairedDelimiter{\abs}{\lvert}{\rvert}
\providecommand{\norm}{}\let\norm\relax\DeclarePairedDelimiter{\norm}{\lVert}{\rVert}
\usepackage[table]{xcolor} % [table] : fournit aussi \rowcolors (alternance de lignes)
\usepackage{pagecolor}
\usepackage{tikz}
\usetikzlibrary{arrows.meta,positioning,calc,shapes.geometric,shapes.misc,shapes.arrows,decorations.pathmorphing,decorations.pathreplacing,decorations.markings,patterns,shadows,mindmap,trees,fit,backgrounds,matrix,intersections,angles,quotes}
\usepackage{pgfplots}
\usepackage[version=4]{mhchem} % équations nucléaires \ce{^{235}_{92}U}
\usepackage[european,siunitx]{circuitikz} % schémas électriques
\pgfplotsset{compat=1.18}
\usepgfplotslibrary{fillbetween,groupplots} % aires entre courbes (calcul intégral)
\usepackage{enumitem}
\usepackage{fancyhdr}
% [most] charge toutes les bibliothèques tcolorbox (listings, minted, raster,
% documentation...) et gonfle inutilement la mémoire TeX sur un document long
% (~300 boîtes) : on ne charge que ce qui est réellement utilisé.
\usepackage[skins,breakable]{tcolorbox}
\usepackage{graphicx}
\usepackage{colortbl}
\usepackage{booktabs}
\usepackage{tabularx}
\usepackage{diagbox} % case d'en-tête diagonale des tableaux à double entrée
% Colonnes extensibles centrée / gauche / droite (C, L, R), souvent utilisées par les modèles
\newcolumntype{C}{>{\centering\arraybackslash}X}
\newcolumntype{L}{>{\raggedright\arraybackslash}X}
\newcolumntype{R}{>{\raggedleft\arraybackslash}X}
\usepackage{array}
\usepackage{multicol}
\usepackage{siunitx}
% parse-numbers=false : accepte aussi \SI{2,5 \times 10^{-3}}{...}
\sisetup{locale=FR,output-decimal-marker={,},group-minimum-digits=4,parse-numbers=false}
\DeclareSIUnit\ohm{\ensuremath{\symup{\Omega}}} % U+2126 absent de la police
\DeclareSIUnit\division{div} % réglages d'oscilloscope (ms/div, V/div)
\DeclareSIUnit\tr{tr} % tours (vitesse de rotation, ex. \SI{1500}{\tr\per\minute})
\DeclareSIUnit\ch{ch} % cheval-vapeur (puissance moteur, unité usuelle hors SI)
\DeclareSIUnit\franccfa{FCFA} % franc CFA (coûts, contexte camerounais)
\DeclareSIUnit\dioptre{\ensuremath{\delta}} % vergence d'une lentille (dioptrie)
\usepackage{qrcode}
% \degree : commande fréquemment utilisée par les modèles pour un angle ou une
% température en dehors de \SI{}{} (non fournie par les paquets chargés).
\providecommand{\degree}{^{\circ}}
% \qed (fin de démonstration), utilisé par les modèles sans amsthm
\providecommand{\qed}{\hfill$\square$}
% \barre : double barre des valeurs interdites dans un tableau de variations (tkz-tab)
\providecommand{\barre}{\ensuremath{\|}}
% environnement proof (amsthm non chargé) utilisé par les modèles
\ifdefined\proof\else\newenvironment{proof}[1][Démonstration]{\par\noindent\textit{#1.}\ }{\qed\par}\fi
\usepackage{lastpage}
\usepackage{hyperref}
\hypersetup{hidelinks}

% ============================================================
% Palette « Pop » OnBuch+ — mêmes hex que la classe P (plus_ui.dart)
% ============================================================
\definecolor{poporange}{HTML}{F59321}
\definecolor{popdark}{HTML}{E07A0C}
\definecolor{popcream}{HTML}{FFF6E8}
\definecolor{popink}{HTML}{1C1714}
\definecolor{poppurple}{HTML}{7A5AE0}
\definecolor{popgreen}{HTML}{2E9E6A}
\definecolor{poppink}{HTML}{E0457B}
\definecolor{popblue}{HTML}{2F7DE1}
\definecolor{popgold}{HTML}{C98A12}
\definecolor{popmuted}{HTML}{8A7F76}
\definecolor{popline}{HTML}{EADFD2}
\definecolor{popgray}{HTML}{8A7F76}
\colorlet{popgrey}{popgray}
% Alias tolérants (le modèle nomme parfois les couleurs différemment)
\colorlet{popwhite}{white}
\colorlet{popblack}{popink}
\colorlet{popred}{poppink}
\colorlet{popyellow}{popgold}
% Teintes claires pour fonds de tableaux / zones
\colorlet{popblueL}{popblue!10!white}
\colorlet{poporangeL}{poporange!14!white}
\colorlet{popgreenL}{popgreen!12!white}
\colorlet{poppurpleL}{poppurple!12!white}
\colorlet{poppinkL}{poppink!12!white}
\colorlet{popgoldL}{popgold!14!white}

\pagecolor{popcream}

% ============================================================
% Typographie
% ============================================================
\defaultfontfeatures{Ligatures=TeX}
\setmainfont{PlusJakartaSans-Medium.ttf}[Path=../../../fonts/,BoldFont=PlusJakartaSans-Bold.ttf]
% Maths en sans-serif (Fira Math) avec chiffres et lettres droites de Plus
% Jakarta, pour que les formules se fondent dans le texte.
\usepackage[math-style=ISO,bold-style=ISO]{unicode-math}
\setmathfont{FiraMath-Regular.otf}[Path=../../../fonts/]
\setmathfont{PlusJakartaSans-Medium.ttf}[Path=../../../fonts/,range={up/num,up/{Latin,latin},\mathup}]
\usepackage{icomma} % virgule décimale sans espace en mode math
% Caractères mathématiques Unicode absents des polices de texte (Plus Jakarta,
% Archivo Black) : rendus par leur équivalent mathématique, en texte comme en
% formule — sinon ils s'affichent en carré vide (ex. « Arithmétique dans ℤ »).
\usepackage{newunicodechar}
\newunicodechar{ℝ}{\ensuremath{\mathbb{R}}}
\newunicodechar{ℤ}{\ensuremath{\mathbb{Z}}}
\newunicodechar{ℂ}{\ensuremath{\mathbb{C}}}
\newunicodechar{ℕ}{\ensuremath{\mathbb{N}}}
\newunicodechar{ℚ}{\ensuremath{\mathbb{Q}}}
\newunicodechar{𝔻}{\ensuremath{\mathbb{D}}}
\newunicodechar{α}{\ensuremath{\alpha}}
\newunicodechar{β}{\ensuremath{\beta}}
\newunicodechar{γ}{\ensuremath{\gamma}}
\newunicodechar{δ}{\ensuremath{\delta}}
\newunicodechar{ε}{\ensuremath{\varepsilon}}
\newunicodechar{θ}{\ensuremath{\theta}}
\newunicodechar{λ}{\ensuremath{\lambda}}
\newunicodechar{μ}{\ensuremath{\mu}}
\newunicodechar{σ}{\ensuremath{\sigma}}
\newunicodechar{τ}{\ensuremath{\tau}}
\newunicodechar{φ}{\ensuremath{\varphi}}
\newunicodechar{ψ}{\ensuremath{\psi}}
\newunicodechar{ω}{\ensuremath{\omega}}
\newunicodechar{Σ}{\ensuremath{\Sigma}}
% majuscules produites par \MakeUppercase dans les titres (α-aminés -> Α)
\newunicodechar{Α}{\ensuremath{\alpha}}
\newunicodechar{Β}{\ensuremath{\beta}}
\newunicodechar{Γ}{\ensuremath{\Gamma}}
\newunicodechar{Δ}{\ensuremath{\Delta}}
\newunicodechar{Ε}{\ensuremath{\varepsilon}}
\newunicodechar{Θ}{\ensuremath{\Theta}}
\newunicodechar{Λ}{\ensuremath{\Lambda}}
\newunicodechar{Μ}{\ensuremath{\mu}}
\newunicodechar{Τ}{\ensuremath{\tau}}
\newunicodechar{Φ}{\ensuremath{\Phi}}
\newunicodechar{Ψ}{\ensuremath{\Psi}}
\newunicodechar{Ω}{\ensuremath{\Omega}}
\newunicodechar{↦}{\ensuremath{\mapsto}}
\newunicodechar{∈}{\ensuremath{\in}}
\newunicodechar{∉}{\ensuremath{\notin}}
\newunicodechar{∧}{\ensuremath{\wedge}}
\newunicodechar{⇔}{\ensuremath{\Leftrightarrow}}
\newunicodechar{⟺}{\ensuremath{\Longleftrightarrow}}
\newunicodechar{⇒}{\ensuremath{\Rightarrow}}
\newunicodechar{⟹}{\ensuremath{\Longrightarrow}}
\newunicodechar{∘}{\ensuremath{\circ}}
\newunicodechar{∪}{\ensuremath{\cup}}
\newunicodechar{∩}{\ensuremath{\cap}}
\newunicodechar{≡}{\ensuremath{\equiv}}
\newunicodechar{⊂}{\ensuremath{\subset}}
\newunicodechar{⊥}{\ensuremath{\perp}}
\newunicodechar{∃}{\ensuremath{\exists}}
\newunicodechar{∀}{\ensuremath{\forall}}
\newunicodechar{∛}{\ensuremath{\sqrt[3]{\,}}}
\newunicodechar{∜}{\ensuremath{\sqrt[4]{\,}}}
\newunicodechar{⃗}{\ensuremath{{}^{\to}}}
\newunicodechar{ⁿ}{\ifmmode^{n}\else\textsuperscript{\ensuremath{n}}\fi}
\newunicodechar{ˣ}{\ifmmode^{x}\else\textsuperscript{\ensuremath{x}}\fi}
\newunicodechar{ᵇ}{\ifmmode^{b}\else\textsuperscript{\ensuremath{b}}\fi}
\newunicodechar{ʳ}{\ifmmode^{r}\else\textsuperscript{\ensuremath{r}}\fi}
\newunicodechar{ᵘ}{\ifmmode^{u}\else\textsuperscript{\ensuremath{u}}\fi}
\newunicodechar{ʸ}{\ifmmode^{y}\else\textsuperscript{\ensuremath{y}}\fi}
\newunicodechar{ᵃ}{\ifmmode^{a}\else\textsuperscript{\ensuremath{a}}\fi}
\newunicodechar{ᵉ}{\ifmmode^{e}\else\textsuperscript{\ensuremath{e}}\fi}
\newunicodechar{ᵏ}{\ifmmode^{k}\else\textsuperscript{\ensuremath{k}}\fi}
\newunicodechar{ᵖ}{\ifmmode^{p}\else\textsuperscript{\ensuremath{p}}\fi}
\newunicodechar{ᴺ}{\ifmmode^{N}\else\textsuperscript{\ensuremath{N}}\fi}
\newunicodechar{ᶜ}{\ifmmode^{c}\else\textsuperscript{\ensuremath{c}}\fi}
\newunicodechar{ʰ}{\ifmmode^{h}\else\textsuperscript{\ensuremath{h}}\fi}
\newunicodechar{ᵠ}{\ifmmode^{q}\else\textsuperscript{\ensuremath{q}}\fi}
\newunicodechar{ₙ}{\ifmmode_{n}\else\textsubscript{\ensuremath{n}}\fi}
\newunicodechar{ₐ}{\ifmmode_{a}\else\textsubscript{\ensuremath{a}}\fi}
\newunicodechar{ᵢ}{\ifmmode_{i}\else\textsubscript{\ensuremath{i}}\fi}
\newunicodechar{ᵣ}{\ifmmode_{r}\else\textsubscript{\ensuremath{r}}\fi}
\newunicodechar{ₖ}{\ifmmode_{k}\else\textsubscript{\ensuremath{k}}\fi}
\newunicodechar{ᵦ}{\ifmmode_{\beta}\else\textsubscript{\ensuremath{\beta}}\fi}
\newunicodechar{ₓ}{\ifmmode_{x}\else\textsubscript{\ensuremath{x}}\fi}
\newfontfamily\popdisplay{ArchivoBlack-Regular.ttf}[Path=../../../fonts/]
\newfontfamily\poplogo{SpaceGrotesk-Bold.ttf}[Path=../../../fonts/]
\newfontfamily\popsemi{PlusJakartaSans-SemiBold.ttf}[Path=../../../fonts/]
\newfontfamily\popxbold{PlusJakartaSans-ExtraBold.ttf}[Path=../../../fonts/]
\newfontfamily\popxboldls{PlusJakartaSans-ExtraBold.ttf}[Path=../../../fonts/,LetterSpace=3.0]

\setlist[enumerate]{leftmargin=1.7em,itemsep=5pt,topsep=4pt}
\setlist[itemize]{leftmargin=1.7em,itemsep=4pt,topsep=4pt,label={\color{poporange}\textbullet}}
\setlength{\parindent}{0pt}
\setlength{\parskip}{5pt}
\renewcommand{\arraystretch}{1.25}

% pgfplots : style Pop par défaut
\pgfplotsset{
  every axis/.append style={axis line style={popink,line width=1pt},tick style={popink},
    grid style={popline,line width=0.5pt},label style={font=\small},tick label style={font=\footnotesize}},
  popcurve/.style={poporange,line width=1.8pt,no markers,smooth},
}
% Même style disponible dans un \draw TikZ simple (hors pgfplots)
\tikzset{popcurve/.style={poporange,line width=1.8pt}}
\tikzset{popgrid/.style={popline,very thin}}
\colorlet{popcurve}{poporange} % le nom du style est parfois utilisé comme couleur

% ============================================================
% Pill / squiggle
% ============================================================
\newcommand{\poppill}[3]{%
  \tikz[baseline=-0.55ex]\node[draw=popink,line width=1.2pt,rounded corners=2.6pt,%
    fill=#2,inner xsep=6.5pt,inner ysep=3pt]{%
    \popxboldls\fontsize{7}{7}\selectfont\color{#3}\MakeUppercase{#1}};%
}
\newcommand{\popsquiggle}[1]{%
  \tikz[baseline=-0.4ex]{
    \draw[#1,line width=2.6pt,line cap=round]
      plot [smooth] coordinates {(0,0.09) (0.34,0.01) (0.68,0.21) (1.02,0.01) (1.36,0.10)};
  }%
}
% Mot-clé surligné (marqueur orange sous le texte)
\newcommand{\cle}[1]{{\popxbold\color{popdark}#1}}

% ============================================================
% En-tête / pied
% ============================================================
\pagestyle{fancy}
\fancyhf{}
\renewcommand{\headrulewidth}{0pt}
\renewcommand{\footrulewidth}{0pt}
\lhead{\raisebox{-0.15ex}{\poplogo\fontsize{13}{13}\selectfont\color{popink}OnBuch\color{poporange}+}}
\rhead{\poppill{\DOCMATIERE\ \textbullet\ \DOCNIVEAU\ \textbullet\ Leçon \DOCLECON}{white}{popink}}
\lfoot{\popxbold\fontsize{7.6}{7.6}\selectfont\color{popmuted}\MakeUppercase{Cours signé OnBuch+}\ \textbullet\ {\popsemi\DOCTITRE}}
\rfoot{\tikz[baseline=-0.6ex]\node[rounded corners=8.5pt,fill=popink,minimum height=17pt,minimum width=17pt,inner xsep=5pt,inner sep=0]{\popxbold\fontsize{8}{8}\selectfont\color{popcream}\thepage\,/\,\pageref*{LastPage}};}

% ============================================================
% Titres
% ============================================================
\newcommand{\chaptitre}[2]{%
  \begin{tcolorbox}[enhanced,colback=poporange,colframe=popink,boxrule=2.6pt,arc=11pt,
      drop shadow=popink,left=15pt,right=15pt,top=12pt,bottom=11pt,before skip=3pt,after skip=18pt]
    \poppill{#2}{popink}{popcream}\par\vspace{9pt}
    {\raggedright\hyphenpenalty=10000\exhyphenpenalty=10000\popdisplay\fontsize{22}{24}\selectfont\color{popcream}\MakeUppercase{#1}\par}
  \end{tcolorbox}
}
% Section numérotée : pastille orange avec le numéro + titre Archivo
\newcounter{popsec}
\newcommand{\coursec}[1]{%
  \stepcounter{popsec}\setcounter{popsub}{0}%
  \par\vspace{16pt}\noindent
  \begin{minipage}{\linewidth}
  \tikz[baseline=-0.7ex]\node[draw=popink,line width=1.6pt,fill=poporange,rounded corners=4pt,minimum size=22pt,inner sep=0]{\popdisplay\fontsize{12}{12}\selectfont\color{popcream}\thepopsec};%
  \hspace{8pt}{\popdisplay\fontsize{15}{17}\selectfont\color{popink}\MakeUppercase{#1}}\par
  \vspace{4pt}\noindent{\color{poporange}\rule{36pt}{3.6pt}}
  \end{minipage}\par\nopagebreak\vspace{8pt}
}
\newcounter{popsub}
\newcommand{\courssub}[1]{%
  \stepcounter{popsub}%
  \par\vspace{9pt}\noindent{\popxbold\fontsize{11.5}{13}\selectfont\color{popink}{\color{poporange}\thepopsec.\thepopsub}\hspace{6pt}#1}\par\nopagebreak\vspace{4pt}
}

% ============================================================
% Boîtes pédagogiques — même recette : fond blanc, bordure épaisse colorée,
% ombre dure encre, badge plein. Titre optionnel [#1].
% ============================================================
\tcbset{popbase/.style={breakable,enhanced,colback=white,boxrule=2.3pt,arc=9pt,
  shadow={3pt}{-3pt}{0pt}{popink},left=14pt,right=14pt,top=11pt,bottom=11pt,before skip=11pt,after skip=15pt,
  fontupper=\popsemi\fontsize{10.6}{14.5}\selectfont\color{popink},
  fontlower=\popsemi\fontsize{10.6}{14.5}\selectfont\color{popink}}}
% Coche dessinée (le glyphe ✓ n'existe pas dans les polices du document)
\newcommand{\popcheck}[1]{\tikz[baseline=-0.5ex]\draw[#1,line width=1.6pt,line cap=round,line join=round](-0.13,0.0)--(-0.04,-0.09)--(0.13,0.1);}
\providecommand{\coche}{\popcheck{popgreen}}
\providecommand{\resultat}[1]{\par\smallskip\noindent\fcolorbox{popgreen}{popgreenL}{\parbox{\dimexpr\linewidth-2\fboxsep-2\fboxrule\relax}{\textbf{Résultat.} #1}}\par\smallskip}
\newcommand{\popboxhead}[3]{\poppill{#1}{#2}{white}\ifx\relax#3\relax\else\hspace{6pt}{\popxbold\fontsize{10.6}{12}\selectfont\color{popink}#3}\fi\par\vspace{7pt}}

\newtcolorbox{aretenir}[1][]{popbase,colframe=popgreen,before upper={\popboxhead{À retenir}{popgreen}{#1}}}
\newtcolorbox{exemplebox}[1][]{popbase,colframe=poporange,before upper={\popboxhead{Exemple}{poporange}{#1}}}
\newtcolorbox{definition}[1][]{popbase,colframe=popblue,before upper={\popboxhead{Définition}{popblue}{#1}}}
\newtcolorbox{propriete}[1][]{popbase,colframe=poppurple,before upper={\popboxhead{Propriété}{poppurple}{#1}}}
\newtcolorbox{methode}[1][]{popbase,colframe=popgold,before upper={\popboxhead{Méthode}{popgold}{#1}}}
\newtcolorbox{attention}[1][]{popbase,colframe=poppink,before upper={\popboxhead{Attention}{poppink}{#1}}}
\newtcolorbox{experience}[1][]{popbase,colframe=popblue,colback=popblueL,before upper={\popboxhead{Expérience}{popblue}{#1}}}
% « Le savais-tu ? » : carte violette pleine (contexte, culture, Cameroun)
\newtcolorbox{savaistu}[1][]{popbase,colback=poppurple,colframe=popink,
  fontupper=\popsemi\fontsize{10.4}{14.2}\selectfont\color{white},
  before upper={\poppill{Le savais-tu\,?}{white}{poppurple}\ifx\relax#1\relax\else\hspace{6pt}{\popxbold\color{white}#1}\fi\par\vspace{7pt}}}
% Exercice résolu : énoncé en haut, \tcblower puis la solution sur fond crème
\newcounter{popexo}\newcounter{popexr}
\newtcolorbox{exoresolu}[1][]{popbase,colframe=poporange,colbacklower=popcream,
  segmentation style={popink,line width=1.2pt,dashed},
  before upper={\stepcounter{popexr}\popboxhead{Exercice résolu \thepopexr}{poporange}{#1}},
  before lower={\poppill{Solution}{popink}{popcream}\par\vspace{6pt}}}
% Exercice (non résolu en place) — la correction est dans la partie Corrigés
\newtcolorbox{exercice}[1][]{popbase,colframe=popink,
  before upper={\stepcounter{popexo}\popboxhead{Exercice \thepopexo}{popink}{#1}}}
% Corrigé
\newtcolorbox{corrige}[1][]{popbase,colframe=popgreen,colback=popgreenL,
  before upper={\popboxhead{Corrigé}{popgreen}{#1}}}
% Objectifs / prérequis / bilan
\newtcolorbox{objectifs}{popbase,colframe=popink,colback=poporangeL,
  before upper={\popboxhead{Objectifs de la leçon}{popink}{}}}
\newtcolorbox{prerequis}{popbase,colframe=popink,colback=popblueL,
  before upper={\popboxhead{Prérequis}{popblue}{}}}
\newtcolorbox{bilan}{popbase,colframe=popink,colback=white,boxrule=2.8pt,
  before upper={\popboxhead{Fiche bilan}{poporange}{L'essentiel à connaître pour l'examen}}}
% Figure encadrée avec légende
\newtcolorbox{popfigure}[1][]{enhanced,breakable=false,colback=white,colframe=popline,boxrule=1.4pt,arc=8pt,
  left=10pt,right=10pt,top=10pt,bottom=8pt,before skip=10pt,after skip=12pt,halign=center,
  lower separated=false,halign lower=center,
  fontlower=\popsemi\fontsize{9}{11.5}\selectfont\color{popmuted},
  #1}
\newcommand{\legende}[1]{\par\vspace{6pt}{\popsemi\fontsize{9}{11.5}\selectfont\color{popmuted}#1}}

% ============================================================
% Signature OnBuch+
% ============================================================
\newcommand{\poplogoinline}[1]{{\poplogo\fontsize{#1}{#1}\selectfont\color{popink}OnBuch\color{poporange}+}}
\newcommand{\signatureonbuch}{%
  \begin{tcolorbox}[enhanced,colback=popink,colframe=popink,boxrule=2.2pt,arc=10pt,
      drop shadow=poporange,left=14pt,right=14pt,top=10pt,bottom=10pt,before skip=0pt,after skip=0pt]
    \tikz[baseline=-0.7ex]\node[circle,fill=popgreen,draw=popcream,line width=1.3pt,minimum size=22pt,inner sep=0]{\popcheck{white}};%
    \hspace{9pt}\begin{minipage}[c]{0.62\linewidth}
      {\popxbold\fontsize{12}{13}\selectfont\color{popcream}Signé\ \textbullet\ {\poplogo OnBuch\color{poporange}+}}\par\vspace{2pt}
      {\raggedright\popsemi\fontsize{8}{10.5}\selectfont\color{popline}\MakeUppercase{Équipe pédagogique OnBuch+}\\\MakeUppercase{Conforme au programme officiel MINESEC}\par}
    \end{minipage}\hfill
    {\poplogo\fontsize{22}{22}\selectfont\color{popcream}OnBuch\color{poporange}+}
  \end{tcolorbox}%
}

% ============================================================
% Page de garde
% #1 titre · #2 matière · #3 niveau · #4 module · #5 n° leçon · #6 durée ·
% #7 description / objectifs (texte ou itemize)
% ============================================================
\newcommand{\pagegarde}[7]{%
  \thispagestyle{empty}
  \newgeometry{a4paper,margin=1.7cm,top=1.6cm,bottom=1.4cm}%
  \begin{tikzpicture}[remember picture,overlay]
    \fill[poporange!18] ([xshift=-3.2cm,yshift=-3.6cm]current page.north east) circle (4.6cm);
    \fill[poppurple!14] ([xshift=2.4cm,yshift=7.6cm]current page.south west) circle (3.8cm);
  \end{tikzpicture}%
  {\poplogo\fontsize{30}{30}\selectfont\color{popink}OnBuch\color{poporange}+}\hfill
  \raisebox{0.6ex}{\poppill{Cours signé \textbullet\ Premium}{popink}{popcream}}\par
  \vspace{1pt}\popsquiggle{poppurple}\par
  \vspace{8pt}
  \begin{tcolorbox}[enhanced,colback=poporange,colframe=popink,boxrule=2.8pt,arc=13pt,
      drop shadow=popink,left=18pt,right=18pt,top=12pt,bottom=13pt,after skip=10pt]
    \poppill{#2 \textbullet\ #3}{popink}{popcream}\hspace{5pt}\poppill{#4}{white}{popink}\par\vspace{8pt}
    {\popdisplay\fontsize{14}{14}\selectfont\color{popink}LEÇON #5}\par\vspace{4pt}
    {\raggedright\hyphenpenalty=10000\exhyphenpenalty=10000\popdisplay\fontsize{25}{27}\selectfont\color{popcream}\MakeUppercase{#1}\par}
    \vspace{8pt}
    {\popxbold\fontsize{9.5}{9.5}\selectfont\color{popink}\MakeUppercase{Durée conseillée : #6}}
  \end{tcolorbox}
  \begin{tcolorbox}[enhanced,colback=white,colframe=popink,boxrule=2.2pt,arc=10pt,
      drop shadow=popink,left=16pt,right=16pt,top=10pt,bottom=10pt,after skip=8pt]
    \poppill{Dans cette leçon}{poppurple}{white}\par\vspace{5pt}
    \popsemi\fontsize{9.3}{12.6}\selectfont\color{popink}#7
  \end{tcolorbox}
  \noindent\begin{minipage}[t]{0.487\linewidth}
    \begin{tcolorbox}[enhanced,colback=popgreen,colframe=popink,boxrule=2.2pt,arc=10pt,
        drop shadow=popink,left=11pt,right=11pt,top=10pt,bottom=10pt,height=3.1cm,valign=center]
      \tcbox[colback=white,colframe=popink,boxrule=1.3pt,arc=5pt,boxsep=0pt,nobeforeafter,
        left=3pt,right=3pt,top=3pt,bottom=3pt,valign=center]{\qrcode[height=1.9cm]{https://play.google.com/store/apps/details?id=com.luvvix.onbuch&hl=fr}}%
      \hspace{8pt}\begin{minipage}[c]{0.5\linewidth}
        {\popdisplay\fontsize{11}{12.5}\selectfont\color{white}\MakeUppercase{Télécharge OnBuch}}\par\vspace{4pt}
        {\raggedright\popsemi\fontsize{8.4}{11}\selectfont\color{white}Gratuit \textbullet\ Google Play\\Tuteur IA, annales, concours, résultats\par}
      \end{minipage}
    \end{tcolorbox}
  \end{minipage}\hfill
  \begin{minipage}[t]{0.487\linewidth}
    \begin{tcolorbox}[enhanced,colback=poppurple,colframe=popink,boxrule=2.2pt,arc=10pt,
        drop shadow=popink,left=11pt,right=11pt,top=10pt,bottom=10pt,height=3.1cm,valign=center]
      \tikz[baseline=-0.7ex]\node[circle,fill=white,draw=popink,line width=1.3pt,minimum size=1.6cm,inner sep=0]{\popdisplay\fontsize{24}{24}\selectfont\color{poppurple}+};%
      \hspace{8pt}\begin{minipage}[c]{0.6\linewidth}
        {\popdisplay\fontsize{11}{12.5}\selectfont\color{white}\MakeUppercase{Passe à OnBuch+}}\par\vspace{4pt}
        {\raggedright\popsemi\fontsize{8.4}{11}\selectfont\color{white}Tous les cours signés, Tuteur IA illimité, fiches PDF — onglet OnBuch+ de l'app\par}
      \end{minipage}
    \end{tcolorbox}
  \end{minipage}
  \vspace{8pt}
  \popcontenu
  \vfill
  \signatureonbuch
  \restoregeometry
  \newpage
}

% Rangée « Ce cours contient » (page de garde)
\newcommand{\poptile}[3]{% #1 couleur #2 gros texte #3 légende
  \begin{tcolorbox}[enhanced,colback=#1,colframe=popink,boxrule=2pt,arc=9pt,shadow={2.5pt}{-2.5pt}{0pt}{popink},
      left=6pt,right=6pt,top=9pt,bottom=9pt,halign=center,valign=center,height=1.75cm,width=\linewidth,nobeforeafter]
    {\popdisplay\fontsize{12.5}{13}\selectfont\color{white}\MakeUppercase{#2}}\par\vspace{3pt}
    {\centering\popsemi\fontsize{7.8}{9.5}\selectfont\color{white}#3\par}
  \end{tcolorbox}}
\newcommand{\popcontenu}{%
  \par\noindent{\popxboldls\fontsize{7.5}{7.5}\selectfont\color{popmuted}CE COURS SIGNÉ CONTIENT}\par\vspace{7pt}
  \noindent\begin{minipage}[t]{0.235\linewidth}\poptile{popblue}{Cours}{complet, illustré, pas à pas}\end{minipage}\hfill
  \begin{minipage}[t]{0.235\linewidth}\poptile{poporange}{Méthodes}{et exercices résolus}\end{minipage}\hfill
  \begin{minipage}[t]{0.235\linewidth}\poptile{poppink}{Exercices}{type examen + corrigés}\end{minipage}\hfill
  \begin{minipage}[t]{0.235\linewidth}\poptile{popgreen}{Fiche bilan}{l'essentiel à retenir}\end{minipage}\par}

% Sommaire
\newtcolorbox{sommaire}{enhanced,colback=white,colframe=popink,boxrule=2.2pt,arc=10pt,
  drop shadow=popink,left=18pt,right=18pt,top=13pt,bottom=13pt,before skip=6pt,after skip=20pt,
  before upper={\poppill{Sommaire}{poppurple}{white}\par\vspace{9pt}\popsemi\fontsize{10.6}{14.5}\selectfont\color{popink}}}

% Signature de fin de cours — poussée tout en bas de la dernière page
\newcommand{\findecours}{%
  \par\vfill\nopagebreak
  \begin{center}\popsquiggle{poporange}\end{center}\vspace{4pt}
  \signatureonbuch
}
\AtBeginDocument{\def\llbracket{\mathopen{[\mkern-3.5mu[}}\def\rrbracket{\mathclose{]\mkern-3.5mu]}}} % intervalles d'entiers (glyphe ⟦ absent de la police)
% Glyphes absents de Fira Math : redéfinis à partir de glyphes disponibles
\AtBeginDocument{%
  \def\setminus{\mathbin{\backslash}}\let\smallsetminus\setminus
  \def\vdots{\mathord{\vbox{\baselineskip3.2pt\lineskiplimit0pt\kern4pt\hbox{.}\hbox{.}\hbox{.}}}}%
  \def\ddots{\mathinner{\mkern1mu\raise6.5pt\hbox{.}\mkern2mu\raise3.6pt\hbox{.}\mkern2mu\raise0.7pt\hbox{.}\mkern1mu}}%
  \def\triangle{\mathord{\scalebox{0.9}{$\symup{\Delta}$}}}%
  \def\nabla{\mathord{\raisebox{\depth}{\scalebox{1}[-1]{$\symup{\Delta}$}}}}%
  \def\checkmark{\popcheck{popdark}}%
  \def\star{\mathbin{\ast}}\def\bigstar{\mathord{\ast}}%
  \def\diamond{\mathbin{\scalebox{0.6}{$\Diamond$}}}%
  \def\bigcup{\mathop{\mathchoice{\raisebox{-0.3em}{\scalebox{1.7}{$\cup$}}}{\scalebox{1.25}{$\cup$}}{\cup}{\cup}}\displaylimits}%
  \def\bigcap{\mathop{\mathchoice{\raisebox{-0.3em}{\scalebox{1.7}{$\cap$}}}{\scalebox{1.25}{$\cap$}}{\cap}{\cap}}\displaylimits}%
  \def\lesseqqgtr{\mathrel{\lessgtr}}\def\bigoplus{\mathop{\oplus}\displaylimits}%
}
\providecommand{\ssi}{si et seulement si} % abréviation fréquente
\AtBeginDocument{\providecommand{\wideparen}{\overparen}} % arc de cercle

%% FIGFIX-BEGIN v5 — correctifs globaux des figures (docs/onbuch-generation/tools/figfix)
\usepackage{adjustbox}\usepackage{etoolbox}\tcbuselibrary{raster}
% toute figure TikZ/pgfplots plus large que la ligne est réduite à la largeur disponible
\BeforeBeginEnvironment{tikzpicture}{\begin{adjustbox}{max width=\linewidth}}
\AfterEndEnvironment{tikzpicture}{\end{adjustbox}}
% légendes pgfplots lisibles par-dessus les courbes
\pgfplotsset{every axis/.append style={legend style={fill=white,fill opacity=0.92,text opacity=1,draw=black!15,inner sep=3pt}}}
% étiquettes d'arêtes (syntaxe edge["..."]) avec fond blanc
\tikzset{every edge quotes/.append style={fill=white,inner sep=1.2pt}}
%% FIGFIX-END
\begin{document}

\pagegarde{Grands mouvements de l'atmosphère et de l'hydrosphère}{SVT}{1ère TI}{Module 3 : L'Éducation à l'Environnement et au Développement Durable}{NCT07}{3 h}{Cette leçon étudie la circulation générale de l'atmosphère (vents, cellules de convection) et la circulation océanique (courants marins), ainsi que leurs interactions. Elle explique l'origine dynamique des phénomènes climatiques dangereux (tempêtes, cyclones) et propose des mesures de protection adaptées.
\vspace{4pt}
\begin{multicols}{2}\small
\begin{itemize}
\item À la fin de cette leçon, je sais décrire sur un schéma les trois cellules de convection atmosphérique de chaque hémisphère.
\item À la fin de cette leçon, je sais identifier et localiser les grands systèmes de vents planétaires (alizés, vents d'ouest, vents polaires).
\item À la fin de cette leçon, je sais décrire un courant marin à partir d'une carte (sens, température, origine).
\item À la fin de cette leçon, je sais expliquer le rôle de la force de Coriolis dans la déviation des vents et des courants.
\item À la fin de cette leçon, je sais relier un phénomène météorologique dangereux (tempête, cyclone) à son origine dynamique.
\item À la fin de cette leçon, je sais expliquer les interactions entre l'océan et l'atmosphère (évaporation, transferts de chaleur).
\end{itemize}\end{multicols}}

\begin{sommaire}
\begin{multicols}{2}
\begin{enumerate}
\item Les moteurs de la circulation atmosphérique
\item La circulation générale de l'atmosphère : cellules et vents planétaires
\item La circulation océanique : les courants marins
\item Interactions océan-atmosphère
\item Les phénomènes climatiques dangereux : tempêtes et cyclones
\item Prévention et protection face aux phénomènes dangereux
\item Activité d'intégration
\item Méthodes et exercices résolus
\item Exercices
\item Corrigés des exercices
\item Fiche bilan
\end{enumerate}
\end{multicols}
\end{sommaire}

\begin{objectifs}
\begin{itemize}
\item À la fin de cette leçon, je sais décrire sur un schéma les trois cellules de convection atmosphérique de chaque hémisphère.
\item À la fin de cette leçon, je sais identifier et localiser les grands systèmes de vents planétaires (alizés, vents d'ouest, vents polaires).
\item À la fin de cette leçon, je sais décrire un courant marin à partir d'une carte (sens, température, origine).
\item À la fin de cette leçon, je sais expliquer le rôle de la force de Coriolis dans la déviation des vents et des courants.
\item À la fin de cette leçon, je sais relier un phénomène météorologique dangereux (tempête, cyclone) à son origine dynamique.
\item À la fin de cette leçon, je sais expliquer les interactions entre l'océan et l'atmosphère (évaporation, transferts de chaleur).
\item À la fin de cette leçon, je sais proposer des mesures de protection face aux tempêtes et cyclones.
\item À la fin de cette leçon, je sais interpréter une carte de situation météorologique simple (anticyclone, dépression).
\end{itemize}
\end{objectifs}

\begin{prerequis}
\begin{itemize}
\item Structure de l'atmosphère et de l'hydrosphère (leçons précédentes du Module 3)
\item Notion de pression atmosphérique et de gradient thermique (Seconde)
\item Répartition inégale de l'énergie solaire à la surface du globe (classe de 3e)
\item Rotation de la Terre sur elle-même et ses conséquences
\item Lecture de cartes et de schémas (légendes, flèches, échelles)
\end{itemize}
\end{prerequis}

\newpage

\coursec{Les moteurs de la circulation atmosphérique}

% Situation d'accroche et problématique
En octobre 2019, des pluies diluviennes accompagnées de vents violents ont provoqué des glissements de terrain meurtriers à Bafoussam. Chaque année, les navires au large de Kribi redoutent les houles liées aux courants marins, et les pêcheurs de Limbé observent que la mer se refroidit ou se réchauffe selon les saisons. D'où viennent ces vents, ces pluies et ces courants ? Pourquoi certaines régions du globe sont-elles régulièrement frappées par des cyclones dévastateurs alors que le Cameroun en est largement épargné ? Comprendre les grands mouvements de l'atmosphère et de l'hydrosphère permet d'expliquer ces phénomènes et de s'en protéger.

\courssub{Inégale répartition de l'énergie solaire : cause première des mouvements atmosphériques}

L'énergie solaire arrive sous forme de rayonnement électromagnétique. À l'équateur, les rayons frappent la surface presque perpendiculairement : la puissance reçue par unité de surface est maximale (moyenne annuelle \textasciitilde \SI{400}{\watt\per\meter\squared}). Aux pôles, l'incidence est rasante, la même énergie se répartit sur une surface plus grande et traverse une épaisseur d'atmosphère plus importante : le flux reçu est bien plus faible (moyenne annuelle \textasciitilde \SI{80}{\watt\per\meter\squared}). Cette \cle{inégale répartition} crée un \textbf{excédent radiatif} net dans la zone intertropicale (absorption $>$ émission) et un \textbf{déficit radiatif} net aux hautes latitudes (émission $>$ absorption). L'atmosphère et l'océan doivent transporter cet excédent vers les pôles pour rétablir l'équilibre thermique planétaire.

\begin{popfigure}
\begin{tikzpicture}[scale=1.2]
% Terre
\draw[popblue, thick] (0,0) circle (2.5);
% Équateur
\draw[poporange, dashed] (-2.5,0) -- (2.5,0) node[right] {\small Équateur};
% Pôles
\node[popdark] at (0,2.8) {Pôle Nord};
\node[popdark] at (0,-2.8) {Pôle Sud};
% Rayons solaires équateur (dense, perpendiculaires)
\foreach \i in {-2.2,-1.8,-1.4,-1.0,-0.6,-0.2,0.2,0.6,1.0,1.4,1.8,2.2}{
  \draw[popgold, ->, thick] (\i,3.5) -- (\i,0.2);
}
\node[popgold] at (0,3.8) {Rayonnement dense \&\ perpendiculaire};
% Rayons solaires pôles (diffus, rasants)
\foreach \i in {-2.0,-1.5,-1.0,-0.5,0,0.5,1.0,1.5,2.0}{
  \draw[popgold!50, ->] (\i,3.5) -- (\i,2.3);
  \draw[popgold!50, ->] (\i,-3.5) -- (\i,-2.3);
}
\node[popgold!50] at (0,-3.8) {Rayonnement diffus \&\ rasant};
% Flèches de bilan radiatif
\draw[popgreen, <->, thick] (-3.2,1.2) -- (-3.2,-1.2) node[midway, left, align=center] {Excédent\\radiatif};
\draw[popred, <->, thick] (3.2,2.5) -- (3.2,1.2) node[midway, right, align=center]{Déficit\\radiatif};
\draw[popred, <->, thick] (3.2,-1.2) -- (3.2,-2.5) node[midway, right, align=center]{Déficit\\radiatif};
% Flèche transport méridien
\draw[popblue, ->, very thick, dashed] (-3.5,0) -- (-3.5,2.5) node[midway, left] {Transport de chaleur};
\draw[popblue, ->, very thick, dashed] (-3.5,0) -- (-3.5,-2.5);
\end{tikzpicture}
\legende{Bilan radiatif terrestre : excédent à l'équateur, déficit aux pôles. L'atmosphère et l'océan assurent le transport méridien de chaleur (flèches bleues en pointillés).}
\end{popfigure}

\courssub{Réchauffement de l'air équatorial, convection et pression atmosphérique}

L'air au contact du sol équatorial se réchauffe par conduction, se dilate et sa \cle{densité} diminue ($\rho = m/V$). D'après le principe d'Archimède appliqué aux fluides, une masse de fluide moins dense que le milieu environnant à la même pression subit une force de poussée vers le haut : c'est la \textbf{convection thermique}. L'air chaud s'élève adiabatiquement, créant en altitude une zone de \textbf{basse pression} (dépression) et au sol une convergence des vents. Symétriquement, aux pôles l'air se refroidit, se contracte, devient plus dense et descend, formant une \textbf{haute pression} (anticyclone) au niveau du sol et une divergence des vents.

\begin{definition}[Pression atmosphérique]
La pression atmosphérique en un point est le poids de la colonne d'air au-dessus de la surface unitaire. Elle s'exprime en hectopascals (hPa) : 1 hPa = 100 Pa = 1 mbar. Une masse d'air chaud correspond à une pression plus faible en altitude (dépression dynamique) ; une masse d'air froid correspond à une pression plus forte au sol (anticyclone thermique).
\end{definition}

\begin{popfigure}
\begin{tikzpicture}[scale=1.1]
% Coupe atmosphère équateur-pôle (schéma de principe cellule de Hadley)
\draw[popblue, thick] (-4,0) -- (4,0) node[midway, below] {Surface terrestre}; % surface
\draw[popblue, thick] (-4,0) -- (-4,3); % bord gauche
\draw[popblue, thick] (4,0) -- (4,3); % bord droit
\draw[popmuted, dashed] (-4,3) -- (4,3) node[midway, above] {Tropopause (\textasciitilde 12-15 km)};
% Équateur à gauche
\node[poporange, align=center] at (-3, -0.6) {Équateur\\(Chaud)};
\draw[popred, ->, thick] (-3,0.3) -- (-3,2.5) node[midway, left, align=center]{Ascendance\\ air chaud};
\draw[popgreen, ->, thick] (-3,2.7) -- (-3,2.9) node[above] {Basse pression (L)};
% Pôle à droite
\node[poppurple, align=center] at (3, -0.6) {Pôle\\(Froid)};
\draw[popblue, ->, thick] (3,2.5) -- (3,0.3) node[midway, right, align=center]{Descendance\\ air froid};
\draw[popred, ->, thick] (3,0) -- (3,-0.3) node[below] {Haute pression (H)};
% Flèches horizontales altitude (vers pôle) et surface (vers équateur)
\draw[popmuted, ->, thick] (-3,2.8) -- (-0.5,2.8) node[midway, above] {Vent d'altitude (vers pôle)};
\draw[popmuted, ->, thick] (0.5,2.8) -- (3,2.8);
\draw[popmuted, ->, thick] (3,-0.2) -- (0.5,-0.2) node[midway, below] {Vent de surface (vers équateur)};
\draw[popmuted, ->, thick] (-0.5,-0.2) -- (-3,-0.2);
% Légendes cellules
\node[popdark, rotate=90] at (-3.7,1.5) {Cellule de Hadley (simplifiée)};
\end{tikzpicture}
\legende{Mécanisme convectif à l'échelle planétaire : ascendance de l'air chaud à l'équateur (basse pression), descendance de l'air froid aux pôles (haute pression), vents de compensation en altitude et au sol.}
\end{popfigure}

\courssub{Du gradient de pression au vent : définition et rôle de la rotation terrestre}

La différence de pression horizontale entre deux zones engendre une \textbf{force de gradient de pression} ($\vec{F}_{grad} = -\frac{1}{\rho}\nabla P$) qui pousse l'air des hautes pressions vers les basses pressions. Ce déplacement d'air est ce que nous appelons le \cle{vent}. Par convention, le vent est nommé d'après sa \textbf{provenance} (ex: un vent du Nord souffle du Nord vers le Sud).

\begin{definition}[Vent]
Le vent est un déplacement d'air d'une zone de haute pression (anticyclone) vers une zone de basse pression (dépression), sous l'effet de la force de gradient de pression. Sa direction est donnée par l'orientation du gradient (des H vers les L) et sa vitesse est proportionnelle à l'intensité du gradient (rapprochement des isobares).
\end{definition}

Cependant, la Terre tourne sur elle-même (rotation vers l'Est). Tout mouvement à la surface du globe subit une \textbf{force apparente} : la \cle{force de Coriolis} ($\vec{F}_c = 2m\vec{\Omega} \times \vec{v}$). Elle dévie la trajectoire du vent vers la \textbf{droite} dans l'hémisphère Nord et vers la \textbf{gauche} dans l'hémisphère Sud. Cette déviation est nulle à l'équateur ($\sin 0^\circ = 0$) et maximale aux pôles ($\sin 90^\circ = 1$).

\begin{definition}[Force de Coriolis]
La force de Coriolis est une force apparente due à la rotation de la Terre, qui dévie tout mouvement horizontal à la surface du globe (admise, non démontrée au programme). Dans l'hémisphère Nord, la déviation est vers la droite ; dans l'hémisphère Sud, vers la gauche. Elle est proportionnelle à la vitesse du mouvement et au sinus de la latitude ($f = 2\Omega\sin\varphi$).
\end{definition}

\begin{propriete}[Équilibre des forces : Vent géostrophique et vent de surface]
\begin{itemize}
  \item \textbf{En altitude (sans frottement) :} Le vent devient \textbf{géostrophique} : force de Coriolis = force de gradient de pression. Il souffle \textbf{parallèle aux isobares} (sens antihoraire autour d'une dépression, horaire autour d'un anticyclone dans l'hémisphère Nord).
  \item \textbf{Au sol (avec frottement) :} Le frottement réduit la vitesse $\rightarrow$ force de Coriolis diminue $\rightarrow$ le vent traverse les isobares vers la basse pression (angle \textasciitilde 15\textdegree--30\textdegree). Convergence au sol dans les dépressions (ascendance), divergence dans les anticyclones (descendance).
\end{itemize}
\end{propriete}

\begin{attention}[Erreurs fréquentes à ne pas commettre]
\begin{itemize}
  \item \textbf{« L'air chaud monte parce qu'il est léger »} : L'air chaud ne monte pas parce qu'il serait « léger » en absolu, mais parce qu'il est \textbf{moins dense} que l'air environnant \textbf{à la même pression}. C'est la différence de densité qui crée la poussée d'Archimède. Dire « l'air chaud est léger » est une simplification trompeuse : à température égale, l'air humide est moins dense que l'air sec, pourtant il ne monte pas systématiquement sans déclencheur.
  \item \textbf{Confondre gradient de pression et Coriolis} : Le gradient \textbf{crée} le mouvement (vitesse), Coriolis ne fait que le \textbf{dévier} (direction). Sans gradient, pas de vent, donc pas de Coriolis.
  \item \textbf{Sens de nomination du vent} : Un « vent d'Ouest » vient de l'Ouest et va vers l'Est. Sur une carte, la flèche du vent pointe vers l'Est.
  \item \textbf{Coriolis à l'équateur} : Elle est nulle. Il n'y a pas de rotation cyclonique possible à l'équateur même.
\end{itemize}
\end{attention}

\courssub{Circulation locale (brises) et circulation générale (planétaire)}

Les mécanismes décrits ci-dessus (gradient de pression + Coriolis) s'appliquent à deux échelles distinctes :

\begin{itemize}
  \item \textbf{Circulation locale (brises) :} Différence de température entre terre et mer (brise de mer le jour, brise de terre la nuit) ou entre versant ensoleillé et versant ombragé (brise de vallée/anabatique le jour, brise de montagne/katabatique la nuit). L'échelle spatiale est de quelques km à quelques dizaines de km, la durée de quelques heures (cycle nycthéméral). La force de Coriolis y est souvent négligeable.
  \item \textbf{Circulation générale (planétaire) :} Organisée en trois grandes cellules de convection méridiennes par hémisphère (\textbf{Hadley} 0\textdegree--30\textdegree, \textbf{Ferrel} 30\textdegree--60\textdegree, \textbf{Polaire} 60\textdegree--90\textdegree) qui ceinturent le globe. Elle résulte de l'inégale répartition solaire à l'échelle planétaire et est structurée par la force de Coriolis. C'est elle qui détermine les \cle{vents planétaires} permanents : \textbf{alizés} (NE au Nord, SE au Sud, convergents à l'équateur $\rightarrow$ ZCIT), \textbf{vents d'ouest} (moyennes latitudes), \textbf{vents polaires} (est).
\end{itemize}

\begin{methode}[Décrire un mouvement atmosphérique sur une carte isobare ou un schéma]
\begin{enumerate}
  \item \textbf{Identifier les centres d'action :} Repérer les zones de haute pression (H, anticyclone) et basse pression (L, dépression) grâce aux isobares (lignes de pression égale) et aux symboles.
  \item \textbf{Tracer la force de gradient de pression :} Dessiner des flèches perpendiculaires aux isobares, orientées des H vers les L (rapprochement des isobares = gradient fort = vent fort).
  \item \textbf{Appliquer la déviation de Coriolis :} Dévier chaque flèche vers la \textbf{droite} (hémisphère Nord) ou vers la \textbf{gauche} (hémisphère Sud).
  \item \textbf{Préciser l'altitude :}
        \begin{itemize}
          \item \textit{En altitude (libre) :} Vent parallèle aux isobares (géostrophique). Sens antihoraire autour d'un L, horaire autour d'un H (Hémisphère Nord).
          \item \textit{Au sol :} Vent traversant les isobares vers le centre du L (convergence) ou s'en écartant depuis le H (divergence), angle \textasciitilde 20\textdegree--30\textdegree.
        \end{itemize}
  \item \textbf{Conclure :} Nommer le vent (provenance), indiquer l'échelle (locale/planétaire) et le lien avec la cellule de convection concernée (Hadley, Ferrel, Polaire).
\end{enumerate}
\end{methode}

\courssub{Exemple concret : brise de mer à Kribi}

\begin{exemplebox}[Brise de mer diurne à Kribi (Cameroun, 3\textdegree N)]
En journée, le sol de la côte de Kribi (sable, roche) se réchauffe plus vite que l'océan Atlantique (inertie thermique forte). L'air au-dessus de la terre devient plus chaud, moins dense, et s'élève par convection : création d'une \textbf{basse pression thermique} locale au sol. L'air plus frais et plus dense au-dessus de la mer constitue une \textbf{haute pression thermique} locale. Le gradient de pression dirige l'air de la mer vers la terre : c'est la \textbf{brise de mer} (vent d'Ouest/Sud-Ouest à Kribi), qui apporte fraîcheur et humidité. La nuit, le sol se refroidit plus vite que la mer : le processus s'inverse (brise de terre, vent d'Est/Nord-Est). Ce cycle quotidien illustre parfaitement le lien \og gradient de pression $\rightarrow$ vent \fg à l'échelle locale, sans influence notable de Coriolis.
\end{exemplebox}

\courssub{Exercice résolu : détermination du sens du vent à partir d'une carte de pression}

\begin{exoresolu}[Vent autour d'une dépression en hémisphère Nord]
\textbf{Énoncé :} Sur une carte isobare de l'hémisphère Nord, une dépression (L, 990 hPa) est centrée à 45\textdegree N. Les isobares (1000, 1004, 1008 hPa) sont circulaires, espacées régulièrement. Indiquer le sens et les caractéristiques du vent au niveau du sol autour de cette dépression.

\textbf{Solution détaillée :}
\begin{enumerate}
  \item \textbf{Analyse du gradient :} La pression augmente de l'intérieur (990 hPa) vers l'extérieur (1008 hPa). La force de gradient de pression est orientée \textbf{radialement vers l'intérieur} (des H vers le L). L'espacement régulier indique un gradient constant, donc une vitesse de vent constante.
  \item \textbf{Application de Coriolis (Hémisphère Nord) :} La force de Coriolis dévie le mouvement vers la \textbf{droite} par rapport à la direction du gradient.
  \item \textbf{Résultante au sol (frottement) :} Le vent souffle \textbf{dans le sens inverse des aiguilles d'une montre (antihoraire)} autour de la dépression, \textbf{croisant les isobares vers l'intérieur} avec un angle d'environ 20\textdegree--30\textdegree (convergence). Cette convergence force l'air à s'élever au centre (ascendance, nuages, précipitations).
  \item \textbf{Nom du vent :} Selon la position sur le pourtour, le vent vient du Nord-Est (côté Est), du Sud-Est (côté Sud), du Sud-Ouest (côté Ouest), du Nord-Ouest (côté Nord).
\end{enumerate}
\emph{Remarque :} En altitude (niveau 500 hPa, sans frottement), le vent serait géostrophique, strictement parallèle aux isobares, sens antihoraire pur, sans convergence.
\end{exoresolu}

\begin{savaistu}[Pourquoi le Cameroun est-il épargné par les cyclones tropicaux ?]
La force de Coriolis est proportionnelle au sinus de la latitude ($f = 2\Omega\sin\varphi$) : elle est \textbf{nulle à l'équateur} ($\varphi=0$) et augmente vers les pôles. Les cyclones tropicaux (ou ouragans/typhons) sont des dépressions chaudes organisées qui ont besoin d'une force de Coriolis \textbf{suffisante} pour initier et maintenir la rotation cyclonique de l'air convergent vers le centre ; ils ne se forment donc jamais à moins d'environ \SI{5}{\degree}--\SI{8}{\degree} de latitude de l'équateur. Le Cameroun, situé entre \SI{2}{\degree} N (sud, Campo) et \SI{13}{\degree} N (nord, Lac Tchad), se trouve en grande partie dans cette zone équatoriale de Coriolis faible, ce qui le protège naturellement de la \textbf{formation} de cyclones. En revanche, il peut subir les effets indirects (pluies diluviennes, vents forts, houle) de systèmes dépressionnaires tropicaux (ondes de l'Est, dépressions du Monsoon) qui se forment plus au nord (Atlantique tropical Nord) ou au sud (océan Indien) et dont les bandes nuageuses s'étendent vers l'équateur.
\end{savaistu}

\courssub{Synthèse des moteurs}

\begin{aretenir}[Moteurs de la circulation atmosphérique]
\begin{itemize}
  \item \textbf{Source d'énergie :} Répartition inégale du rayonnement solaire (excédent radiatif net intertropical, déficit net polaire). Température moyenne équateur \textasciitilde 27\textdegree C, pôles \textasciitilde -30\textdegree C.
  \item \textbf{Mécanisme thermique :} Réchauffement sol $\rightarrow$ conduction air $\rightarrow$ dilatation $\rightarrow$ diminution densité $\rightarrow$ convection (ascendance air chaud, descendance air froid). Principe d'Archimède dans l'air.
  \item \textbf{Pression :} Air chaud ascendant = basse pression (L, dépression) au sol ; Air froid descendant = haute pression (H, anticyclone) au sol.
  \item \textbf{Moteur dynamique (Vent) :} Force de gradient de pression (H $\rightarrow$ L) $\rightarrow$ vitesse du vent proportionnelle au gradient (rapprochement isobares).
  \item \textbf{Déviation (Direction) :} Force de Coriolis (rotation Terre) $\rightarrow$ Droite (Nord), Gauche (Sud), Nulle (Équateur). Dépend de la vitesse et de la latitude.
  \item \textbf{Échelles :} Locale (brises, km, heures, Coriolis négligeable) vs Générale (cellules Hadley/Ferrel/Polaire, 10 000 km, permanente, Coriolis structurant).
  \item \textbf{Vents planétaires résultants :} Alizés (NE/SE, convergents $\rightarrow$ ZCIT), Vents d'Ouest (moyennes latitudes), Vents polaires (Est).
\end{itemize}
\end{aretenir}

\coursec{La circulation générale de l'atmosphère : cellules et vents planétaires}

\courssub{Observation et questionnement : pourquoi les vents soufflent-ils selon des directions préférentielles ?}

Imagine un pêcheur sur les côtes de Kribi ou un paysan dans le Nord du Cameroun. Tous deux observent que les vents ne soufflent pas au hasard : à certaines périodes, un vent régulier vient du nord-est, à d'autres du sud-ouest. Les marins d'autrefois, pour traverser l'Atlantique, utilisaient des \cle{vents constants} qu'ils appelaient \emph{alizés} — vents qui poussent les navires des Canaries vers les Caraïbes. Mais pourquoi ces vents sont-ils si réguliers ? Pourquoi trouve-t-on des zones de calmes (les \emph{calmes équatoriaux} et les \emph{calmes des chevaux}) et des zones de tempêtes récurrentes ?

\begin{experience}[Observation des vents de surface à l'échelle planétaire]
\textbf{Matériel :} Carte des vents moyens de surface (janvier et juillet) issue des réanalyses météorologiques (ex. ERA5).\\
\textbf{Protocole :}
\begin{enumerate}
\item Repère sur la carte les zones où les vents soufflent de manière constante en direction et en force.
\item Note les latitudes où les vents changent brutalement de direction.
\item Compare les hémisphères Nord et Sud.
\end{enumerate}
\textbf{Observations :}
\begin{itemize}
\item Entre l'équateur et $\approx 30^\circ$ de latitude, les vents convergent vers l'équateur : au nord-est dans l'hémisphère Nord, au sud-est dans l'hémisphère Sud (alizés).
\item Vers $30^\circ$ N et S, les vents faiblissent (zones de calmes, hautes pressions subtropicales).
\item Entre $30^\circ$ et $60^\circ$, les vents soufflent dominamment d'ouest en est (vents d'ouest).
\item Au-delà de $60^\circ$, vents d'est polaires.
\item Une bande de convergence équatoriale (ZCIT) se déplace au rythme des saisons.
\end{itemize}
\textbf{Question :} Quelle organisation de la circulation atmosphérique produit ce motif reproductible ?
\end{experience}

\courssub{Les cellules de convection : le moteur de la circulation générale}

\begin{definition}[Cellule de convection atmosphérique]
Une \cle{cellule de convection} est une boucle fermée de circulation de l'air entre deux latitudes : l'air chaud monte (ascendance), se déplace en altitude, redescend (subsidence) et revient en surface vers son point de départ. Elle transforme l'énergie thermique (différence de température) en énergie cinétique (vent).
\end{definition}

\begin{propriete}[Modèle des trois cellules par hémisphère (admis)]
La circulation atmosphérique s'organise en \textbf{trois cellules de convection par hémisphère}, séparées par des zones de hautes et de basses pressions :
\begin{itemize}
\item \textbf{Cellule de Hadley} ($0^\circ$–$30^\circ$) : ascendance équatoriale, subsidence subtropicale.
\item \textbf{Cellule de Ferrel} ($30^\circ$–$60^\circ$) : circulation indirecte (moteur : instabilité barocline, non convection pure).
\item \textbf{Cellule polaire} ($60^\circ$–$90^\circ$) : ascendance subpolaire, subsidence polaire.
\end{itemize}
\end{propriete}

\begin{methode}[Démarche d'investigation : comprendre la cellule de Hadley]
\begin{enumerate}
\item \textbf{Observation :} L'équateur reçoit un rayonnement solaire intense toute l'année $\rightarrow$ fort réchauffement de la surface $\rightarrow$ air chaud, léger, humide.
\item \textbf{Hypothèse :} Cet air monte, créant une basse pression au sol ; en altitude, il diverge vers les pôles ; quelque part, il doit redescendre pour fermer la boucle.
\item \textbf{Données (radiosondages, satellites) :} 
\begin{itemize}
\item Ascendance vigoureuse jusqu'à la tropopause ($\approx 15-18$ km) à l'équateur.
\item Divergence en altitude vers $30^\circ$ N/S.
\item Subsidence marquée vers $30^\circ$ (air sec, stable, haute pression au sol).
\item Retour en surface vers l'équateur : alizés.
\end{itemize}
\item \textbf{Interprétation :} La cellule de Hadley est une machine thermique directe : elle transporte la chaleur de l'équateur vers les subtropiques.
\item \textbf{Conclusion :} Le gradient de température équateur-pôle est le moteur ultime ; la rotation de la Terre (force de Coriolis) dévie les flux, créant les vents planétaires.
\end{enumerate}
\end{methode}

\begin{popfigure}
\begin{tikzpicture}[scale=0.9, every node/.style={font=\small}]
% Earth surface
\draw[thick, popblue!80!popdark] (-6,0) -- (6,0) node[midway, below=2pt] {Surface terrestre};
% Latitude markers
\foreach \lat/\x in {0/0, 30/2.5, 60/5, 90/7.5} {
  \draw[dashed, popmuted] (\x,0) -- (\x,5);
  \draw[dashed, popmuted] (-\x,0) -- (-\x,5);
  \node[below] at (\x,0) {\lat°};
  \node[below] at (-\x,0) {\lat°};
}
% Hadley cell (0-30)
\draw[poporange, thick, ->] (0,0) -- (0,3.5) node[midway, right] {Ascendance};
\draw[poporange, thick, ->] (0,3.5) -- (2.5,3.5) node[midway, above] {Divergence};
\draw[poporange, thick, ->] (2.5,3.5) -- (2.5,0) node[midway, right] {Subsidence};
\draw[poporange, thick, ->] (2.5,0) -- (0,0) node[midway, below] {Alizés (retour)};
% Ferrel cell (30-60)
\draw[popgreen, thick, ->] (2.5,0) -- (2.5,2.5) node[midway, right] {Ascendance};
\draw[popgreen, thick, ->] (2.5,2.5) -- (5,2.5) node[midway, above] {Divergence};
\draw[popgreen, thick, ->] (5,2.5) -- (5,0) node[midway, right] {Subsidence};
\draw[popgreen, thick, ->] (5,0) -- (2.5,0) node[midway, below] {Vents d'ouest};
% Polar cell (60-90)
\draw[poppurple, thick, ->] (5,0) -- (5,1.5) node[midway, right] {Ascendance};
\draw[poppurple, thick, ->] (5,1.5) -- (7.5,1.5) node[midway, above] {Divergence};
\draw[poppurple, thick, ->] (7.5,1.5) -- (7.5,0) node[midway, right] {Subsidence};
\draw[poppurple, thick, ->] (7.5,0) -- (5,0) node[midway, below] {Vents d'est polaires};
% Pressure labels
\node[popblue, align=center] at (0,4.2) {Basse pression\\équatoriale};
\node[poporange, align=center] at (2.5,4.2) {Haute pression\\subtropicale};
\node[popgreen, align=center] at (5,3.5) {Basse pression\\subpolaire};
\node[poppurple, align=center] at (7.5,2.5) {Haute pression\\polaire};
% Hemisphere label
\node[popdark, font=\bfseries] at (-7,2.5) {Hémisphère Nord};
\end{tikzpicture}
\legende{Coupe méridienne simplifiée de la circulation générale dans l'hémisphère Nord. Les trois cellules (Hadley en orange, Ferrel en vert, Polaire en violet) montrent les ascendance, divergence, subsidence et retour en surface. Les zones de pression au sol sont indiquées.}
\end{popfigure}

\begin{attention}[Piège fréquent : alizés vs contre-alizés]
Ne confonds pas :
\begin{itemize}
\item \textbf{Alizés} : vents de \emph{surface} (retour de la cellule de Hadley), soufflant \emph{vers l'équateur} (NE au Nord, SE au Sud).
\item \textbf{Contre-alizés} : courants d'\emph{altitude} (branche montante et divergente de la cellule de Hadley), soufflant \emph{vers les pôles} (SW au Nord, NW au Sud).
\end{itemize}
Sur un schéma, les flèches en bas (près du sol) sont les alizés ; celles en haut sont les contre-alizés.
\end{attention}

\courssub{Les grands vents planétaires et la zone de convergence intertropicale (ZCIT)}

La rotation de la Terre dévie les mouvements d'air (force de Coriolis) : dans l'hémisphère Nord, déviation vers la \emph{droite} ; dans l'hémisphère Sud, vers la \emph{gauche}. Cela transforme les flux méridiens (Nord-Sud) en vents zonaux (Est-Ouest).

\begin{center}
\begin{tabularx}{\linewidth}{|l|X|X|X|}
\hline
\rowcolor{popblueL} \textbf{Cellule} & \textbf{Latitudes} & \textbf{Vent de surface (nom \& direction)} & \textbf{Pression au sol} \\
\hline
Hadley & $0^\circ$–$30^\circ$ & Alizés : NE (HN) / SE (HS) $\rightarrow$ vers l'équateur & Haute subtropicale ($\approx 30^\circ$) \\
Ferrel & $30^\circ$–$60^\circ$ & Vents d'ouest (SW $\rightarrow$ NE en HN ; NW $\rightarrow$ SE en HS) & Basse subpolaire ($\approx 60^\circ$) \\
Polaire & $60^\circ$–$90^\circ$ & Vents d'est polaires (NE $\rightarrow$ SW en HN ; SE $\rightarrow$ NW en HS) & Haute polaire ($90^\circ$) \\
\hline
\end{tabularx}
\end{center}

\begin{definition}[Zone de convergence intertropicale (ZCIT ou ITCZ)]
La \cle{ZCIT} est la bande étroite où les alizés de l'hémisphère Nord et de l'hémisphère Sud se rencontrent. La convergence forcée provoque une ascendance massive, une instabilité convective et des \textbf{précipitations intenses et régulières}. Elle suit le maximum d'insolation solaire (zénith) avec un décalage d'environ 1 à 2 mois.
\end{definition}

\begin{exemplebox}[Saisons au Cameroun et déplacement de la ZCIT]
Au Cameroun, le régime des pluies est dicté par la ZCIT :
\begin{itemize}
\item \textbf{Janvier :} ZCIT au sud de l'équateur ($\approx 5^\circ$ S) $\rightarrow$ saison sèche au Nord (harmattan), courte saison sèche au Sud.
\item \textbf{Avril--Juin :} ZCIT remonte vers le nord, traverse le pays $\rightarrow$ début des pluies au Sud, puis au Centre, enfin au Nord (grande saison des pluies au Nord en juillet--août).
\item \textbf{Juillet--Août :} ZCIT au maximum nord ($\approx 10^\circ$–$12^\circ$ N) $\rightarrow$ forte pluviométrie au Nord (région de Maroua, Garoua), « petite saison sèche » au Sud (effet de la subsidence hadleyenne).
\item \textbf{Octobre--Novembre :} ZCIT redescend $\rightarrow$ deuxième saison des pluies au Sud, fin des pluies au Nord.
\end{itemize}
Ce balancier annuel explique pourquoi Yaoundé a deux saisons des pluies et le Nord une seule.
\end{exemplebox}

\begin{popfigure}
\begin{tikzpicture}[scale=0.55]
% World map simplified: equator, tropics, polar circles
\draw[popblue!50!popdark, thick] (-180,0) -- (180,0) node[midway, below=2pt] {Équateur};
\draw[popmuted, dashed] (-180,23.5) -- (180,23.5) node[midway, right] {Tropique Cancer};
\draw[popmuted, dashed] (-180,-23.5) -- (180,-23.5) node[midway, right] {Tropique Capricorne};
\draw[popmuted, dashed] (-180,66.5) -- (180,66.5);
\draw[popmuted, dashed] (-180,-66.5) -- (180,-66.5);
% ITCZ positions
\draw[poppink, thick, decorate, decoration={snake, amplitude=2pt, segment length=10pt}] (-180,5) -- (180,5) node[midway, above] {ZCIT juillet};
\draw[poporange, thick, decorate, decoration={snake, amplitude=2pt, segment length=10pt}] (-180,-5) -- (180,-5) node[midway, below] {ZCIT janvier};
% Trade wind arrows (deflected by Coriolis)
% Northern Hemisphere trades (NE -> SW)
\foreach \lon in {-150,-90,-30,30,90,150} {
  \draw[poporange, thick, ->] (\lon,30) -- (\lon-5,15) node[midway, left, font=\tiny] {Alizés NE};
  \draw[poporange, thick, ->] (\lon,15) -- (\lon-5,5);
}
% Southern Hemisphere trades (SE -> NW)
\foreach \lon in {-150,-90,-30,30,90,150} {
  \draw[poporange, thick, ->] (\lon,-30) -- (\lon+5,-15) node[midway, right, font=\tiny] {Alizés SE};
  \draw[poporange, thick, ->] (\lon,-15) -- (\lon+5,-5);
}
% Westerlies (mid-latitudes)
\foreach \lat in {45,-45} {
  \foreach \lon in {-150,-90,-30,30,90,150} {
    \draw[popgreen, thick, ->] (\lon,\lat) -- (\lon+15,\lat) node[midway, above, font=\tiny] {Vents d'ouest};
  }
}
% Polar easterlies
\foreach \lat in {75,-75} {
  \foreach \lon in {-150,-90,-30,30,90,150} {
    \draw[poppurple, thick, ->] (\lon,\lat) -- (\lon-10,\lat) node[midway, above, font=\tiny] {Vents d'est polaires};
  }
}
% Cameroon marker
\fill[popred] (12,4) circle (3pt) node[above right] {Cameroun};
\end{tikzpicture}
\legende{Carte simplifiée des vents planétaires et position saisonnière de la ZCIT. Les alizés (orange) convergent vers la ZCIT (ligne ondulée rose en juillet, orange en janvier). Les vents d'ouest (vert) dominent les moyennes latitudes, les vents d'est polaires (violet) aux hautes latitudes. Le Cameroun (point rouge) subit le passage de la ZCIT.}
\end{popfigure}

\begin{methode}[Décrire la circulation générale sur un schéma ou une carte]
Pour réussir l'exercice de description (savoir-faire exigible au Probatoire) :
\begin{enumerate}
\item \textbf{Repère l'équateur et les pôles} ; identifie l'hémisphère concerné.
\item \textbf{Nomme chaque cellule} de l'équateur vers le pôle : Hadley, Ferrel, Polaire.
\item \textbf{Indique le sens de rotation} de chaque cellule (horaire/anti-horaire selon hémisphère) et les branches : ascendance, divergence, subsidence, retour.
\item \textbf{Localise les zones de pression} : basse équatoriale, hautes subtropicales, basses subpolaires, hautes polaires.
\item \textbf{Nomme les vents de surface} dans chaque cellule (alizés, vents d'ouest, vents d'est polaires) et leur direction (flèches).
\item \textbf{Signale la ZCIT} (convergence des alizés) et son déplacement saisonnier si la carte montre deux positions.
\item \textbf{Utilise le vocabulaire précis} : « ascendance convective », « subsidence adiabatique », « divergence/convergence », « force de Coriolis ».
\end{enumerate}
\end{methode}

\begin{exoresolu}[Description d'un schéma de circulation générale]
\textbf{Énoncé :} À partir du schéma de la coupe méridienne (figure ci-dessus), décris la circulation atmosphérique dans l'hémisphère Nord en utilisant les termes : cellule de Hadley, cellule de Ferrel, cellule polaire, alizés, vents d'ouest, vents d'est polaires, haute pression subtropicale, basse pression subpolaire, ZCIT.

\textbf{Solution détaillée :}
\begin{enumerate}
\item \textbf{Cellule de Hadley ($0^\circ$–$30^\circ$ N)} : L'air chaud et humide monte à l'équateur (basse pression équatoriale, ZCIT) $\rightarrow$ divergence en altitude vers le nord $\rightarrow$ subsidence vers $30^\circ$ N (haute pression subtropicale, air sec) $\rightarrow$ retour en surface vers l'équateur sous forme d'\textbf{alizés de nord-est} (déviés par Coriolis).
\item \textbf{Cellule de Ferrel ($30^\circ$–$60^\circ$ N)} : Circulation indirecte. L'air de surface monte vers $60^\circ$ N (basse pression subpolaire) $\rightarrow$ divergence en altitude vers le sud et le nord $\rightarrow$ subsidence vers $30^\circ$ N et vers le pôle $\rightarrow$ vents de surface d'\textbf{ouest} (sud-ouest $\rightarrow$ nord-est) dus à la déviation de Coriolis sur le flux polaire-équatorial.
\item \textbf{Cellule polaire ($60^\circ$–$90^\circ$ N)} : Ascendance vers $60^\circ$ N, divergence en altitude vers le pôle, subsidence au pôle (haute pression polaire), retour en surface sous forme de \textbf{vents d'est polaires} (nord-est $\rightarrow$ sud-ouest).
\item \textbf{ZCIT} : Bande de convergence des alizés hémisphériques, siège de forts orages, se déplace au rythme du zénith solaire.
\end{enumerate}
\end{exoresolu}

\begin{savaistu}[Le rôle du Cameroun dans la surveillance de la ZCIT]
L'ASECNA (Agence pour la Sécurité de la Navigation Aérienne en Afrique et à Madagascar), dont le siège est à Dakar mais avec une forte présence à Douala et Yaoundé, utilise les données des radiosondages de Garoua, Ngaoundéré, Yaoundé et Douala pour suivre en temps réel la position de la ZCIT. Ces données alimentent les modèles de prévision des orages violents et des crues du Bénoué et du Logone, cruciaux pour l'agriculture et la prévention des inondations dans le Nord.
\end{savaistu}

\begin{exercice}[Entraînement : lecture de carte des vents]
Sur la carte mondiale fournie (figure 2), identifie :
\begin{enumerate}
\item La latitude approximative de la ZCIT en juillet et en janvier.
\item Le sens des alizés au large des côtes du Cameroun (golfe de Guinée) en janvier.
\item La zone où les vents d'ouest sont les plus forts (indice : « quarantièmes rugissants »).
\item Explique pourquoi les alizés ne soufflent pas exactement du nord au sud mais sont déviés.
\end{enumerate}
\end{exercice}

\begin{corrige}[Exercice 1]
\begin{enumerate}
\item Juillet : ZCIT $\approx 5^\circ$–$10^\circ$ N (ligne rose) ; Janvier : ZCIT $\approx 5^\circ$ S (ligne orange).
\item En janvier, le Cameroun est au nord de la ZCIT $\rightarrow$ alizés de nord-est (harmattan) soufflent du continent vers l'océan : vent sec, poussiéreux, direction NE $\rightarrow$ SW.
\item Les « quarantièmes rugissants » ($40^\circ$–$50^\circ$ S) : vents d'ouest très forts, peu de terres pour les freiner.
\item Force de Coriolis : dans l'hémisphère Nord, déviation vers la droite $\rightarrow$ vent du nord devient nord-est ; dans l'hémisphère Sud, déviation vers la gauche $\rightarrow$ vent du sud devient sud-est.
\end{enumerate}
\end{corrige}

\coursec{La circulation océanique : les courants marins}

Dans la section précédente, nous avons vu que l'atmosphère est parcourue par de grands vents planétaires — alizés, vents d'ouest — organisés en cellules de convection. Mais ces vents ne soufflent pas sur une Terre immobile : ils frottent la surface des océans et, par simple entraînement, mettent l'eau en mouvement. L'océan devient alors, à son tour, un immense système de circulation. C'est cette circulation océanique que nous allons maintenant étudier : elle joue un rôle capital dans le climat des côtes camerounaises, dans la pêche et dans la genèse des phénomènes dangereux que nous aborderons ensuite.

\courssub{Qu'est-ce qu'un courant marin ?}

\textbf{Une observation pour commencer.}\par
En 1992, un conteneur transportant des milliers de canards en plastique s'est déversé dans le Pacifique Nord lors d'une tempête. Les années suivantes, ces jouets ont été retrouvés sur les côtes de l'Alaska, d'Hawaï, puis de l'Angleterre, après un voyage de plus de quinze ans. Les canards n'ont pas dérivé au hasard : ils ont suivi des « routes » d'eau bien précises. De même, au large de Kribi ou de Limbé, les pêcheurs savent qu'une pirogue à la dérive ne part pas dans n'importe quelle direction. L'océan n'est pas une eau stagnante : il est traversé par de véritables fleuves marins.

\begin{definition}[Courant marin]
Un \cle{courant marin} est un déplacement régulier et durable d'une masse d'eau de mer, à la surface ou en profondeur, caractérisé par une direction, une vitesse et une température propres.
\end{definition}

On distingue deux grandes catégories de courants marins :

\begin{itemize}
\item les \cle{courants de surface} (ou courants d'origine éolienne) : ils concernent les 200 à 400 premiers mètres de l'océan et sont engendrés par les vents dominants ;
\item les \cle{courants profonds} (ou courants thermohalins) : ils circulent en profondeur et sont dus aux différences de température et de salinité entre les masses d'eau.
\end{itemize}

\begin{attention}[Ne pas confondre courant et vague]
Une vague est un mouvement oscillatoire de l'eau : l'eau monte et descend mais ne se déplace quasiment pas (un bouchon sur une vague reste sur place). Un courant, lui, est un véritable \emph{transport} d'eau d'une région vers une autre. Au Probatoire, cette confusion est fréquente et coûteuse.
\end{attention}

\courssub{Les causes des courants de surface}

\textbf{Démarche d'investigation.}\par
\emph{Observation} : sur les cartes mondiales, les grands courants de surface suivent des trajectoires qui ressemblent étrangement à celles des vents planétaires étudiés précédemment. \emph{Hypothèse} : les vents dominants entraînent l'eau de surface. \emph{Vérification expérimentale} (réalisable en classe) : on remplit une bassine d'eau, on dépose des copeaux de liège à la surface et l'on souffle régulièrement dans une direction fixe à l'aide d'un ventilateur. Les copeaux se déplacent dans le sens du souffle, puis sont déviés en heurtant les parois. \emph{Interprétation} : le frottement de l'air sur l'eau transfère une partie de l'énergie du vent à la surface liquide ; les obstacles (parois de la bassine, comme les continents dans l'océan) détournent le mouvement. \emph{Conclusion} : les vents planétaires sont bien le moteur principal des courants de surface.

\begin{propriete}[Origine des courants de surface — admise]
Les courants marins de surface sont principalement engendrés par les \cle{vents planétaires} (alizés, vents d'ouest) qui entraînent l'eau par frottement. Leur trajectoire est ensuite déviée par :
\begin{itemize}
\item la \cle{force de Coriolis}, due à la rotation de la Terre (déviation vers la droite dans l'hémisphère Nord, vers la gauche dans l'hémisphère Sud) ;
\item la \cle{forme des continents}, qui bloquent et réorientent les masses d'eau.
\end{itemize}
Cette propriété est admise : sa démonstration physique n'est pas exigible au Probatoire, mais tu dois savoir l'utiliser pour expliquer la direction d'un courant.
\end{propriete}

Ainsi, sous les tropiques, les \emph{alizés} poussent l'eau d'est en ouest : ce sont les courants équatoriaux. Arrivés contre les continents occidentaux des océans, ces courants bifurquent vers les pôles. Aux latitudes moyennes, les \emph{vents d'ouest} ramènent l'eau vers l'est. La boucle est bouclée : c'est la naissance des grands tourbillons océaniques.

\courssub{Courants chauds et courants froids}

\begin{definition}[Courant chaud, courant froid]
Un \cle{courant chaud} est un courant qui transporte de l'eau des basses latitudes (régions chaudes, proches de l'équateur) vers les hautes latitudes ; son eau est plus chaude que celle des régions qu'il traverse. Un \cle{courant froid} transporte de l'eau des hautes latitudes vers les basses latitudes ; son eau est plus froide que celle des régions traversées.
\end{definition}

\begin{attention}[Chaud ou froid : une notion relative]
Un courant est dit chaud ou froid \emph{par rapport à l'eau environnante}, pas en valeur absolue. Le Gulf Stream reste un courant « chaud » même lorsqu'il atteint les côtes de la Norvège avec une eau à \SI{10}{\celsius}, car l'eau environnante y est beaucoup plus froide.
\end{attention}

Conséquence climatique majeure : les courants chauds réchauffent et humidifient les côtes qu'ils longent (climat doux et humide), tandis que les courants froids les refroidissent et stabilisent l'atmosphère, ce qui limite les pluies (climat sec, voire désertique).

\courssub{Les grands courants de l'Atlantique}

L'océan Atlantique, qui borde le Cameroun, offre les exemples les plus célèbres :

\begin{itemize}
\item le \cle{Gulf Stream} : courant chaud né dans le golfe du Mexique, il longe les côtes américaines puis traverse l'Atlantique Nord vers l'Europe occidentale, qu'il adoucit considérablement ;
\item le \cle{courant des Canaries} : courant froid qui descend le long des côtes du Maroc et de l'Afrique de l'Ouest, vers l'équateur ;
\item le \cle{courant de Benguela} : courant froid qui remonte le long des côtes de Namibie et d'Angola, dans l'Atlantique Sud ;
\item le \cle{courant de Guinée} : courant chaud qui coule d'ouest en est dans le golfe de Guinée, baignant directement les côtes camerounaises (Limbé, Kribi, Douala) et contribuant à leur climat chaud et humide.
\end{itemize}

\begin{popfigure}
\begin{center}
\begin{tikzpicture}[scale=0.9,>=Stealth]
% Continents simplifiés
\fill[popgreenL,draw=popgreen!60!black] (-4.6,2.6) .. controls (-3.6,3.4) and (-2.6,3.2) .. (-2.2,2.4) .. controls (-2.6,1.6) and (-2.4,0.8) .. (-2.8,0.2) .. controls (-3.4,-0.8) and (-3.2,-1.8) .. (-3.6,-2.8) .. controls (-4.2,-3.2) and (-4.6,-2.4) .. (-4.4,-1.2) .. controls (-4.8,0.2) and (-4.8,1.6) .. (-4.6,2.6) -- cycle;
\fill[popgreenL,draw=popgreen!60!black] (1.6,2.8) .. controls (2.8,3.2) and (4.2,2.8) .. (4.4,2.0) .. controls (4.2,1.0) and (3.6,0.4) .. (3.4,-0.6) .. controls (3.8,-1.6) and (3.4,-2.6) .. (2.8,-3.0) .. controls (2.2,-2.6) and (2.4,-1.4) .. (2.2,-0.4) .. controls (1.8,0.6) and (1.6,1.8) .. (1.6,2.8) -- cycle;
% Gulf Stream (chaud)
\draw[->,very thick,poporange] (-2.6,0.6) .. controls (-1.6,1.0) and (-0.6,1.4) .. (0.2,2.0) .. controls (0.8,2.4) and (1.2,2.6) .. (1.7,2.7);
\node[poporange,font=\small\bfseries] at (-0.4,1.35) {Gulf Stream (chaud)};
% Canaries (froid)
\draw[->,very thick,popblue] (2.6,2.2) .. controls (2.8,1.4) and (2.8,0.6) .. (2.6,-0.1);
\node[popblue,font=\small\bfseries,align=center] at (3.6,1.1) {Canaries\\(froid)};
% Guinée (chaud)
\draw[->,very thick,poporange] (-2.0,-0.3) .. controls (-0.8,-0.5) and (0.6,-0.5) .. (1.9,-0.4);
\node[poporange,font=\small\bfseries] at (0,-0.85) {Courant de Guinée (chaud)};
% Benguela (froid)
\draw[->,very thick,popblue] (2.7,-2.7) .. controls (2.9,-2.0) and (2.9,-1.3) .. (2.7,-0.7);
\node[popblue,font=\small\bfseries,align=center] at (3.8,-1.8) {Benguela\\(froid)};
% Équateur
\draw[dashed,popmuted] (-4.5,-0.2) -- (4.5,-0.2);
\node[popmuted,font=\scriptsize] at (-4.0,-0.45) {Équateur};
% Labels continents
\node[font=\small\itshape] at (-3.7,1.4) {Amérique};
\node[font=\small\itshape] at (3.1,2.6) {Europe};
\node[font=\small\itshape] at (3.0,0.6) {Afrique};
\node[font=\small\itshape] at (0.3,2.9) {Atlantique Nord};
\node[font=\small\itshape] at (0,-2.4) {Atlantique Sud};
% Cameroun
\fill[poppink] (2.35,-0.35) circle (1.5pt);
\node[poppink,font=\scriptsize\bfseries] at (1.5,-0.15) {Cameroun};
\end{tikzpicture}
\end{center}
\legende{Carte simplifiée de l'océan Atlantique : principaux courants marins. Flèches orange = courants chauds (Gulf Stream, courant de Guinée) ; flèches bleues = courants froids (Canaries, Benguela). Le point rose situe le Cameroun, baigné par le courant chaud de Guinée.}
\end{popfigure}

\begin{exemplebox}[Le courant de Benguela et le désert du Namib]
Le courant froid de \cle{Benguela} remonte des hautes latitudes de l'Atlantique Sud le long des côtes de la Namibie et de l'Angola. En refroidissant l'air côtier, il empêche la formation de nuages de pluie : l'air froid et stable ne s'élève guère, inhibant la convection. Résultat : le désert du Namib, l'un des plus arides du monde, longe pourtant l'océan sur plus de \SI{1500}{km} ! Ce même courant froid, riche en nutriments remontés des profondeurs (phénomène d'upwelling), nourrit en revanche d'immenses bancs de poissons, base d'une pêche abondante en Angola et en Namibie.
\end{exemplebox}

\begin{savaistu}[Un fleuve dans l'océan]
Le Gulf Stream transporte environ 30 millions de mètres cubes d'eau par seconde (30 Sv), soit plus de 100 fois le débit de tous les fleuves du monde réunis ! Sans lui, l'Europe occidentale aurait un climat comparable à celui du Canada à la même latitude. À l'échelle du Cameroun, c'est le courant chaud de Guinée qui maintient la chaleur et l'humidité du littoral, favorables à la mangrove et à la forêt dense côtière.
\end{savaistu}

\courssub{L'organisation en gyres subtropicaux}

En reliant les courants sur une carte, on constate qu'ils dessinent d'immenses boucles appelées \cle{gyres} (ou tourbillons) subtropicaux.

\begin{propriete}[Sens de rotation des gyres — admise]
Sous l'action combinée des alizés, des vents d'ouest, de la force de Coriolis et des continents, les courants de surface forment des gyres subtropicaux qui tournent :
\begin{itemize}
\item dans le \cle{sens horaire} (sens des aiguilles d'une montre) dans l'hémisphère Nord ;
\item dans le \cle{sens antihoraire} dans l'hémisphère Sud.
\end{itemize}
\end{propriete}

\begin{popfigure}
\begin{center}
\begin{tikzpicture}[scale=1.0,>=Stealth]
% Océan
\fill[popblueL,rounded corners=8pt] (-4,-2.6) rectangle (4,2.6);
% Continents
\fill[popgreenL] (-4,-2.6) rectangle (-3.1,2.6);
\fill[popgreenL] (3.1,-2.6) rectangle (4,2.6);
\node[font=\scriptsize\itshape,rotate=90] at (-3.55,0) {continent Ouest};
\node[font=\scriptsize\itshape,rotate=90] at (3.55,0) {continent Est};
% Vents
\draw[->,thick,popmuted,dashed] (-2.8,1.6) -- (2.8,1.6) node[midway,above,font=\scriptsize]{vents d'ouest};
\draw[->,thick,popmuted,dashed] (2.8,-1.4) -- (-2.8,-1.4) node[midway,below,font=\scriptsize]{alizés};
% Gyre (sens horaire)
\draw[->,very thick,poporange] (-2.4,-1.4) .. controls (-2.9,-0.6) and (-2.9,0.6) .. (-2.4,1.6);
\draw[->,very thick,poporange] (-2.4,1.6) .. controls (-1.0,2.0) and (1.0,2.0) .. (2.4,1.6);
\draw[->,very thick,popblue] (2.4,1.6) .. controls (2.9,0.6) and (2.9,-0.6) .. (2.4,-1.4);
\draw[->,very thick,popblue] (2.4,-1.4) .. controls (1.0,-1.8) and (-1.0,-1.8) .. (-2.4,-1.4);
\node[font=\small\bfseries,popdark] at (0,0.15) {Gyre subtropical};
\node[font=\scriptsize,popdark] at (0,-0.35) {(hémisphère Nord : sens horaire)};
% Légende couleurs
\node[poporange,font=\scriptsize] at (-1.6,2.35) {courant chaud};
\node[popblue,font=\scriptsize] at (1.7,-2.3) {courant froid};
\end{tikzpicture}
\end{center}
\legende{Schéma d'un gyre subtropical de l'hémisphère Nord. Les alizés poussent l'eau vers l'ouest (courant équatorial), les continents la dévient vers les pôles (courant chaud, orange), les vents d'ouest la ramènent vers l'est, puis elle redescend vers l'équateur (courant froid, bleu). La boucle tourne dans le sens horaire.}
\end{popfigure}

On retiendra la logique de la boucle : \emph{alizés} $\rightarrow$ courant équatorial vers l'ouest $\rightarrow$ continent $\rightarrow$ courant chaud vers les pôles $\rightarrow$ \emph{vents d'ouest} $\rightarrow$ retour vers l'est $\rightarrow$ courant froid vers l'équateur. Dans l'Atlantique Nord, le gyre associe ainsi le Gulf Stream (chaud) et le courant des Canaries (froid) ; dans l'Atlantique Sud, le gyre antihoraire associe le courant du Brésil (chaud) et le courant de Benguela (froid).

\courssub{La circulation thermohaline : les courants profonds}

\begin{definition}[Circulation thermohaline]
La \cle{circulation thermohaline} est la circulation profonde des océans, provoquée par les différences de \cle{densité} de l'eau de mer, elles-mêmes dues aux différences de température (préfixe \emph{thermo-}) et de salinité (préfixe \emph{-halin}, du grec \emph{hals}, sel).
\end{definition}

Le mécanisme est le suivant : aux hautes latitudes de l'Atlantique Nord, l'eau de surface se refroidit fortement ; sa salinité augmente lors de la formation de la banquise (la glace emprisonne peu de sel, qui enrichit l'eau restante par rejet de saumure). Cette eau froide et salée devient plus dense : elle \emph{plonge} vers les profondeurs (formation de l'eau profonde de l'Atlantique Nord), puis s'écoule lentement vers le sud, traverse les océans, et remonte enfin à la surface dans certaines régions chaudes (upwelling), bouclant un cycle qui dure environ mille ans. L'ensemble est surnommé le \cle{« tapis roulant » océanique} mondial, car il relie tous les océans comme un convoyeur continu.

\begin{attention}[Complément utile, pas cœur du programme]
La circulation thermohaline est un complément précieux pour enrichir une copie de Probatoire, mais l'exigibilité porte surtout sur les courants de surface. Retiens l'essentiel : courants de surface = moteur \emph{éolien} ; courants profonds = moteur \emph{thermohalin} (température + salinité).
\end{attention}

\courssub{Savoir-faire : décrire un courant marin à partir d'une carte}

\begin{methode}[Décrire un courant marin sur une carte]
Face à une carte de courants, rédige ta description en cinq points, dans l'ordre :
\begin{enumerate}
\item \textbf{Nommer} le courant (ex. : courant de Benguela) ;
\item préciser sa \textbf{zone d'origine} (ex. : hautes latitudes de l'Atlantique Sud) ;
\item indiquer sa \textbf{direction} d'écoulement (ex. : du sud vers le nord, le long de la côte occidentale africaine) ;
\item qualifier sa \textbf{température} : chaud (venant des basses latitudes) ou froid (venant des hautes latitudes) ;
\item préciser sa \textbf{zone d'arrivée} et, si possible, une \textbf{conséquence} climatique (ex. : côtes namibiennes froides et arides).
\end{enumerate}
\end{methode}

\begin{exoresolu}[Décrire le courant de Guinée]
\textbf{Énoncé.} À partir de la carte des courants de l'Atlantique, décris le courant de Guinée et indique une conséquence pour le littoral camerounais.
\tcblower
\textbf{Solution détaillée.} Le courant de Guinée est un courant marin de surface qui prend naissance dans l'Atlantique équatorial occidental (basses latitudes). Il s'écoule d'ouest en est en longeant les côtes du golfe de Guinée, jusqu'au littoral du Nigeria et du Cameroun. Venant des régions équatoriales chaudes, c'est un courant chaud. En arrivant sur les côtes camerounaises (Douala, Limbé, Kribi), il réchauffe et humidifie l'air côtier : il contribue ainsi au climat chaud et humide du littoral, marqué par des précipitations abondantes (plus de \SI{4000}{mm} par an à Debundscha, au pied du Mont Cameroun).
\end{exoresolu}

\begin{exercice}[À toi de jouer]
À l'aide de la carte de l'Atlantique, décris le courant des Canaries (nom, origine, direction, température, arrivée) et explique pourquoi les côtes du Maroc et de la Mauritanie qu'il longe sont relativement sèches malgré la proximité de l'océan.
\end{exercice}

\begin{corrige}[Exercice 1]
Le courant des Canaries naît aux latitudes moyennes de l'Atlantique Nord, au large de la péninsule Ibérique. Il s'écoule du nord vers le sud, le long des côtes nord-ouest africaines (Maroc, Sahara occidental, Mauritanie, Sénégal), jusqu'aux îles du Cap-Vert où il rejoint le courant équatorial. Venant des hautes latitudes, c'est un courant froid. En refroidissant l'air côtier, il le rend stable et empêche l'ascendance de l'air humide : les nuages de pluie ne se forment pas, d'où la sécheresse des côtes qu'il longe, prolongeant le désert saharien jusqu'à l'océan.
\end{corrige}

\begin{aretenir}[L'essentiel de la section]
\begin{itemize}
\item Un courant marin est un déplacement régulier et durable d'une masse d'eau de mer, en surface ou en profondeur.
\item Les courants de surface sont engendrés par les vents planétaires, déviés par la force de Coriolis et les continents ; ils s'organisent en gyres : sens horaire au Nord, antihoraire au Sud.
\item Courants chauds : des basses vers les hautes latitudes (Gulf Stream, courant de Guinée) ; courants froids : des hautes vers les basses latitudes (Canaries, Benguela).
\item Les courants profonds forment la circulation thermohaline, le « tapis roulant » océanique, piloté par la température et la salinité.
\item Pour décrire un courant : nom, origine, direction, température, arrivée.
\end{itemize}
\end{aretenir}

Ces courants ne circulent pas indépendamment de l'atmosphère : océan et atmosphère échangent en permanence chaleur et humidité. C'est précisément de ces interactions océan-atmosphère, que nous étudierons dans la section suivante, que naissent les phénomènes climatiques les plus redoutables.

\coursec{Interactions océan-atmosphère}

\courssub{Transition : des courants marins aux échanges globaux}

Dans la section précédente, nous avons décrit la \cle{circulation océanique} comme un vaste tapis roulant mus par le vent et les différences de densité (thermohaline). Nous avons identifié les grands \cle{courants de surface} (chauds ou froids) et les \cle{courants profonds}. Mais l'océan ne se contente pas de transporter de l'eau ; il transporte de l'\textbf{énergie} (chaleur) et de la \textbf{matière} (vapeur d'eau, gaz). C'est au contact de l'atmosphère que ces transferts s'opèrent, modelant le climat de la planète. Observons d'abord une situation concrète vécue au Cameroun.

\begin{experience}[Observation locale : contraste climatique entre Kribi et Douala]
\textbf{Matériel :} Données climatiques moyennes mensuelles (température, pluviométrie) pour Kribi (côte Sud, courant de Guinée chaud) et Douala (côte Littoral, zone de convergence intertropicale influencée par l'océan).
\textbf{Protocole :} Comparer l'évolution mensuelle des températures minimales/maximales et des hauteurs de pluie sur l'année.
\textbf{Observations :}
\begin{itemize}
    \item À Kribi, les températures sont stables (24--28 °C) toute l'année ; les pluies sont abondantes (\textgreater 2500 mm/an) avec deux maxima.
    \item À Douala, les températures sont proches mais l'amplitude thermique diurne est légèrement plus marquée en saison sèche ; les pluies sont très abondantes (\textasciitilde 3600 mm/an) avec une saison des pluies très intense (juillet--août).
\end{itemize}
\textbf{Question :} Pourquoi le littoral camerounais, baigné par le même océan (Golfe de Guinée), présente-t-il des régimes pluviométriques si intenses alors que d'autres côtes tropicales (comme le Nord du Chili ou la Namibie) sont des déserts côtiers ?
\textbf{Hypothèse :} La nature du courant marin qui baigne la côte (chaud vs froid) et la température de surface de l'océan contrôlent l'évaporation et la stabilité de l'air, donc les précipitations.
\textbf{Interprétation :} Le courant de Guinée (chaud) réchauffe l'air au contact de l'eau, favorisant l'évaporation et l'instabilité convective $\rightarrow$ pluies abondantes. Un courant froid (ex: courant de Benguela en Namibie) refroidit l'air basal, créant une inversion de température qui bloque la convection $\rightarrow$ aridité, brouillards.
\textbf{Conclusion :} L'océan n'est pas un simple décor ; sa température de surface (TSM), dictée par les courants, est le \textbf{premier moteur} du temps et du climat côtier.
\end{experience}

\courssub{L'océan : réservoir thermique et moteur de l'évaporation}

\begin{propriete}[Capacité thermique de l'océan -- Admise]
L'océan couvre \textasciitilde 71 \% de la surface terrestre. Sa \cle{capacité thermique massique} ($c_{eau} \approx 4180\,\text{J}\cdot\text{kg}^{-1}\cdot\text{K}^{-1}$) est environ 4 fois celle de l'air ($c_{air} \approx 1005\,\text{J}\cdot\text{kg}^{-1}\cdot\text{K}^{-1}$) et sa masse est \textasciitilde 270 fois plus grande. \textbf{Les premiers 3 mètres de la couche superficielle de l'océan contiennent autant de chaleur que toute l'atmosphère.}
\end{propriete}

\begin{savaistu}[L'océan, « radiateur » de la Terre]
C'est pourquoi l'océan absorbe \textasciitilde 90 \% de l'excès d'énergie piégé par l'effet de serre anthropique. Il \textbf{stocke} cette chaleur en été (ou dans les tropiques) et la \textbf{restitue} en hiver (ou vers les pôles) via les courants. Sans cet amortisseur, les amplitudes thermiques jour/nuit et été/hiver seraient insupportables pour la vie.
\end{savaistu}

\textbf{Mécanisme : le cycle de l'eau océanique.}\par
Le rayonnement solaire (ondes courtes) pénètre les premiers mètres de l'océan. L'eau se réchauffe $\rightarrow$ \cle{évaporation} (changement de phase liquide $\rightarrow$ gaz). Cette évaporation prélève de l'\textbf{énergie latente de vaporisation} ($L_v \approx 2,26 \times 10^6\,\text{J}\cdot\text{kg}^{-1}$) à l'océan : c'est un \textbf{refroidissement} net de la surface océanique. La vapeur d'eau monte, se refroidit, \cle{condense} (nuages, précipitations) et rend cette énergie latente à l'atmosphère : c'est un \textbf{réchauffement} de l'air. \textbf{L'océan alimente ainsi l'atmosphère en chaleur (latente + sensible) et en eau.}

\begin{popfigure}
\begin{tikzpicture}[scale=0.9, every node/.style={font=\small}]
% Couleurs
\colorlet{ocean}{popblue!60!popdark}
\colorlet{atmos}{popblueL!50!white}
\colorlet{fluxR}{poporange}
\colorlet{fluxL}{popred}
\colorlet{fluxS}{popgold}
% Océan
\fill[ocean!30] (-5,-3) rectangle (5,0);
\draw[thick, ocean] (-5,0) -- (5,0) node[midway, above=2pt, black] {Surface océan (TSM)};
% Atmosphère
\fill[atmos] (-5,0) rectangle (5,4);
% Rayonnement solaire
\draw[fluxR, very thick, ->] (0,4.5) -- (0,3.5) node[midway, right] {\textbf{Rayonnement solaire} (ondes courtes)};
\draw[fluxR, thick, ->] (0,3.5) -- (0,0.5) node[midway, right] {Pénétration};
% Évaporation / Chaleur latente
\draw[fluxL, very thick, ->, decorate, decoration={snake, amplitude=1.5pt, segment length=4pt}] (-3,0) -- (-3,2.5) node[midway, left, align=right] {\textbf{Évaporation}\\\textbf{Chaleur latente} ($L_v$)};
% Condensation / Précipitations
\draw[fluxL, very thick, ->, decorate, decoration={snake, amplitude=1.5pt, segment length=4pt}] (3,2.5) -- (3,0) node[midway, right, align=left] {\textbf{Condensation}\\\textbf{Pluies} (restitution $L_v$)};
% Chaleur sensible (conduction/convection)
\draw[fluxS, thick, <->] (-1,0) -- (-1,1.5) node[midway, left] {\textbf{Chaleur sensible}};
% Courant marin
\draw[popgreen, very thick, ->] (4.5,-1.5) -- (4.5,-0.5) node[midway, right] {\textbf{Courant chaud}};
\draw[popgreen, very thick, ->] (-4.5,-1.5) -- (-4.5,-0.5) node[midway, left] {\textbf{Courant froid}};
% Rayonnement infrarouge
\draw[poppurple, thick, <->] (2,0) -- (2,1.5) node[midway, right] {\textbf{IR (serre)}};
% Légendes boîtes
\node[draw, fill=white, rounded corners, align=center, anchor=north] at (0,4.7) {\textbf{ATMOSPHÈRE}\\\textit{Reçoit chaleur \& eau}};
\node[draw, fill=white, rounded corners, align=center, anchor=south] at (0,-3.3) {\textbf{OCÉAN}\\\textit{Stocke \& transporte chaleur}};
\end{tikzpicture}
\legende{Schéma des échanges d'énergie et d'eau à l'interface océan-atmosphère. L'océan absorbe le rayonnement solaire, perd de la chaleur par évaporation (latente) et conduction (sensible), et gagne/perd de la chaleur par le rayonnement infrarouge. Les courants marins redistribuent horizontalement la chaleur stockée.}
\end{popfigure}

\courssub{Influence des courants sur les climats littoraux}

La température de surface de la mer (TSM), imposée par les courants, dicte le climat des côtes voisines.

\begin{propriete}[Effet des courants de surface sur le climat côtier -- Admis]
\begin{itemize}
    \item \textbf{Courant chaud} (ex: Gulf Stream, Kuroshio, Courant de Guinée) : TSM élevée $\rightarrow$ fort flux d'évaporation + air chaud et instable $\rightarrow$ \textbf{climat doux, humide, pluvieux} sur les côtes adjacentes (hivers adoucis, étés frais).
    \item \textbf{Courant froid} (ex: Courant de Californie, Humboldt, Benguela, Canaries) : TSM basse $\rightarrow$ faible évaporation + refroidissement de la couche d'air basale $\rightarrow$ \textbf{inversion de température} (air plus chaud au-dessus d'air froid) $\rightarrow$ stabilité atmosphérique $\rightarrow$ \textbf{climat sec, brumeux, désertique} (déserts côtiers : Atacama, Namib, Sahara occidental).
\end{itemize}
\end{propriete}

\begin{exemplebox}[Le Gulf Stream : le « radiateur » de l'Europe occidentale]
Sans le \cle{Gulf Stream} (et son prolongement, la Dérive Nord-Atlantique), l'Europe de l'Ouest (Paris, Londres, Bergen) aurait un climat \textbf{continental froid}, comparable à celui de \textbf{Terre-Neuve (Canada)} ou de la \textbf{Sibérie méridionale} à la même latitude (45°--60° N).
\begin{itemize}
    \item Le courant transporte \textasciitilde $1,4 \times 10^{15}\,\text{W}$ de chaleur vers le Nord-Est.
    \item L'air maritime réchauffé par le courant apporte douceur et humidité : hivers tempérés (gelées rares à Londres), étés modérés.
    \item \textbf{Preuve paléoclimatique :} Lors du \emph{Younger Dryas} (\textasciitilde 12 900--11 700 ans av. J.-C.), un afflux d'eau douce (fonte glaciaire) a ralenti la circulation thermohaline (AMOC) $\rightarrow$ refroidissement brutal de l'Europe de 5 à 10 °C en quelques décennies.
\end{itemize}
\end{exemplebox}

\begin{attention}[Piège fréquent : « Un courant froid apporte la pluie »]
\textbf{FAUX.} Un courant froid \textbf{stabilise} l'atmosphère (inversion thermique : air froid au sol, air plus chaud en altitude). Il inhibe la convection $\rightarrow$ \textbf{peu ou pas de pluie}. Il favorise les \textbf{brouillards} (condensation au contact de l'air froid) et la \textbf{sécheresse côtière} (déserts de l'Atacama, du Namib). C'est le courant \textbf{chaud} qui, par évaporation intense et instabilité, génère les fortes pluies tropicales (côte du Cameroun, Golfe de Guinée).
\end{attention}

\begin{popfigure}
\begin{tikzpicture}[scale=0.85, every node/.style={font=\footnotesize}]
% Monde simplifié (projection équirectangulaire approximative)
% Axes
\draw[thin, gray] (-9,-4.5) grid (9,4.5);
% Continents (simplifiés)
\fill[popgreen!20!popcream] (-9,1) .. controls (-6,1.5) and (-5,3) .. (-3,4) .. controls (-1,3.5) and (0,2) .. (2,3) .. controls (4,2) and (5,1) .. (9,0.5) -- (9,-4.5) -- (-9,-4.5) -- cycle; % Afrique/AmSud approx
\fill[popgreen!20!popcream] (-9,4.5) -- (-9,2) .. controls (-6,2.5) .. (-3,4.5) -- cycle; % Europe/Asie approx
\fill[popgreen!20!popcream] (-9,4.5) -- (-7,4.5) .. controls (-5,3) .. (-3,4.5); % Groenland/Arctique
\fill[popgreen!20!popcream] (5,-1) .. controls (7,0) and (8,-1) .. (9,-1.5) -- (9,-4.5) -- (5,-4.5) -- cycle; % AmNord approx
% Courants CHAUDS (Rouge/Orange)
\draw[poporange, very thick, ->] (-7,-3) .. controls (-5,-2) and (-3,-1) .. (-1,0) node[above right, poporange, font=\bfseries] {Gulf Stream};
\draw[poporange, very thick, ->] (-1,0) .. controls (1,1) and (3,2) .. (5,3) node[above, poporange, font=\bfseries] {Dérive N.-Atl.};
\draw[poporange, very thick, ->] (6,-2) .. controls (7,-1) .. (8,0.5) node[above right, poporange, font=\bfseries] {Kuroshio};
\draw[poporange, very thick, ->] (-7,-1) .. controls (-5,-0.5) .. (-3,0.5) node[above, poporange, font=\bfseries] {C. Guinée};
\draw[poporange, very thick, ->] (0,-4) .. controls (2,-3) .. (4,-2) node[right, poporange, font=\bfseries] {C. Brésil};
\draw[poporange, very thick, ->] (-4,-4) .. controls (-2,-3.5) .. (0,-3) node[above, poporange, font=\bfseries] {C. Équat.};
% Courants FROIDS (Bleu)
\draw[popblue, very thick, ->] (8,3) .. controls (6,2) and (4,1) .. (2,0) node[below left, popblue, font=\bfseries] {C. Californie};
\draw[popblue, very thick, ->] (2,-1) .. controls (0,-2) and (-2,-3) .. (-4,-3.5) node[below, popblue, font=\bfseries] {C. Canaries};
\draw[popblue, very thick, ->] (-2,-2) .. controls (-4,-3) .. (-6,-3.5) node[below, popblue, font=\bfseries] {C. Benguela};
\draw[popblue, very thick, ->] (-6,1) .. controls (-7,0) .. (-8,-1) node[left, popblue, font=\bfseries] {C. Humboldt};
\draw[popblue, very thick, ->] (4,3) .. controls (2,2) .. (0,1) node[left, popblue, font=\bfseries] {C. Labrador};
% Étiquettes conséquences
\node[draw, fill=poporange!15, rounded corners, align=center, anchor=south] at (-1,3.5) {\textbf{Europe de l'Ouest}\\\textit{Climat tempéré océanique}\\\textit{Hivers doux, pluies régulières}};
\node[draw, fill=popblue!15, rounded corners, align=center, anchor=north] at (6,-2.5) {\textbf{Californie / Pérou}\\\textit{Déserts côtiers / Brumes}\\\textit{(Atacama, Namib)}};
\node[draw, fill=poporange!15, rounded corners, align=center, anchor=south] at (-5,1.5) {\textbf{Golfe de Guinée}\\\textit{(Cameroun)}\\\textit{Climat équatorial humide}};
\node[draw, fill=popblue!15, rounded corners, align=center, anchor=north] at (-7,-2) {\textbf{Namibie / Angola}\\\textit{Désert du Namib}\\\textit{Brouillards, aridité}};
% Légende
\node[draw, fill=white, rounded corners, anchor=south east, align=center] at (8.5,-4){\textbf{Légende :}\\\textcolor{poporange}{\rule{1cm}{2pt}} Courant chaud (adoucit, humidifie)\\\textcolor{popblue}{\rule{1cm}{2pt}} Courant froid (refroidit, assèche)};
\end{tikzpicture}
\legende{Carte simplifiée des principaux courants de surface et de leurs influences climatiques contrastées sur les littoraux. Les courants chauds (orange) adoucissent les climats tempérés et alimentent les pluies tropicales. Les courants froids (bleu) génèrent des déserts côtiers et des brouillards persistants.}
\end{popfigure}

\courssub{Les moussons : respiration saisonnière continent-océan}

\begin{definition}[Mousson]
La \cle{mousson} est une \textbf{inversion saisonnière des vents} à l'échelle planétaire, due au \textbf{contraste thermique} entre un continent et l'océan adjacent. Elle se manifeste par un vent humide venu de la mer en été (pluies) et un vent sec venu du continent en hiver (saison sèche).
\end{definition}

\begin{methode}[Mécanisme de la mousson d'été (ex: Asie du Sud, Afrique de l'Ouest)]
\begin{enumerate}
    \item \textbf{Printemps/Été :} Le continent (ex: sous-continent indien, Sahara) se réchauffe \textbf{rapidement} (faible inertie thermique) $\rightarrow$ surchauffe de l'air $\rightarrow$ \textbf{basse pression thermique} au sol.
    \item L'océan (océan Indien, Atlantique Sud) se réchauffe \textbf{lentement} (forte inertie) $\rightarrow$ air plus frais $\rightarrow$ \textbf{haute pression} relative en surface.
    \item Le gradient de pression pousse l'air maritime (chargé en vapeur d'eau par évaporation intense) vers le continent : \textbf{vent de mousson (Sud-Ouest en hémisphère Nord)}.
    \item Forçage orographique (Himalaya, reliefs du Cameroun/Adamawa) + convergence intertropicale $\rightarrow$ ascendance massive $\rightarrow$ \textbf{condensation, orages, pluies diluviennes}.
\end{enumerate}
\end{methode}

\begin{exemplebox}[La mousson au Cameroun : le « flux de Sud-Ouest »]
Au Cameroun, la mousson est le \textbf{flux de Sud-Ouest} (alizé maritime) qui pénètre depuis le Golfe de Guinée (courant de Guinée \textbf{chaud}) vers le continent africain surchauffé (Sahara, Tchad) entre \textbf{avril et octobre}.
\begin{itemize}
    \item \textbf{Arrivée :} Sud (Kribi, Ebolowa) $\rightarrow$ Mars/Avril. Nord (Maroua, Garoua) $\rightarrow$ Mai/Juin.
    \item \textbf{Retrait :} Nord $\rightarrow$ Septembre/Octobre. Sud $\rightarrow$ Novembre.
    \item \textbf{Conséquence agricole :} Calendrier cultural calé sur l'onde pluviométrique (maïs, coton au Nord ; cacao, café, bananier au Sud).
    \item \textbf{Risque :} Retard ou arrêt précoce $\rightarrow$ \textbf{sécheresse agricole} (ex: années 1970--80, 2010--2011).
\end{itemize}
\end{exemplebox}

\begin{experience}[Modélisation simple : contraste thermique terre/mer]
\textbf{Matériel :} Deux bacs identiques : l'un rempli de sable sec (continent), l'autre d'eau (océan). Deux thermomètres à sonde (ou infrarouge). Lampe puissante (soleil).
\textbf{Protocole :} Exposer les deux bacs à la lampe pendant 30 min. Relever la température du sable et de l'eau toutes les 5 min. Éteindre la lampe et suivre le refroidissement.
\textbf{Résultats attendus :} Le sable monte en T° très vite (faible $c$, pas de mélange) et redescend vite. L'eau monte lentement (forte $c$, mélange convectif) et reste tiède longtemps.
\textbf{Interprétation :} Ce différentiel d'inertie thermique crée le gradient de pression moteur de la mousson. L'océan est le « réservoir » qui alimente le continent en eau quand celui-ci « aspire » l'air.
\end{experience}

\courssub{El Niño / La Niña : l'interaction océan-atmosphère à l'échelle planétaire}

\begin{definition}[ENSO (El Niño - Oscillation Australe)]
L'\cle{ENSO} (\emph{El Niño Southern Oscillation}) est un \textbf{mode de variabilité climatique naturel} couplé océan-atmosphère dans le Pacifique équatorial, de période 2 à 7 ans. Il module les climats mondiaux (téléconnexions).
\end{definition}

\begin{propriete}[Mécanisme simplifié d'El Niño -- Admis]
\begin{enumerate}
    \item \textbf{État normal (La Niña / Neutre) :} \cle{Vents d'Est (alizés)} forts $\rightarrow$ poussent l'eau chaude vers l'Ouest (Indonésie, Australie) $\rightarrow$ \textbf{remontée d'eau froide (upwelling)} à l'Est (côte Pérou/Chili) riche en nutriments $\rightarrow$ pêche abondante. Convection forte à l'Ouest (pluies), subsidence à l'Est (sec).
    \item \textbf{Déclenchement El Niño :} Affaiblissement / inversion des alizés (cause complexe : onde de Kelvin, réchauffement aléatoire). L'eau chaude « rebondit » vers l'Est.
    \item \textbf{État El Niño :} Eau chaude en surface à l'Est (côte Amérique du Sud) $\rightarrow$ \textbf{disparition de l'upwelling} $\rightarrow$ \textbf{effondrement de la pêche} (anchois). Convection se déplace vers l'Est $\rightarrow$ \textbf{pluies torrentielles} (Pérou, Équateur, Californie) + \textbf{sécheresses} à l'Ouest (Indonésie, Australie, Afrique australe, Nordeste brésilien).
    \item \textbf{Retour (La Niña) :} Renforcement des alizés $\rightarrow$ situation inverse (froid intense à l'Est, pluies extrêmes à l'Ouest).
\end{enumerate}
\end{propriete}

\begin{savaistu}[El Niño et le Cameroun : téléconnexions]
Bien que lointain, l'ENSO influence le Cameroun :
\begin{itemize}
    \item \textbf{El Niño} : Tendance à une \textbf{saison des pluies plus faible / retardée} au Nord (Soudano-sahélien) et parfois au Centre (région de l'Adamaoua), lié à un affaiblissement de la mousson africaine.
    \item \textbf{La Niña} : Tendance à des \textbf{pluies excédentaires}, risques d'inondations (ex: inondations de 2010, 2012, 2020 au Nord et Extrême-Nord).
    \item La prévision saisonnière (Météo Cameroun, ACMAD) utilise l'indice ENSO (ONI) pour alerter les agriculteurs et les gestionnaires de barrages (Lagdo, Song Loulou, Memve'ele).
\end{itemize}
\end{savaistu}

\begin{exoresolu}[Analyse d'une carte d'anomalie de TSM (Température de Surface de la Mer) - El Niño 1997-1998]
\textbf{Énoncé :} On donne une carte (simulée ci-dessous) montrant les anomalies de TSM (°C) dans le Pacifique équatorial en décembre 1997 (pic d'un « Super El Niño »). Anomalies positives (rouge) jusqu'à +4°C au large du Pérou ; anomalies négatives (bleu) à l'Ouest (Australie/Indonésie).
\textbf{Questions :}
\begin{enumerate}
    \item Identifier la phase de l'ENSO et justifier par l'anomalie de TSM.
    \item Expliquer la conséquence sur l'upwelling au large du Pérou et sur la pêche à l'anchois.
    \item Prédire le régime pluviométrique à Jakarta (Indonésie) et à Lima (Pérou) ce mois-là.
\end{enumerate}
\tcblower
\textbf{Solution détaillée :}
\begin{enumerate}
    \item \textbf{Phase : El Niño fort.} Justification : Anomalie \textbf{positive} majeure (\textgreater +2°C, ici +4°C) dans le Pacifique \textbf{central et oriental} (région Niño 3.4 et 1+2). Cela signale la présence d'eau anormalement chaude à l'Est, signature du reflux de la « piscine chaude » occidentale.
    \item \textbf{Upwelling :} L'eau chaude en surface (stratification forte) \textbf{bloque la remontée d'eau profonde froide et nutritive}. L'upwelling s'effondre. \textbf{Conséquence :} Disparition du phytoplancton $\rightarrow$ effondrement de la chaîne alimentaire $\rightarrow$ \textbf{crise majeure de la pêche à l'anchois} (pêche la plus importante du monde en tonnage habituellement).
    \item \textbf{Pluviométrie :}
    \begin{itemize}
        \item \textbf{Jakarta (Ouest) :} Zone d'anomalie négative (eau plus froide) + subsidence (air descendant) $\rightarrow$ \textbf{Sécheresse sévère}, risques d'incendies de forêt (tourbières).
        \item \textbf{Lima (Est) :} Zone d'anomalie positive (eau très chaude) + convection intense $\rightarrow$ \textbf{Pluies torrentielles, inondations, glissements de terrain} (phénomène rare dans ce désert côtier hyper-aride).
    \end{itemize}
\end{enumerate}
\end{exoresolu}

\begin{aretenir}[Synthèse : Interactions Océan-Atmosphère]
\begin{itemize}
    \item \textbf{Couplage thermique :} L'océan (réservoir, inertie forte) amortit et redistribue l'énergie solaire ; l'atmosphère (réactivité rapide) transporte la chaleur latente (vapeur) et sensible.
    \item \textbf{Courants de surface = Climat côtier :} Chauds $\rightarrow$ doux/humide (Gulf Stream, Guinée). Froids $\rightarrow$ sec/brumeux (Humboldt, Benguela, Canaries).
    \item \textbf{Mousson :} Moteur = Contraste d'inertie thermique Continent (faible) vs Océan (forte) $\rightarrow$ Inversion saisonnière des vents (Sud-Ouest humide / Nord-Est sec au Cameroun).
    \item \textbf{ENSO :} Oscillation couplée Pacifique équatorial (2-7 ans). El Niño = Eau chaude à l'Est $\rightarrow$ Téléconnexions mondiales (sécheresses Australie/Indonésie/Afrique Australe, pluies Amériques). La Niña = Inverse.
    \item \textbf{Enjeu :} Comprendre ces interactions est vital pour la \textbf{prévision saisonnière}, la \textbf{gestion de l'eau/agriculture} et l'\textbf{adaptation au changement climatique} (réchauffement océan $\rightarrow$ cyclones plus intenses, moussons modifiées, AMOC ralenti ?).
\end{itemize}
\end{aretenir}

\begin{exercice}[Entraînement : Courants et climats]
\begin{enumerate}
    \item En vous aidant de la carte des courants (figure ci-dessus), associez chaque ville à son courant dominant et déduisez son type de climat côtier :
    \begin{itemize}
        \item San Francisco (USA, Californie)
        \item Bergen (Norvège)
        \item Walvis Bay (Namibie)
        \item Rio de Janeiro (Brésil)
    \end{itemize}
    \item Expliquer pourquoi la ville de \textbf{Seattle} (USA, État de Washington, 47° N) a un hiver beaucoup plus doux que \textbf{Québec} (Canada, 46° N) bien qu'elles soient à latitude proche.
    \item Lors d'un épisode \textbf{La Niña} fort, quel impact prévisible sur la campagne agricole cotonnière dans le Nord-Cameroun (région de l'Extrême-Nord) ? Justifier par la téléconnexion.
\end{enumerate}
\end{exercice}

\begin{corrige}[Exercice : Courants et climats]
\begin{enumerate}
    \item
    \begin{center}
    \begin{tabularx}{\linewidth}{|l|X|X|}\hline
    \rowcolor{popblueL} \textbf{Ville} & \textbf{Courant dominant} & \textbf{Climat côtier résultant} \\ \hline
    San Francisco & Courant de Californie (Froid) & Désertique / Méditerranéen sec, brouillards fréquents (stratus), été frais. \\ \hline
    Bergen & Dérive Nord-Atlantique / Gulf Stream (Chaud) & Tempéré océanique très humide, hivers très doux pour la latitude (60° N), étés frais. \\ \hline
    Walvis Bay & Courant de Benguela (Froid) & Désert côtier hyper-aride (Namib), brouillards côtiers quotidiens, quasi pas de pluie. \\ \hline
    Rio de Janeiro & Courant du Brésil (Chaud) & Tropical humide (Aw/Am), hivers doux, étés chauds et pluvieux. \\ \hline
    \end{tabularx}
    \end{center}
    \item \textbf{Seattle} est sous l'influence directe du \textbf{courant chaud du Pacifique Nord (Kuroshio/Dérive Nord-Pacifique)} et des vents d'Ouest qui apportent l'air maritime doux. \textbf{Québec} est sous l'influence du \textbf{courant froid du Labrador} et des vents de Nord-Ouest continentaux (air polaire continental). L'océan Pacifique agit comme un radiateur ; le continent nord-américain + courant froid agissent comme un réfrigérateur.
    \item \textbf{La Niña} $\rightarrow$ Renforcement des alizés / mousson africaine souvent \textbf{renforcée} $\rightarrow$ \textbf{Pluies excédentaires} probables sur le Sahel et le Nord-Cameroun (Soudano-sahélien). \textbf{Conséquence agricole :} Bonnes conditions pour le coton et le mil/sorgho \textbf{SI} la répartition temporelle est bonne. \textbf{Risque :} Excès d'eau $\rightarrow$ inondations (champs, routes, barrages), pourriture des capsules de coton, difficulté de récolte. \textit{Mesure :} Surveillance bulletins ACMAD/Météo Cameroun, variétés à cycle court, drainage.
\end{enumerate}
\end{corrige}

\coursec{Les phénomènes climatiques dangereux : tempêtes et cyclones}

\noindent
Dans la section précédente, nous avons vu comment l'océan et l'atmosphère échangent chaleur, humidité et quantité de mouvement. Ces échanges sont le moteur des grands systèmes météorologiques qui, lorsqu'ils s'intensifient, deviennent des phénomènes dangereux pour les populations. Nous allons maintenant analyser ces systèmes violents : les dépressions tropicales (cyclones, ouragans, typhons) et les tempêtes des latitudes moyennes, en comprenant leur origine dynamique, leur structure et les moyens de s'en protéger.

\courssub{Dépressions, anticyclones et lecture d'une carte météorologique}

\noindent
Sur une carte météorologique, la pression atmosphérique au niveau de la mer est représentée par des \cle{isobares} (lignes reliant les points de même pression). 

\begin{definition}[Dépression et anticyclone]
Une \textbf{dépression} (ou cyclone extratropical) est une zone de pression basse (ex. \SI{990}{\hecto\pascal}) entourée d'isobares fermées ; les vents y tournent dans le sens inverse des aiguilles d'une montre dans l'hémisphère nord (effet Coriolis). Un \textbf{anticyclone} est une zone de pression haute (ex. \SI{1025}{\hecto\pascal}) avec des vents divergents tournant dans le sens des aiguilles d'une montre (hémisphère nord).
\end{definition}

\noindent
\textbf{Observation :} Sur une carte, des isobares serrées indiquent un gradient de pression fort, donc des vents violents. Des isobares espacées signalent un temps calme.

\begin{methode}[Lire une carte de pression au niveau de la mer]
\begin{enumerate}
\item Repérer les centres de basse pression (lettre \textbf{B} ou \textbf{L}) et de haute pression (\textbf{H} ou \textbf{A}).
\item Observer l'écartement des isobares : plus elles sont rapprochées, plus le vent est fort.
\item Déduire le sens de rotation des vents grâce à la règle de Buys‑Ballot (dos au vent, la basse pression est à gauche dans l'hémisphère nord).
\item Identifier les fronts (lignes de discontinuité de température/humidité) souvent associés aux dépressions des latitudes moyennes.
\end{enumerate}
\end{methode}

\courssub{Formation et mécanisme d'un cyclone tropical}

\noindent
Un cyclone tropical est un cas particulier de dépression : il tire son énergie non pas du contraste de masses d'air, mais de la \cle{chaleur latente} libérée par la condensation d'énormes quantités de vapeur d'eau évaporée au-dessus d'un océan très chaud.

\begin{propriete}[Source d'énergie du cyclone tropical -- admise au Probatoire]
L'énergie cinétique d'un cyclone tropical provient quasi exclusivement de la chaleur latente libérée lors de la condensation de la vapeur d'eau évaporée à la surface d'un océan dont la température dépasse \SI{26,5}{\celsius} sur une couche d'au moins \SI{50}{\meter} d'épaisseur.
\end{propriete}

\noindent
\textbf{Démarche d'investigation (guidée)}  
\begin{enumerate}
\item \textbf{Observation :} Les cyclones ne se forment que dans les bandes \SI{5}{\degree}–\SI{25}{\degree} de latitude, au-dessus d'eaux chaudes, jamais à l'équateur.
\item \textbf{Hypothèse :} La chaleur de l'océan alimente le système ; la rotation terrestre (force de Coriolis) est nécessaire pour organiser la convection en tourbillon.
\item \textbf{Données :} Cartes de température de surface de la mer (TSM) et trajectoires historiques montrent une coïncidence stricte entre TSM > \SI{26,5}{\celsius} et zones de genèse.
\item \textbf{Interprétation :} L'évaporation massive fournit de l'air chaud et humide qui s'élève, se condense, libère de la chaleur, abaisse la pression en surface, aspire plus d'air humide → boucle de rétroaction positive.
\item \textbf{Conclusion :} Trois conditions sont indispensables : (1) eau chaude profonde, (2) forte évaporation / instabilité convective, (3) Coriolis non nul (latitude > \SI{5}{\degree}) pour initier la rotation.
\end{enumerate}

\begin{methode}[Relier un phénomène dangereux à son origine dynamique -- trois questions clés]
\begin{enumerate}
\item \textbf{Quelle source d'énergie ?} (chaleur latente pour cyclone, contraste de masses d'air pour tempête frontale).
\item \textbf{Quelle circulation d'air ?} (convergence en surface, ascendance, divergence en altitude pour cyclone ; front chaud/froid avec cisaillement pour tempête).
\item \textbf{Quelles conditions géographiques ?} (océan chaud + latitude suffisante pour cyclone ; zone de conflit polaire‑tropicale pour tempête).
\end{enumerate}
\end{methode}

\begin{attention}[Piège fréquent -- affaiblissement du cyclone]
Un cyclone tropical \textbf{meurt} dès qu'il touche la terre ferme ou une eau plus froide (< \SI{26}{\celsius}) : il est privé de sa source d'énergie (évaporation) et le frottement au sol détruit la circulation tourbillonnaire. Ne pas confondre « affaiblissement » et « dissipation immédiate » ; les pluies diluviennes persistent souvent après l'atterrissage (\emph{landfall}).
\end{attention}

\begin{savaistu}[L'œil du cyclone]
L'œil, zone de calme apparent de \SI{20}{\kilo\meter} à \SI{50}{\kilo\meter} de diamètre, est entouré du \textbf{mur de l'œil} où soufflent les vents les plus violents (pouvant dépasser \SI{250}{\kilo\meter\per\hour}). La pression y est minimale (parfois < \SI{900}{\hecto\pascal}). L'air y descend lentement, se réchauffe par compression adiabatique, ce qui explique l'absence de nuages.
\end{savaistu}

\begin{popfigure}
\begin{tikzpicture}[scale=0.9, every node/.style={font=\small}]
% Earth surface
\draw[fill=popblue!20] (-6,-1) rectangle (6,0);
\draw[thick] (-6,0) -- (6,0) node[midway,below=2pt] {Surface de l'océan (TSM > \SI{26,5}{\celsius})};
% Eye
\draw[fill=popcream] (0,0) ellipse (1.2 and 0.4);
\draw[thick] (0,0) ellipse (1.2 and 0.4);
\node[align=center] at (0,0.2){\textbf{ŒIL}\\(calme, air descendant)};
% Eyewall
\draw[fill=poporange!30] (-1.2,0) .. controls (-2,2) and (2,2) .. (1.2,0) -- cycle;
\draw[thick] (-1.2,0) .. controls (-2,2) and (2,2) .. (1.2,0);
\node[align=center] at (0,1.6){\textbf{MUR DE L'ŒIL}\\(vents max, forte pluie)};
% Updrafts
\foreach \x in {-1.5,-0.5,0.5,1.5} {
  \draw[->,thick,popgreen] (\x,0) -- (\x,3.5) node[above] {Ascendance};
}
% Outflow at top
\draw[->,thick,poporange] (-4,3.5) -- (-5,4) node[left] {Divergence en altitude};
\draw[->,thick,poporange] (4,3.5) -- (5,4) node[right] {Divergence en altitude};
% Inflow at surface
\draw[->,thick,popblue] (-5,-0.5) -- (-2,-0.5) node[midway,below] {Alimentation (alizés)};
\draw[->,thick,popblue] (5,-0.5) -- (2,-0.5) node[midway,below] {Alimentation (alizés)};
% Spiral rainbands
\foreach \r in {2.5,3.5,4.5} {
  \draw[popblue!60,thick,domain=0:360,smooth,variable=\t] plot ({\r*cos(\t)*0.3},{\r*sin(\t)*0.2+0.5});
}
\node at (0,-2.5) {Coupe verticale schématique d'un cyclone tropical};
\end{tikzpicture}
\legende{Coupe verticale d'un cyclone tropical : œil central (air descendant), mur de l'œil (ascendances violentes, vents maximaux), bandes pluvieuses spiralées, divergence en altitude et convergence des alizés en surface.}
\end{popfigure}

\courssub{Structure d'un cyclone tropical : œil, mur de l'œil, bandes pluvieuses}

\noindent
Le cyclone mature présente une structure tripartite bien visible sur les images satellites et les radars :

\begin{itemize}
\item \textbf{Œil (eye)} : zone circulaire de \SI{20}{\kilo\meter}–\SI{50}{\kilo\meter} de diamètre, ciel dégagé, vent faible, pression minimale. L'air y descend, se réchauffe adiabatiquement, inhibant la convection.
\item \textbf{Mur de l'œil (eyewall)} : anneau de nuages convectifs (cumulonimbus) de \SI{10}{\kilo\meter}–\SI{20}{\kilo\meter} d'épaisseur. C'est là que se produisent les ascendances les plus intenses (jusqu'à \SI{20}{\meter\per\second}), les précipitations les plus fortes et les vents maximaux (souvent > \SI{118}{\kilo\meter\per\hour}, catégorie 1 sur l'échelle de Saffir‑Simpson).
\item \textbf{Bandes pluvieuses spiralées (rainbands)} : rangées de nuages convectifs s'enroulant vers le centre, apportant des pluies torrentielles et des rafales de vent sur des centaines de kilomètres.
\end{itemize}

\courssub{Vocabulaire régional : cyclone, ouragan, typhon}

\noindent
Il s'agit du \textbf{même phénomène physique} ; seul le nom change selon le bassin océanique :

\begin{center}
\begin{tabularx}{\linewidth}{|l|X|}\hline
\rowcolor{popblueL}
\textbf{Bassin} & \textbf{Appellation} \\ \hline
Océan Indien Sud‑Ouest (dont canal de Mozambique) & Cyclone tropical \\ \hline
Atlantique Nord, Pacifique Nord‑Est & Ouragan (hurricane) \\ \hline
Pacifique Nord‑Ouest & Typhon (typhoon) \\ \hline
Pacifique Sud‑Ouest & Cyclone tropical \\ \hline
\end{tabularx}
\end{center}

\courssub{Tempêtes des latitudes moyennes : fronts et cisaillement}

\noindent
Contrairement aux cyclones tropicaux, les \textbf{tempêtes extratropicales} (ou dépressions des latitudes moyennes) naissent au niveau du \textbf{front polaire}, zone de rencontre entre une masse d'air polaire froide et une masse d'air tropicale chaude.

\begin{experience}[Modèle conceptuel de la genèse d'une dépression frontale -- admis]
\begin{itemize}
\item \textbf{Milieu initial :} Front stationnaire séparant air froid (nord) et air chaud (sud).
\item \textbf{Perturbation :} Une ondulation (onde frontale) se forme sous l'effet du jet‑stream.
\item \textbf{Évolution :} L'air chaud remonte au nord (front chaud), l'air froid descend au sud (front froid). La pression chute au centre de l'onde → dépression.
\item \textbf{Maturité :} Le front froid rattrape le front chaud → occlusion. La dépression atteint son intensité maximale (vents forts, pluies continues).
\item \textbf{Dissipation :} L'air froid entoure totalement le centre ; la source d'énergie (contraste thermique) s'épuise.
\end{itemize}
\end{experience}

\noindent
Ces tempêtes apportent des vents violents (souvent \SI{80}{\kilo\meter\per\hour}–\SI{120}{\kilo\meter\per\hour}), des précipitations durables et, en hiver, de la neige ou du verglas.

\courssub{Conséquences et dangers pour les populations}

\noindent
Qu'il s'agisse d'un cyclone tropical ou d'une tempête frontale, les impacts majeurs sont :

\begin{itemize}
\item \textbf{Vents destructeurs} (> \SI{118}{\kilo\meter\per\hour} pour un cyclone, rafales > \SI{150}{\kilo\meter\per\hour}) : arrachage de toitures, chute d'arbres, projectiles.
\item \textbf{Pluies diluviennes} (\SI{200}{\milli\meter}–\SI{500}{\milli\meter} en \SI{24}{\hour}) : crues éclaires, inondations urbaines et rurales.
\item \textbf{Ondes de tempête (surcotes marines)} : élévation du niveau de la mer de \SI{2}{\meter} à \SI{6}{\meter} (pression basse + vent de terre) submergeant les côtes basses.
\item \textbf{Glissements de terrain et coulées de boue} : saturation des sols sur les versants (ex. région de l'Ouest camerounais, Mont Cameroun).
\item \textbf{Risques sanitaires post‑événement} : épidémies de choléra, paludisme, blessures, stress psychologique.
\end{itemize}

\begin{exemplebox}[Cyclone Idai — Mozambique, mars 2019]
Le cyclone Idai (catégorie 3) a touché Beira avec des vents de \SI{185}{\kilo\meter\per\hour} et des pluies > \SI{300}{\milli\meter}. L'onde de tempête a submergé 90 \% de la ville. Bilan : plus de 1 000 morts, 2,2 millions de sinistrés, destruction des infrastructures sanitaires et agricoles. La majorité des décès sont dus aux inondations et aux coulées de boue dans les zones rurales mal drainées.
\end{exemplebox}

\courssub{Relier un phénomène dangereux à son origine dynamique -- application}

\begin{exoresolu}[Analyse d'une situation météorologique]
\textbf{Énoncé :} Une carte de pression au 12 h UTC montre une dépression centrée à 15° N, 60° O avec isobares serrées (écartement 2 hPa/100 km). La température de surface de la mer est de \SI{28}{\celsius}. Un front froid s'étend du Texas vers le golfe du Mexique. Identifiez la nature de chaque système et expliquez l'origine dynamique de leurs vents violents.

\textbf{Solution détaillée :}
\begin{enumerate}
\item \textbf{Système tropical (15° N, 60° O)} : 
\begin{itemize}
\item \textit{Source d'énergie :} chaleur latente (TSM = \SI{28}{\celsius} > \SI{26,5}{\celsius}) → évaporation massive.
\item \textit{Circulation :} convergence des alizés en surface, forte ascendance dans le mur de l'œil, divergence en altitude.
\item \textit{Condition géographique :} latitude 15° N → Coriolis suffisant pour organiser la rotation.
\item \Rightarrow \textbf{Cyclone tropical (ouragan)}.
\end{itemize}
\item \textbf{Système frontal (golfe du Mexique)} :
\begin{itemize}
\item \textit{Source d'énergie :} contraste thermique air polaire / air tropical (gradient de température horizontal).
\item \textit{Circulation :} cisaillement vertical du vent (jet‑stream), front chaud et front froid, occlusion naissante.
\item \textit{Condition géographique :} latitudes moyennes (30°–40° N), zone de conflit des masses d'air.
\item \Rightarrow \textbf{Tempête extratropicale (dépression frontale)}.
\end{itemize}
\end{enumerate}
\noindent Les vents violents du cyclone viennent de la libération de chaleur latente ; ceux de la tempête frontale du gradient de pression horizontal intense lié au contraste de masse d'air.
\end{exoresolu}

\courssub{Prévention et protection face aux phénomènes dangereux}

\noindent
La réduction du risque repose sur trois piliers : \textbf{surveillance}, \textbf{préparation} et \textbf{réaction}.

\begin{methode}[Mesures de protection -- niveau individuel et collectif]
\begin{enumerate}
\item \textbf{Surveillance et alerte précoce} : suivre les bulletins de la météo nationale (ex. \emph{Direction de la Météorologie Nationale} au Cameroun) et les alertes cycloniques (niveaux jaune, orange, rouge). Utiliser les applications d'alerte SMS.
\item \textbf{Aménagement du territoire} : interdire l'habitat en zone inondable (plaines alluviales du Wouri, bas‑fond de Yaoundé), préserver les mangroves et forêts côtières qui amortissent l'onde de tempête.
\item \textbf{Construction résiliente} : toitures arrimées, ouvertures protégées (volets cycloniques), fondations surélevées en zone côtière.
\item \textbf{Plans d'urgence familiaux et communautaires} : point de rassemblement, kit de survie (eau, nourriture, médicaments, radio à piles), évacuation vers des abris officiels (écoles, mairies renforcées).
\item \textbf{Gestion post‑catastrophe} : drainage rapide, surveillance épidémiologique (choléra, paludisme), soutien psychosocial.
\end{enumerate}
\end{methode}

\begin{attention}[Erreur à ne pas commettre lors du Probatoire]
Ne pas confondre \textbf{cyclone tropical} (énergie = chaleur latente, symétrique, œil) et \textbf{tempête frontale} (énergie = contraste thermique, asymétrique, fronts). La question « Quelle est l'origine dynamique des vents violents ? » attend une réponse précise selon le système.
\end{attention}

\begin{popfigure}
\begin{tikzpicture}[scale=0.55, every node/.style={font=\tiny}]
% World map simplified: equator, tropics, continents outlines
\draw[popmuted] (-180,-90) rectangle (180,90);
% Equator and tropics
\draw[popblue!50,dashed] (-180,0) -- (180,0) node[right] {Équateur};
\draw[popblue!50,dashed] (-180,23.5) -- (180,23.5) node[right] {Tropique Cancer};
\draw[popblue!50,dashed] (-180,-23.5) -- (180,-23.5) node[right] {Tropique Capricorne};
% 5° and 25° latitude bands
\draw[poporange!70] (-180,5) -- (180,5);
\draw[poporange!70] (-180,25) -- (180,25);
\draw[poporange!70] (-180,-5) -- (180,-5);
\draw[poporange!70] (-180,-25) -- (180,-25);
% Continents rough shapes
\draw[fill=popgreen!30,draw=popgreen!80] (-120,60) .. controls (-100,50) and (-80,55) .. (-60,60) .. controls (-40,65) and (-20,60) .. (0,65) .. controls (20,60) and (40,65) .. (60,60) .. controls (80,55) and (100,50) .. (120,60) -- (120,90) -- (-120,90) -- cycle; % North America / Eurasia top
\draw[fill=popgreen!30,draw=popgreen!80] (-80,-60) .. controls (-60,-50) and (-40,-55) .. (-20,-60) .. controls (0,-65) and (20,-60) .. (40,-60) .. controls (60,-55) and (80,-50) .. (100,-60) -- (100,-90) -- (-80,-90) -- cycle; % South America / Africa bottom
% Cyclone genesis zones (hatched)
\foreach \lat in {10,15,20} {
  \foreach \lon in {-150,-120,-90,-60,-30,0,30,60,90,120,150} {
    \fill[poporange!30] (\lon-5,\lat-2.5) rectangle (\lon+5,\lat+2.5);
    \fill[poporange!30] (\lon-5,-\lat-2.5) rectangle (\lon+5,-\lat+2.5);
  }
}
% Typical tracks (arrows)
\draw[->,thick,poppurple] (-150,10) .. controls (-120,15) and (-90,20) .. (-60,25);
\draw[->,thick,poppurple] (-120,-10) .. controls (-90,-15) and (-60,-20) .. (-30,-25);
\draw[->,thick,poppurple] (30,10) .. controls (60,15) and (90,20) .. (120,25);
\draw[->,thick,poppurple] (60,-10) .. controls (90,-15) and (120,-20) .. (150,-25);
% Labels
\node[poporange] at (0,30) {Zones de genèse (5°–25° N/S)};
\node[poppurple] at (0,-30) {Trajectoires types (ouest → nord/est)};
\end{tikzpicture}
\legende{Carte schématique mondiale des zones de formation des cyclones tropicaux (bandes hachurées entre 5° et 25° de latitude nord et sud) et de leurs trajectoires typiques (flèches violettes). Les cyclones naissent au-dessus d'eaux chaudes (> \SI{26,5}{\celsius}) et sont entraînés vers l'ouest par les alizés, puis recourbent vers les pôles sous l'effet du flux de moyenne altitude.}
\end{popfigure}

\noindent
En résumé, les phénomènes climatiques dangereux — cyclones tropicaux et tempêtes des latitudes moyennes — trouvent leur origine dans des mécanismes dynamiques distincts mais tous deux liés aux grands mouvements de l'atmosphère et de l'hydrosphère étudiés précédemment. Comprendre ces mécanismes permet d'anticiper les risques, d'alerter les populations et de mettre en œuvre des stratégies d'adaptation indispensables face au changement climatique qui tend à intensifier ces événements.

\coursec{Prévention et protection face aux phénomènes dangereux}

Après avoir analysé la genèse et la dynamique des tempêtes et des cyclones, nous devons maintenant répondre à une question vitale : \textbf{comment protéger les populations et les biens face à ces phénomènes inévitables ?} L'histoire récente du Cameroun, comme les inondations dévastatrices de Douala en 2020 ou les glissements de terrain meurtriers dans l'Ouest (Bafoussam, 2019), nous rappelle que la connaissance scientifique ne suffit pas ; elle doit se traduire en actions concrètes de prévision, d'alerte, d'aménagement du territoire et d'éducation des citoyens. C'est l'objet de cette dernière section.

\courssub{La prévision météorologique : l'œil scientifique sur l'atmosphère et l'océan}

\emph{Observation :} Avant l'ère spatiale, les prévisions ne dépassaient guère 24 heures et reposaient sur des observations de surface éparses. Aujourd'hui, nous suivons la naissance d'une dépression tropicale au large du Cap-Vert une semaine avant qu'elle n'atteigne les Antilles ou les côtes africaines. Comment est-ce possible ?

\emph{Hypothèse :} La prévision moderne repose sur l'assimilation de millions de données d'observation dans des modèles numériques de la circulation atmosphérique et océanique (les équations de la dynamique des fluides appliquées à la Terre).

\begin{experience}[La chaîne d'observation et de modélisation météorologique]
\textbf{Matériel (conceptuel) :} Satellites géostationnaires (ex. Meteosat pour l'Afrique) et polaires, bouées dérivantes et ancrées (réseau Argo, DBCP), radars météorologiques au sol (bande C ou S), radiosondages (ballons-sondes), stations au sol (synop), avions de reconnaissance (pour les cyclones).
\textbf{Protocole :}
\begin{enumerate}
    \item \textbf{Observation globale continue :} Les satellites fournissent des images dans le visible, l'infrarouge et le micro-ondes (température de surface de la mer, couverture nuageuse, profil de température/humidité, vents de surface par diffusométrie). Les bouées mesurent la pression, la température de l'eau, le vent, les vagues. Les radars détectent la pluviométrie et la vitesse des hydrométéores (effet Doppler) à haute résolution (quelques km, 5 min).
    \item \textbf{Assimilation de données (analyse) :} Toutes ces observations hétérogènes sont injectées dans un « premier coup d'œil » (arrière-plan) issu du modèle précédent pour produire l'état initial le plus réaliste possible de l'atmosphère et de l'océan (analyse 3D-Var ou 4D-Var / EnKF).
    \item \textbf{Intégration numérique :} Le modèle (ex. ARPEGE de Météo-France, IFS de l'ECMWF, GFS américain, WRF pour la haute résolution locale) résout les équations primitives (Navier-Stokes, thermodynamique, continuité, état) sur une grille 3D (horizontal ~9-13 km global, ~1-3 km régional ; vertical ~90 niveaux). Les paramétrisations représentent les processus sous-maille (convection, rayonnement, couche limite, microphysique des nuages).
    \item \textbf{Prévision d'ensemble (EPS) :} On lance 50 à 100 simulations avec des conditions initiales légèrement perturbées et des paramétrages stochastiques pour quantifier l'incertitude (probabilité d'occurrence d'un vent > 100 km/h, pluviométrie > 100 mm/24h).
    \item \textbf{Produits pour les prévisionnistes :} Cartes de géopotentiel, vents, température, pluviométrie, risque orageux, trajectoires cycloniques (cônes d'incertitude), bulletins d'alerte.
\end{enumerate}
\textbf{Observations :} La prévisibilité déterministe est limitée à ~7-10 jours (théorème de Lorenz : chaos atmosphérique). Au-delà, seule l'approche probabiliste (ensemble) a du sens. La prévision des cyclones (trajectoire) a gagné ~1 jour de précision par décennie ; l'intensité reste plus difficile. Les phénomènes convectifs violents (orages stationnaires causant crues éclairs) restent difficiles à localiser précisément > 6-12h à l'avance.
\textbf{Interprétation :} La prévision est une chaîne : \textbf{Observation $\rightarrow$ Assimilation $\rightarrow$ Modélisation $\rightarrow$ Expertise humaine $\rightarrow$ Diffusion}. La rupture à n'importe quel maillon (panne satellite, modèle biaisé, communication défaillante) dégrade le service rendu à la population.
\textbf{Conclusion :} La prévision météorologique moderne est une prouesse technologique et scientifique internationale (coordination OMM), mais elle reste probabiliste. Son utilité sociale dépend entièrement de la chaîne d'alerte qui suit.
\end{experience}

\begin{popfigure}
\begin{tikzpicture}[
    node distance=1.2cm and 1.5cm,
    box/.style={draw=popdark, thick, rounded corners, fill=popcream, text width=2.8cm, align=center, minimum height=1.8cm, font=\small},
    arrow/.style={-Stealth, thick, popblue},
    labelnode/.style={font=\footnotesize\itshape, text=popmuted, align=center}
]
% Nodes
\node[box, align=center] (obs){Observation\\Satellites, bouées,\\radars, radiosondes,\\stations au sol};
\node[box, right=of obs, align=center] (assim){Assimilation\\ \& Analyse\\État initial 3D\\de l'atmosphère/\\océan};
\node[box, right=of assim, align=center] (model){Modélisation\\Numérique\\Intégration des équations\\Prévision d'ensemble (EPS)};
\node[box, right=of model, align=center] (expert){Expertise\\Prévisionniste\\Interprétation,\\validation, localisation};
\node[box, right=of expert, align=center] (alert){Alerte \& Diffusion\\Vigilance, bulletins,\\SMS, médias, sirènes};
\node[box, below=of alert, xshift=-3.5cm, align=center] (protect){Protection\\Civile\\Évacuation, secours,\\mise à l'abri};

% Arrows
\draw[arrow] (obs) -- (assim) node[midway, above, labelnode] {Données brutes};
\draw[arrow] (assim) -- (model) node[midway, above, labelnode] {Champs analysés};
\draw[arrow] (model) -- (expert) node[midway, above, labelnode, align=center]{Sorties brutes\\(probabilistes)};
\draw[arrow] (expert) -- (alert) node[midway, above, labelnode, align=center]{Bulletins\\validés};
\draw[arrow] (alert) -- (protect) node[midway, right, labelnode, align=center]{Ordres,\\consignes};

% Feedback loop
\draw[arrow, poporange, dashed] (protect.west) -- ++(-1,0) |- (obs.west) node[midway, left, labelnode, align=center]{Retour\\d'expérience\\(RETEX)};
\end{tikzpicture}
\legende{La chaîne d'alerte météorologique : de l'observation à la protection des populations. Chaque maillon doit être robuste ; la faiblesse d'un seul (ex. : radars en panne, population non informée) fait échouer le système. Le retour d'expérience (RETEX) après chaque événement améliore le maillon suivant.}
\end{popfigure}

\begin{definition}[Vigilance météorologique (échelle nationale)]
Système d'information du public et des autorités sur l'évolution d'une situation météorologique dangereuse. Au Cameroun, la \textbf{Direction de la Météorologie Nationale (DMN)}, placée sous la tutelle du Ministère des Transports (MINT), émet des bulletins de vigilance (couleurs : vert, jaune, orange, rouge) basés sur les seuils de vent, pluie, orage, vague-submersion.
\end{definition}

\begin{attention}[Piège fréquent : Confondre prévision et certitude]
Une prévision « orage fort probable (80\%) » ne signifie pas qu'il y \emph{aura} un orage chez vous, ni qu'il n'y \emph{aura pas} d'orage si la probabilité est 30\%. Au Probatoire, on attend que vous distinguiez \textbf{prévision déterministe} (une valeur : 45 mm de pluie) et \textbf{prévision probabiliste} (risque de dépasser 50 mm : 60\%). Les décisions de protection (évacuation) se basent sur l'analyse du risque (aléa $\times$ vulnérabilité $\times$ enjeu), pas sur une certitude illusoire.
\end{attention}

\courssub{Les systèmes d'alerte : de l'information à l'action}

\emph{Problème :} Une prévision parfaite est inutile si elle n'atteint pas le dernier kilomètre : le citoyen en zone de risque, souvent illettré, sans smartphone, ou méfiant envers les autorités.

\begin{methode}[Organiser un système d'alerte précoce efficace (norme OMM/CAP)]
\textbf{Quatre piliers indissociables (Cadre de Sendai, OMM) :}
\begin{enumerate}
    \item \textbf{Connaissance des risques :} Cartographie des aléas (zones inondables, couloirs d'éboulement, zones de submersion marine) et des enjeux (population, infrastructures critiques : hôpitaux, écoles, réseaux eau/électricité). \emph{Au Cameroun : carte des zones à risque de Douala (quartiers Makèpè, Bépanda, New-Bell), Yaoundé (Mvog-Ada, Melen), Ouest (versants du Mont Bamboutos).}
    \item \textbf{Service d'alerte et de prévision :} Capacité technique (DMN) à détecter, prévoir et qualifier le phénomène (localisation, intensité, horaire, probabilité). Standardisation des messages (format CAP - Common Alerting Protocol).
    \item \textbf{Diffusion et communication :} Canaux redondants : radio/TV nationale (CRTV), radios communautaires (langues locales), SMS Cell Broadcast (alerte géolocalisée), sirènes (usines, mairies), relais communautaires (chefs de quartier, leaders religieux, ASC - Agents de Santé Communautaires). \textbf{Mot-clé :} \cle{accessibilité} (langue, format, handicap).
    \item \textbf{Capacité de réaction :} Plans d'urgence communaux (PUC), exercices de simulation (simulacres), identification des sites d'hébergement d'urgence (écoles, églises, stades), prépositionnement des secours (Protection Civile, Sapeurs-Pompiers, Croix-Rouge, forces de défense).
\end{enumerate}
\textbf{Critère de succès :} Le « dernier kilomètre » : le temps entre la réception de l'alerte et la mise en sécurité effective doit être inférieur au temps d'arrivée du phénomène (délai de mise en sécurité < délai d'arrivée).
\end{methode}

\begin{exemplebox}[Gestion d'une alerte crue à Douala - Scénario type]
\emph{Situation :} Une perturbation équatoriale stationnaire annonce 150 mm en 24h sur le bassin du Wouri (vigilance ORANGE pluies).
\textbf{Chronologie idéale :}
\begin{itemize}
    \item \textbf{J-1 (18h) :} DMN émet Bulletin de Vigilance Orange. Mairie de Douala active le Centre Opérationnel Communal (COC).
    \item \textbf{J-1 (20h) :} Message radio/TV en français, douala, bassa : « Risque crue Wouri. Quartiers riverains : préparez kits survie, repérez point haut. »
    \item \textbf{J (06h) :} Pluies intenses débutent. Radars/Bouées confirment. Vigilance passée ROUGE. Ordre d'évacuation préventive zones basses (Makèpè, Akwa-Nord). Sirènes + porte-à-porte ASC.
    \item \textbf{J (10h) :} Eau monte. Population rejoint sites d'hébergement (stade, lycées). Coupure préventive électricité/eau zones inondées (ENEO/Camwater).
    \item \textbf{J+1 :} Décrue. Évaluation dégâts. Retour progressif. RETEX dans les 15 jours.
\end{itemize}
\emph{Point critique réel :} Souvent, l'ordre d'évacuation arrive trop tard (eau déjà à 50 cm) ou les sites d'hébergement manquent de latrines/eau potable $\rightarrow$ risque choléra post-inondation.
\end{exemplebox}

\courssub{Mesures de protection structurelles : aménager le territoire pour réduire la vulnérabilité}

\emph{Observation :} Face à un cyclone ou une crue centennale, la résistance pure (murs épais) a des limites économiques et physiques. La stratégie moderne combine \textbf{résistance} (renforcer l'existant), \textbf{évitement} (ne pas construire en zone rouge) et \textbf{résilience} (accepter l'eau mais limiter les dégâts).

\begin{propriete}[Mesures structurelles (ou « grises » et « vertes »)]
\begin{itemize}
    \item \textbf{Protection côtière hybride :} Digues/épis en béton/roches (\emph{grises}) + \textbf{Mangroves et forêts littorales} (\emph{vertes/bleues}). La mangrove dissipe l'énergie des vagues (coefficient de transmission $K_t < 0,3$ pour 100 m de largeur), piège les sédiments, stocke le carbone, et sert de nurserie halieutique.
    \item \textbf{Urbanisme résilient :} \textbf{Zone de retrait (setback line)} : interdiction de construire sur une bande littorale (ex. 80-100 m au Cameroun selon loi 96/12). Surélévation des habitations (pilotis, rez-de-chaussée non habité), toitures renforcées (crochets anti-cycloniques, pente $> 30^\circ$), ouvertures protégées (volets, verre feuilleté).
    \item \textbf{Gestion des eaux pluviales urbaines :} Caniveaux dimensionnés (période de retour 10-20 ans), bassins de rétention/détention (stockage temporaire pour lisser le pic de crue), désimperméabilisation (noues, pavés drainants, toitures végétalisées).
    \item \textbf{Stabilisation des versants :} Drainage profond (puits, drains horizontaux), paratonnerres/parablocs, revêtements végétaux (vetiver, arbres à racines pivotantes), terrasses culturales (agriculture de pente).
\end{itemize}
\end{propriete}

\begin{popfigure}
\begin{tikzpicture}[
    scale=0.9,
    every node/.style={font=\small},
    layer/.style={draw=popdark, thick, fill=popcream, rounded corners=3pt, minimum height=1.2cm},
    water/.style={fill=popblue!30, draw=popblue, thick},
    veg/.style={fill=popgreen!40, draw=popgreen!60!black, thick, pattern=north east lines, pattern color=popgreen!60!black},
    sand/.style={fill=popgold!40, draw=poporange, thick},
    constru/.style={fill=poppink!30, draw=poppurple, thick, minimum width=1.5cm, minimum height=1cm},
    arrow/.style={-Stealth, poporange, thick},
    label/.style={font=\footnotesize, text=popdark, align=center}
]
% Mer
\fill[water] (0,0) rectangle (14,2);
\foreach \x in {0.5,1.5,...,13.5} {
    \draw[popblue!50, decorate, decoration={snake, amplitude=0.1cm, segment length=0.5cm}] (\x,0.5) -- (\x+0.3,0.5);
}
\node[label] at (7,1) {Océan Atlantique / Estuaire du Wouri};

% Mangrove (barrière naturelle)
\fill[veg] (0,2) rectangle (3,4);
\node[label, popgreen!50!black, align=center] at (1.5,3){MANGROVE\\(Avacennia, Rhizophora)\\Dissipation houle $\approx$ 70-90\%\\Piégeage sédiments};

% Zone tampon / Digue végétalisée
\fill[sand] (3,2) rectangle (5,4);
\draw[popgreen!60!black, thick] (3,2) -- (5,2);
\node[label, align=center] at (4,3){Zone tampon\\Végétation halophile\\Dune fixe};

% Digue de protection (structure grise)
\draw[fill=popmuted!30, draw=popdark, thick] (5,2) -- (5.3,2) -- (5.3,5) -- (5,5) -- cycle;
\node[label, rotate=90, align=center] at (4.7,3.5){DIGUE / ÉPI\\Béton / Enrochement};

% Zone de retrait (non constructible)
\fill[popcream] (5.3,2) rectangle (8,4);
\draw[dashed, poporange, thick] (5.3,2) -- (8,2);
\node[label, poporange, align=center] at (6.65,3){ZONE DE RETRAIT\\(80-100 m)\\INTERDICTION DE CONSTRUIRE\\(Loi 96/12, Code urbanisme)};

% Zone urbanisée résiliente (sur pilotis)
\foreach \x in {8.5, 10.5, 12.5} {
    \node[constru, align=center] (h\x) at (\x,3.5){Maison\\sur pilotis};
    \draw[popdark, thick] (\x-0.6,2) -- (\x-0.6,3) -- (\x+0.6,3) -- (\x+0.6,2);
    \node[label, align=center] at (\x,2.3){Rez-de-chaussée\\libre (stockage,\\parking)};
    \draw[arrow] (\x,4.5) -- (\x,5) node[above, label, align=center]{Toiture\\renforcée};
}

% Flèches d'interaction
\draw[arrow, decorate, decoration={snake, amplitude=0.15cm}] (2.5,1) -- (2.5,2) node[midway, right, label, align=center]{Houle cyclonique\\atténuée};
\draw[arrow] (5.15,3.5) -- (5.3,3.5) node[midway, above, label, align=center]{Protection\\structurelle};
\draw[arrow] (6.65,4) -- (6.65,4.5) node[above, label, align=center]{Espace de\\liberté de l'eau};
\end{tikzpicture}
\legende{Profil transversal d'un littoral protégé (type estuaire du Wouri, Douala). La résilience repose sur la combinaison : (1) Mangrove préservée = 1ère ligne de défense naturelle, (2) Digue/épis = 2ème ligne pour les surcotes extrêmes, (3) Zone de retrait respectée = espace de crue sans enjeu humain, (4) Habitat adapté (pilotis, toitures) = réduction de la vulnérabilité résiduelle. La destruction de la mangrove (bois de chauffe, extension portuaire, urbanisation sauvage) supprime le premier amortisseur et expose la digue et la ville à l'érosion directe.}
\end{popfigure}

\begin{exemplebox}[La mangrove de l'estuaire du Wouri (Douala) : un capital naturel menacé]
L'estuaire du Wouri abrite l'une des plus vastes mangroves d'Afrique centrale (~18 000 ha). Elle protège la ville de Douala (capitale économique, 3-4 millions d'habitants) des surcotes de tempête et de l'érosion côtière.
\textbf{Services écosystémiques de protection :}
\begin{itemize}
    \item \textbf{Atténuation des vagues :} Une bande de 500 m de \emph{Rhizophora} réduit la hauteur des vagues de 50 à 90\% (coefficient $K_t \approx 0,1-0,5$ selon densité).
    \item \textbf{Stabilisation des berges :} Racines échasses (pneumatophores) et contreforts fixent les vases, limitant le recul du trait de côte (actuellement 1-3 m/an par endroits).
    \item \textbf{Tampon inondation :} Zone d'expansion de crue naturelle lors des crues du Wouri + marée haute + surcote.
\end{itemize}
\textbf{Menaces :} Coupe pour bois de chauffe/fumage poisson, extension zone portuaire (Kribi/Douala), urbanisation anarchique (remblais), pollution hydrocarbures/métaux lourds.
\textbf{Conséquence directe :} La régression de la mangrove (perte ~30\% en 30 ans) augmente la fréquence des submersions dans les quartiers bas (Youpwé, Bépanda) et le coût de maintenance des digues artificielles. \textbf{Restaurer la mangrove est une mesure de protection structurelle « verte » à haut rendement économique (coût/bénéfice $\gg$ 1).}
\end{exemplebox}

\courssub{Mesures non structurelles : la gouvernance du risque}

Les murs ne suffisent pas. La loi, l'éducation et la planification sont souvent plus efficaces et moins coûteuses.

\begin{propriete}[Mesures non structurelles (ou « douces »)]
\begin{itemize}
    \item \textbf{Cadre juridique et réglementaire :} Loi-cadre sur l'environnement (1996), Loi sur l'urbanisme (2004), Code de l'urbanisme et de l'habitat. \textbf{Outils :} Plan Local d'Urbanisme (PLU), Plan de Prévention des Risques Naturels (PPRN) -- \emph{peu appliqués au Cameroun à ce jour}. Interdiction formelle de lotir/construire en zone inondable (lit majeur, zones de ruissellement concentré).
    \item \textbf{Gestion intégrée des bassins versants (GIRE) :} Reboisement des têtes de bassin (régulation débit de pointe), agriculture de conservation (couvert végétal, labour en courbes de niveau) pour réduire l'érosion et le ruissellement rapide. Lutte contre l'orpaillage illégal (destruction berges, mercure).
    \item \textbf{Éducation aux risques et culture de la prévention :} Intégration dans les programmes scolaires (SVT, Géographie, Éducation à la Citoyenneté), exercices d'évacuation annuels dans les écoles/hôpitaux, formation des leaders d'opinion, campagnes médias en langues locales (« \emph{Nyanga na mboka} » - la pluie et la ville).
    \item \textbf{Assurance et solidarité :} Développement de l'assurance indicielle (paramétrique) pour l'agriculture (sécheresse/inondation), fonds de catastrophe (Fonds National de Solidarité), micro-crédit post-catastrophe pour relance économique.
    \item \textbf{Planification d'urgence :} Plan Communal de Sauvegarde (PCS), Plan d'Organisation des Secours (PLANORSEC), annuaires des personnes vulnérables (PMR, personnes âgées isolées), conventions avec le privé (logistique, engins, carburant).
\end{itemize}
\end{propriete}

\begin{attention}[Erreur majeure : « L'alerte suffit »]
Beaucoup d'élèves (et de décideurs) pensent qu'émettre un bulletin d'alerte rouge suffit à protéger. \textbf{Faux.} Une alerte sans \textbf{plan d'évacuation connu}, sans \textbf{site d'hébergement identifié et équipé}, sans \textbf{moyens de transport} pour les personnes à mobilité réduite, sans \textbf{confiance} dans l'institution qui alerte, ne sauvera personne. L'éducation aux risques (savoir \emph{quoi faire}, \emph{où aller}, \emph{quand partir}) est une mesure de protection à part entière, souvent plus rentable qu'une digue. Au Probatoire, citez toujours le triptyque : \textbf{Informer $\neq$ Alerter $\neq$ Protéger}.
\end{attention}

\courssub{Comportements à adopter : la responsabilité individuelle et collective}

\emph{Situation :} Le cyclone/tempête arrive, l'alerte est rouge, l'eau monte, le vent arrache les toits. Que faites-vous \emph{maintenant} ?

\begin{aretenir}[Conduites à tenir PENDANT l'événement (Cyclone / Tempête violente / Crue éclair)]
\begin{itemize}
    \item \textbf{Avant l'arrivée (veille / heures pré-cycloniques) :} Couper l'électricité (disjoncteur général) et le gaz/arrivée d'eau. Fermer solidement portes, fenêtres, volets (planches/contreplaqué si pas de volets). Ranger/attacher tout objet volant (tôles, bidons, meubles jardin). Préparer « sac d'urgence » (papiers plastifiés, médicaments, eau, biscuits, lampe torche, radio piles, sifflet). Rejoindre le lieu d'hébergement \textbf{avant} la nuit ou la montée des eaux.
    \item \textbf{Pendant le paroxysme :} \textbf{RESTEZ À L'INTÉRIEUR} (pièce la plus centrale, sans fenêtre, au rez-de-chaussée surélevé ou étage si inondation). N'ouvrez pas pour « voir ». Écoutez la radio (consignes officielles). Ne téléphonez pas (saturez le réseau secours) sauf urgence vitale (112 / 118).
    \item \textbf{FACE À L'EAU (CRUE / SUBMERSION) :} \textbf{JAMAIS} traverser une route inondée (à pied : 30 cm courant fort = chute ; en voiture : 40-50 cm = flottaison/embarquement). 50 cm d'eau = 1 tonne de poussée sur une portière. Ne pas entrer dans les sous-sols/parkings (pièges mortels : montée fulgurante, pas d'issue, électricité). Montez au point haut le plus proche (toit terrasse accessible, étage).
    \item \textbf{APRÈS (décrue / calme apparent) :} Attendez le feu vert officiel (levée vigilance). Attention : \emph{œil du cyclone} = accalmie trompeuse (vent redouble en sens inverse). Eau du robinet = non potable (ébullition/chlore). Électricité = danger électrocution (ne rallumez qu'après contrôle pro). Signalez dangers (fils nus, murs fissurés, cadavres animaux). Documentez dégâts (photos) pour assurance/aide.
\end{itemize}
\end{aretenir}

\courssub{Application au contexte camerounais : réalités et acteurs}

Le Cameroun, par sa position géographique (golfe de Guinée, zone de convergence intertropicale, reliefs volcaniques), est exposé à une multiplicité de risques hydro-météorologiques.

\begin{center}
\begin{tabularx}{\linewidth}{|l|X|X|}\hline
\rowcolor{popblueL}
\textbf{Risque majeur} & \textbf{Zones concernées (exemples)} & \textbf{Facteurs aggravants locaux} \\ \hline
\textbf{Inondations urbaines} & Douala (Wouri, Dibamba), Yaoundé (Mfoundi, Mefou), Garoua (Bénoué), Maroua (Kaliao) & Imperméabilisation sauvage, absence/défaut réseaux drainage, occupation lit majeur, déchets bouchant caniveaux, élévation niveau mer. \\ \hline
\textbf{Glissements de terrain} & Ouest (Bafoussam, Dschang, Santchou - versants Mont Bamboutos, Koupe), Nord-Ouest (Bamenda), Sud-Ouest (Buéa - flancs Mont Cameroun) & Pentes fortes, sols latéritiques profonds lessivés, déforestation (cultures, bois de chauffe), urbanisation non réglementée sur versants, séismes volcaniques (Mont Cameroun). \\ \hline
\textbf{Vents violents / Tornades} & Adamoua, Nord, Extrême-Nord (ligne de grain, orages multicellulaires), Littoral (rafales descendantes orages) & Habitat précaire (tôles mal fixées, cases banco), arbres non élagués près lignes électriques/habitations. \\ \hline
\textbf{Sécheresse / Canicule} & Extrême-Nord, Nord, Adamoua (Sahel soudanien) & Variabilité climatique, dégradation sols, pression démographique, gestion eau conflictuelle. \\ \hline
\end{tabularx}
\end{center}

\begin{savaistu}[Le Bangladesh : la preuve par l'exemple que la prévention sauve des vies]
En 1970, le cyclone \emph{Bhola} tue \textbf{300 000 à 500 000 personnes} (vent 185 km/h, surcote 6-10 m, pas d'alerte, pas d'abris, population dense sur îles basses).
En 2007, le cyclone \emph{Sidr} (intensité comparable, cat. 4-5) tue \textbf{$\approx$ 3 500 personnes}.
En 2020, le cyclone \emph{Amphan} (cat. 5 en mer, affaibli à l'abord) tue \textbf{$\approx$ 100 personnes} au Bangladesh.
\textbf{Pourquoi ?} Investissement massif (Banque Mondiale, ONG, État) sur 40 ans dans : (1) \textbf{Réseau de 14 000 abris cycloniques} (écoles/centres communautaires sur pilotis, accessibles à pied < 30 min), (2) \textbf{Système d'alerte précoce communautaire} (volontaires \emph{CPP - Cyclone Preparedness Program} : drapeaux, sirènes, porte-à-porte, mégaphones, SMS), (3) \textbf{Replantation mangrove} (ceinture verte), (4) \textbf{Éducation continue} (simulacres réguliers, programmes scolaires). \textbf{Leçon :} La mortalité n'est pas une fatalité météorologique, c'est un choix de gouvernance et d'investissement dans la résilience.
\end{savaistu}

\begin{definition}[Acteurs clés de la gestion des risques au Cameroun]
\begin{itemize}
    \item \textbf{Direction de la Météorologie Nationale (DMN / MINTRANSPORT) :} Observations, prévisions, vigilance, climatologie. Centre de prévision de Douala (aéroport) et Yaoundé.
    \item \textbf{Direction de la Protection Civile (MINAT / DPC) :} Coordination secours, plans d'urgence (PLANORSEC), formation sapeurs-pompiers, gestion abris, logistique d'urgence.
    \item \textbf{Collectivités Territoriales Décentralisées (Communes, Régions) :} \textbf{Responsables de premier plan} (Loi 2004/017). Élaborent Plan Communal de Sauvegarde (PCS), gèrent voirie/drainage, délivrent permis de construire (contrôle zones à risque), organisent secours de proximité.
    \item \textbf{Ministère de l'Habitat et du Développement Urbain (MINHDU) :} Urbanisme, PLU, résorption habitat insalubre, normes construction parasismique/cyclonique.
    \item \textbf{Société Civile / ONG / Croix-Rouge Camerounaise / OMS / PNUD / Banque Mondiale :} Appui technique, financement projets résilience (ex. Projet de Gestion des Risques de Catastrophes et de Changement Climatique - PGRC), aide humanitaire post-événement.
\end{itemize}
\end{definition}

\courssub{Exercer le savoir-faire : proposer des mesures adaptées à une situation donnée}

\emph{Compétence exigible au Probatoire :} À partir d'une description de situation (carte de risque, photo aérienne, récit de catastrophe), proposer un ensemble cohérent de mesures de protection structurées (Avant / Pendant / Après).

\begin{methode}[Répondre à une consigne « Proposer des mesures de protection »]
\textbf{Étape 1 : Analyser la situation (Diagnostic).} Identifier : l'aléa (type, intensité, fréquence), les enjeux exposés (population, bâtiments, infrastructures), la vulnérabilité (type habitat, niveau pauvreté, connaissance du risque), les mesures existantes (lacunes).
\textbf{Étape 2 : Structurer la réponse selon les 3 temps de la gestion du risque.}
\begin{itemize}
    \item \textbf{AVANT (Prévention / Aménagement / Préparation) :} Mesures structurelles (digue, mangrove, drainage, normes construction), non structurelles (PPRN, assurance, éducation, exercices, plan alerte).
    \item \textbf{PENDANT (Gestion de crise / Alerte / Secours) :} Chaîne d'alerte (moyens diffusion), évacuation (itinéraires, sites, transport PMR), secours (moyens, coordination), continuité services vitaux (eau, santé, énergie).
    \item \textbf{APRÈS (Relèvement / Reconstruction résiliente / RETEX) :} Évaluation dégâts, relogement temporaire, reconstruction \emph{« mieux qu'avant »} (Build Back Better), soutien psychologique, enquête pour améliorer prévention.
\end{itemize}
\textbf{Étape 3 : Justifier chaque mesure} par son lien direct avec le diagnostic (ex. : « Comme le quartier est en zone de ruissellement concentré (pente 15\%, sols imperméables), je propose des noues végétalisées en amont pour infiltrer l'eau à la source »).
\textbf{Étape 4 : Intégrer la dimension locale} (acteurs locaux, matériaux locaux, culture du risque, genre, handicap).
\end{methode}

\begin{exoresolu}[Proposition de mesures pour un quartier informel en zone inondable (Type Douala - Makèpè)]
\textbf{Énoncé :} Le quartier « Marigot » (fictif, inspiré de Makèpè) est un quartier informel de 15 000 hab. en bordure de mangrove dégradée, zone marécageuse remblayée, densément bâti en matériaux précaires (bois, tôles, banco). Réseau d'assainissement inexistant. Inondations annuelles (hauteur 0,5 à 1,5 m, durée 2-5 jours). Épidémies choléra récurrentes post-crue. Taux de pauvreté élevé. Proposer un plan de protection intégré sur 5 ans.

\tcblower

\textbf{Solution détaillée :}
\textbf{1. Diagnostic partagé (Participatif) :} Ateliers avec riverains, chefs de quartiers, ONG, Mairie, DMN. Cartographie participative des points bas, issues de secours, points d'eau, latrines. Identification personnes vulnérables.
\textbf{2. AVANT (Prévention structurelle \& non structurelle) :}
\begin{itemize}
    \item \textbf{Urgence (An 1) :} Curage manuel caniveaux existants + création fossés de drainage primaires (main-d'œuvre intensive HIMO - jeunes locaux payés). Distribution kits hygiène (chlore, savon, moustiquaires).
    \item \textbf{Court terme (An 1-2) :} \textbf{Restauration mangrove} (10 ha) en zone tampon riveraine (pépinières villageoises \emph{Rhizophora}, plantation participative). \textbf{Surélévation plots} : appui technique/financier (micro-crédit) pour rehausser sols maisons (remblai latéritique compacté + géotextile) + toitures tôles fixées (crochets/torons). \textbf{Points d'eau surélevés} (forages pompes manuelles sur plots 1,5 m). \textbf{Latrines étanches surélevées} (type ECOSAN ou fosses septiques hermétiques).
    \item \textbf{Moyen terme (An 2-4) :} \textbf{Drainage structurant} : collecteur principal bétonné + bassins de rétention en zone basse (parc de détente/infiltration). \textbf{Plan Local d'Urbanisme (PLU) révisé} : zonage « zone à risque maîtrisé » (construction autorisée seulement sur pilotis + normes), « zone de retrait » (non constructible, parc mangrove). \textbf{Assurance indicielle inondation} (paiement automatique si pluviométrie satellite > seuil).
    \item \textbf{Éducation continue :} Club « Jeunes Sentinelles Climat » dans écoles/quartiers (surveillance niveau eau, nettoyage caniveaux, relais alerte radio/WhatsApp). Simulacre évacuation annuel.
\end{itemize}
\textbf{3. PENDANT (Gestion de crise) :}
\begin{itemize}
    \item Système d'alerte local : 3 drapeaux (Vert/Jaune/Orange/Rouge) sur mât central + SMS groupé + crieurs publics. Seuil déclenchement : niveau eau marigot > cote critique (règle à lecture directe peinte sur pilier).
    \item Sites d'hébergement identifiés/pré-équipés : École primaire (surélevée), Église, Centre de santé (étage). Stocks prépositionnés (riz, eau, kits cuisine, médicaments, chlore) gérés par comité de vigilance.
    \item Brigades d'évacuation formées (jeunes valides) avec brancards, barques légères pour PMR/enfants.
\end{itemize}
\textbf{4. APRÈS (Relèvement résilient) :}
\begin{itemize}
    \item Évaluation participative dégâts (grille standardisée). Aide ciblée (cash-for-work reconstruction).
    \item \textbf{Build Back Better} : Reconstruction uniquement selon normes résilientes (pilotis, matériaux durables). Appui foncier (titres d'occupation sécurisés) pour inciter investissement.
    \item RETEX communautaire (réunion bilan) $\rightarrow$ mise à jour PCS et PLU.
\end{itemize}
\textbf{Justification globale :} Approche hybride (grise/verte), participative (appropriation), phasée (réaliste financièrement), intégrant santé (choléra) et genre (femmes gestionnaires eau/latrines). Coût estimé << coût dégâts récurrents + vies sauvées.
\end{exoresolu}

\begin{exercice}[Entraînement au Probatoire - Situation : Glissement de terrain à Bafoussam-Ouest]
Un quartier d'habitat spontané s'est développé sur un versant raide ($35^\circ$) en matériaux volcaniques altérés (latérite épaisse), au pied d'une carrière artisanale active. En saison des pluies 2023, un glissement rotatif a détruit 12 maisons, fait 4 morts. La mairie a décrété « zone non constructible » mais n'a pas relogé les familles qui sont revenues.
\textbf{Consigne :} En vous appuyant sur la méthode « Avant / Pendant / Après », proposez 3 mesures structurelles et 3 mesures non structurelles prioritaires pour réduire le risque résiduel. Justifiez chaque mesure par un facteur de vulnérabilité ou d'aléa identifié dans l'énoncé.
\end{exercice}

\begin{corrige}[Exercice : Glissement de terrain à Bafoussam-Ouest]
\textbf{Analyse :} Aléa : Glissement rotatif réactivable (pente forte, géologie sensible, sous-creusement carrière, pluviométrie forte). Vulnérabilité : Habitat précaire sur zone instable, absence drainage, population revenue par nécessité (pas d'alternative foncière), ignorance risque / fatalisme.
\textbf{Mesures Structurelles (Avant) :}
1. \textbf{Arrêt immédiat carrière + reprofilage/stabilisation tête de versant :} Supprime le facteur déclencheur anthropique (sous-creusement) et réduit la poussée amont. \textit{Justification : cause directe identifiée.}
2. \textbf{Drainage profond (forages drainants + drains horizontaux) + drainage de surface (rigoles bétonnées en courbes de niveau) :} Abaisse la nappe phréatique (facteur déclencheur principal pluies) et évacue ruissellement. \textit{Justification : mécanisme physique du glissement (pression interstitielle).}
3. \textbf{Ouvrages de retenue au pied (berme en enrochement / mur de soutènement gabions) + revégétalisation versants (Vetiver, arbres pionniers) :} Butée mécanique + renforcement racinaire sols. \textit{Justification : protection directe enjeux en pied de pente.}
\textbf{Mesures Non Structurelles :}
1. \textbf{Relogement effectif + sécurisation foncière site d'accueil (Avant) :} Seule mesure \textbf{radicale} supprimant l'enjeu humain en zone rouge. Nécessite volonté politique, budget, concertation. \textit{Justification : vulnérabilité extrême (vies perdues), inefficacité mesures structurelles seules sur pente $35^\circ$ remaniée.}
2. \textbf{Système d'alerte précoce local : inclinomètres / piézomètres automatiques + seuil pluie cumulée 72h + sirène/relais quartier (Pendant) :} Permet évacuation préventive heures/jours avant rupture. \textit{Justification : délai entre pluie intense et glissement (heures à jours) exploitable.}
3. \textbf{Éducation au risque + formation comités de vigilance + exercice évacuation annuel (Avant/Pendant) :} Change la perception fataliste, assure réaction rapide. \textit{Justification : population revenue = besoin compréhension risque + savoir-faire réaction.}
\textbf{Mesure Après :} \textbf{RETEX technique + révision PLU communal intégrant carte aléa glissement (LIDAR/drone) + interdiction stricte nouvelle occupation zone instable.}
\end{corrige}

\begin{aretenir}[Synthèse : La résilience se construit au quotidien]
Face aux grands mouvements atmosphériques et océaniques (tempêtes, cyclones, crues, glissements), la protection ne repose pas sur une « solution miracle » mais sur une \textbf{stratégie intégrée, multi-échelle et continue} :
\begin{itemize}
    \item \textbf{La science (Prévision) :} Indispensable mais insuffisante. Incertitude inhérente (chaos).
    \item \textbf{La technique (Structurel) :} Nécessaire pour les enjeux forts, mais coûteuse, rigide, parfois contre-productive (effet digue $\rightarrow$ fausse sécurité $\rightarrow$ urbanisation accrue). \textbf{Privilégier les solutions fondées sur la nature (SbN) : mangroves, zones d'expansion de crue, forêts de protection.}
    \item \textbf{La gouvernance (Non structurel) :} Le levier le plus puissant et le moins cher : urbanisme (PPRN, PLU), éducation, assurance, planification d'urgence. \textbf{La loi non appliquée ne protège personne.}
    \item \textbf{Le citoyen (Comportement) :} Maillon final. Connaître les gestes qui sauvent (ne pas traverser l'eau, couper électricité, écouter la radio, rejoindre le point haut) est une compétence vitale, enseignée à l'école et répétée dans la communauté.
\end{itemize}
Au Cameroun, comme partout, \textbf{la réduction du risque de catastrophe est un choix de développement}, pas seulement une affaire de secours d'urgence. C'est l'essence de l'Éducation à l'Environnement et au Développement Durable.
\end{aretenir}

\newpage

\coursec{Activité d'intégration}

\courssub{Objectifs de l'activité}

À l'issue de cette activité d'intégration, tu seras capable de :

\begin{itemize}
\item \cle{Décrire} un mouvement atmosphérique ou océanique à partir d'un schéma ou d'une carte (cellules de convection, vents, courant marin) ;
\item \cle{Relier} un phénomène météorologique dangereux (pluies diluviennes, vents violents) à son origine dynamique (convergence des vents, évaporation océanique, convection) ;
\item \cle{Proposer} des mesures de protection adaptées, organisées dans le temps (avant, pendant, après le phénomène) ;
\item Travailler en équipe, argumenter oralement en un temps limité et évaluer la production d'un autre groupe à l'aide d'une grille commune.
\end{itemize}

Cette activité mobilise l'ensemble des savoir-faire de la leçon dans une \cle{situation authentique} : la protection d'une grande ville camerounaise face à un épisode pluvieux extrême.

\courssub{Situation}

\textbf{Mission météo : alerter la ville de Douala.}

Tu es stagiaire à la Direction de la Météorologie Nationale, basée à Douala. Nous sommes le \textbf{12 août}, en pleine saison des pluies. Le chef de service te convoque avec trois autres stagiaires : les images satellites et les mesures des stations côtières sont préoccupantes. Voici le dossier qu'il remet à ton équipe.

\textbf{Document 1 — Carte simplifiée du golfe de Guinée (mois d'août).} La carte montre : le continent africain au nord, l'océan Atlantique équatorial au sud ; le \cle{courant de Guinée}, courant marin chaud qui longe la côte du golfe de Guinée d'ouest en est ; la position de la \cle{ZCIT} (Zone de Convergence InterTropicale), décalée vers le nord en août, au niveau du littoral camerounais ; la ville de Douala (environ 4° de latitude Nord), située au fond de la baie de Wouri.

\textbf{Document 2 — Schéma de la circulation atmosphérique à compléter.} Une coupe verticale de l'atmosphère entre l'équateur et le pôle Nord, avec seulement les indications « équateur », « 30° N », « 60° N », « pôle Nord ». Il manque : les flèches des mouvements verticaux et horizontaux de l'air, les noms des cellules, les zones de hautes et de basses pressions, les noms des vents de surface.

\textbf{Document 3 — Données mesurées et prévisions.}

\begin{center}
\begin{tabularx}{\linewidth}{|l|X|}\hline
\rowcolor{popblueL} \textbf{Paramètre} & \textbf{Valeur mesurée ou prévue} \\ \hline
Température de surface de la mer (golfe de Guinée) & 28 °C \\ \hline
Vents au sud de la ZCIT & mousson de sud-ouest (alizés de sud-est déviés, humides) \\ \hline
Vents au nord de la ZCIT & alizés de nord-est (harmattan, plus sec) \\ \hline
Humidité relative de l'air au-dessus de la mer & 90 \% \\ \hline
Cumul de pluie prévu sur Douala en 72 h & 250 à 300 mm \\ \hline
Rafales de vent prévues & 80 à 100 km/h \\ \hline
\end{tabularx}
\end{center}

\textbf{Contexte.} Douala, capitale économique du Cameroun, compte plus de 3 millions d'habitants. De nombreux quartiers (Bépanda, Deïdo, Ndogpassi, Bonabéri) sont construits dans des zones basses et marécageuses, avec un réseau de drainage insuffisant. En août 2020, des pluies diluviennes avaient déjà provoqué des inondations meurtrières dans la ville.

\textbf{Mission.} Le chef de service attend de ton équipe : un schéma de circulation complété, une analyse du rôle du courant de Guinée, un \cle{bulletin d'alerte} destiné à la radio et à la télévision, et un \cle{plan de protection} de la ville. La restitution orale durera 3 minutes par groupe, évaluée avec la grille commune affichée au tableau.

\courssub{Consignes}

Par groupes de 4, répondez aux tâches suivantes, dans l'ordre :

\begin{enumerate}
\item \textbf{Compléter le schéma de circulation} (Document 2) : dessinez les flèches des mouvements de l'air dans la cellule de Hadley (entre l'équateur et 30° N), indiquez les zones de basses pressions (équateur) et de hautes pressions (30° N), puis nommez les vents de surface (alizés). Légendez chaque élément.
\item \textbf{Décrire le courant de Guinée} à partir du Document 1 : sens d'écoulement, caractère (chaud ou froid), et expliquez son rôle dans l'apport de \cle{chaleur} et d'\cle{humidité} à l'atmosphère au-dessus du golfe de Guinée.
\item \textbf{Expliquer l'origine dynamique du phénomène annoncé} : pourquoi les vents convergent-ils au niveau de Douala ? Pourquoi l'air chaud et humide s'élève-t-il ? Quel lien avec les pluies et les vents violents prévus ?
\item \textbf{Rédiger le bulletin d'alerte} (environ 10 lignes) : il doit être compréhensible par la population, annoncer le phénomène, expliquer brièvement son origine et indiquer les comportements à adopter.
\item \textbf{Proposer un plan de protection en trois volets} : mesures \emph{avant} le phénomène (prévention), \emph{pendant} (conduite à tenir), \emph{après} (secours et remise en état).
\item \textbf{Préparer la restitution orale} de 3 minutes : répartissez la parole entre les quatre membres du groupe.
\end{enumerate}

\courssub{Ressources à mobiliser}

\begin{aretenir}[Ce que tu dois réactiver pour réussir la mission]
\begin{itemize}
\item \textbf{Moteur thermique :} l'air chaud, moins dense, s'\cle{élève} (basses pressions) ; l'air froid, plus dense, \cle{descend} (hautes pressions). Le vent est un déplacement d'air des hautes vers les basses pressions.
\item \textbf{Cellule de Hadley :} air ascendant à l'équateur, déplacement en altitude vers les tropiques, descente vers 30° de latitude, retour au sol par les \cle{alizés}.
\item \textbf{ZCIT :} zone de convergence des alizés des deux hémisphères ; l'air y est chaud, humide et ascendant, d'où des pluies abondantes. En août, elle est décalée vers le nord, au niveau du Cameroun.
\item \textbf{Courant de Guinée :} courant chaud d'ouest en est le long du golfe de Guinée ; il réchauffe l'air et augmente l'\cle{évaporation}, donc l'humidité de l'air.
\item \textbf{Chaîne causale du phénomène :} mer chaude $\rightarrow$ évaporation $\rightarrow$ air humide $+$ convergence des vents $\rightarrow$ ascendance $\rightarrow$ refroidissement, condensation $\rightarrow$ nuages, pluies intenses et vents violents.
\item \textbf{Protection :} prévention (information, aménagement), conduite à tenir pendant, secours et assainissement après.
\end{itemize}
\end{aretenir}

\begin{attention}[Pièges fréquents à éviter dans ta production]
\begin{itemize}
\item Ne confonds pas \textbf{vent} et \textbf{courant marin} : le vent est un déplacement d'\emph{air}, le courant marin un déplacement d'\emph{eau} de mer.
\item Le vent de surface va des \textbf{hautes} vers les \textbf{basses} pressions, jamais l'inverse.
\item Les alizés de l'hémisphère Nord sont des vents de \textbf{nord-est} (ils viennent du nord-est) ; la mousson qui atteint Douala vient du \textbf{sud-ouest}. Un vent est nommé d'après la direction \emph{d'où il vient}.
\item La pluie ne tombe pas « parce qu'il fait chaud » : elle résulte de l'\textbf{ascendance} de l'air humide, de son refroidissement et de la \textbf{condensation} de la vapeur d'eau. Chaque étape de la chaîne causale doit être citée.
\item La ZCIT n'est pas fixe : elle se déplace au cours de l'année en suivant le zénith du Soleil (vers le nord en août, vers le sud en janvier).
\end{itemize}
\end{attention}

\courssub{Démarche et corrigé commenté}

\begin{exoresolu}[Corrigé commenté de la mission]
\textbf{Étape 1 — Schéma de circulation complété.}

Dans la cellule de Hadley : au niveau de l'équateur, l'air fortement réchauffé s'élève (flèche verticale ascendante) : c'est la zone des \textbf{basses pressions} équatoriales. En altitude, l'air se déplace vers le nord, se refroidit et redescend vers 30° N (flèche verticale descendante) : zone des \textbf{hautes pressions} subtropicales. Au sol, l'air revient vers l'équateur : ce sont les \textbf{alizés de nord-est} dans l'hémisphère Nord. On complète par la cellule de Ferrel (30°--60° N), à circulation inverse, et la cellule polaire (60° N--pôle), avec les vents d'ouest des latitudes moyennes et les vents d'est polaires.

\emph{Justification :} l'air chaud est moins dense, il s'élève ; l'air froid est plus dense, il descend. Le vent de surface va toujours des hautes vers les basses pressions.

\textbf{Étape 2 — Description du courant de Guinée.}

Le courant de Guinée est un courant marin \textbf{chaud} qui s'écoule d'\textbf{ouest en est} le long des côtes du golfe de Guinée. Avec une température de surface de 28 °C, il agit comme un « radiateur » sous l'atmosphère : il réchauffe les couches basses de l'air et alimente une forte \textbf{évaporation}. L'air au-dessus du golfe devient chaud et très humide (90 \% d'humidité relative) : c'est le « réservoir » de chaleur et de vapeur d'eau qui alimentera les pluies de la saison.

\textbf{Étape 3 — Origine dynamique du phénomène.}

En août, la ZCIT est positionnée sur le littoral camerounais. De part et d'autre, les vents \textbf{convergent} : la mousson de sud-ouest, chargée d'humidité océanique, et les alizés de nord-est, plus secs. Cette convergence force l'air chaud et humide à \textbf{s'élever}. En montant, l'air se détend et se refroidit ; la vapeur d'eau se \textbf{condense} en gouttelettes : formation d'épais nuages de convection (cumulonimbus), pluies torrentielles (250 à 300 mm en 72 h) et vents violents (rafales de 80 à 100 km/h). La chaîne causale est donc : mer chaude $+$ courant de Guinée $\rightarrow$ évaporation $\rightarrow$ convergence à la ZCIT $\rightarrow$ ascendance $\rightarrow$ condensation $\rightarrow$ pluies diluviennes et vents destructeurs.

\textbf{Étape 4 — Bulletin d'alerte (exemple de rédaction).}

« Alerte météo — ville de Douala, 12 août. La Météorologie Nationale annonce pour les 72 prochaines heures des pluies diluviennes, de 250 à 300 millimètres, accompagnées de vents violents pouvant atteindre 100 kilomètres par heure. Ce phénomène s'explique par la rencontre, au-dessus de notre région, de vents humides venus de l'océan et de vents venus du continent : l'air chaud et chargé d'eau s'élève et forme de violents orages. Risques majeurs : inondations des quartiers bas (Bépanda, Deïdo, Ndogpassi, Bonabéri), chutes d'arbres et de toitures, coupures d'électricité. Consignes : évitez tout déplacement inutile, ne traversez jamais une rue inondée, éloignez-vous des arbres et des lignes électriques, mettez vos documents et biens en hauteur, écoutez la radio. Numéros d'urgence : sapeurs-pompiers 118, SAMU 119. Restez vigilants et suivez les consignes des autorités. »

\textbf{Étape 5 — Plan de protection en trois volets.}

\emph{Avant (prévention) :} diffusion de l'alerte par radio, télévision et réseaux sociaux ; curage des caniveaux et des évacuateurs d'eau ; identification des sites d'accueil (écoles, salles des fêtes) ; stockage d'eau potable, de vivres et de médicaments ; sensibilisation des quartiers à risque.

\emph{Pendant :} interdiction de traverser les zones inondées ; coupure préventive de l'électricité dans les quartiers submergés ; évacuation des habitants des zones les plus basses vers les sites d'accueil ; maintien de l'écoute des bulletins météo.

\emph{Après :} secours et recensement des victimes ; assainissement et désinfection des zones inondées (risque de choléra et de paludisme) ; remise en état des routes et réseaux ; bilan et retour d'expérience pour améliorer le plan communal de protection.

\textbf{Étape 6 — Restitution orale.}

Répartition conseillée : élève 1 présente le schéma, élève 2 le courant de Guinée, élève 3 lit le bulletin, élève 4 expose le plan de protection. Grille d'évaluation commune : exactitude scientifique (4 pts), qualité du schéma (4 pts), clarté du bulletin (4 pts), réalisme du plan de protection (4 pts), expression orale et respect du temps (4 pts).
\end{exoresolu}

\begin{popfigure}
\begin{center}
\begin{tikzpicture}[scale=1.0, >=Stealth]
% Sol
\draw[thick, popdark] (0,0) -- (12,0);
\foreach \x/\lab in {0/{Équateur},4/{30° N},8/{60° N},12/{Pôle Nord}}
  \draw[popdark] (\x,0) -- (\x,-0.15) node[below, font=\small] {\lab};
% Cellule de Hadley
\draw[->, very thick, poporange] (0.4,0.3) -- (0.4,3.2);
\draw[->, very thick, poporange] (0.4,3.4) -- (3.6,3.4);
\draw[->, very thick, poporange] (3.8,3.2) -- (3.8,0.3);
\draw[->, very thick, popblue] (3.6,0.25) -- (0.6,0.25);
\node[poporange, font=\small\bfseries] at (2,2.2) {Cellule de Hadley};
% Cellule de Ferrel
\draw[->, very thick, popgreen] (4.2,0.3) -- (4.2,3.2);
\draw[->, very thick, popgreen] (4.4,3.4) -- (7.6,3.4);
\draw[->, very thick, popgreen] (7.8,3.2) -- (7.8,0.3);
\draw[->, very thick, popblue] (7.6,0.25) -- (4.4,0.25);
\node[popgreen, font=\small\bfseries] at (6,2.2) {Cellule de Ferrel};
% Cellule polaire
\draw[->, very thick, poppurple] (8.4,0.3) -- (8.4,3.2);
\draw[->, very thick, poppurple] (8.6,3.4) -- (11.4,3.4);
\draw[->, very thick, poppurple] (11.6,3.2) -- (11.6,0.3);
\draw[->, very thick, popblue] (11.4,0.25) -- (8.6,0.25);
\node[poppurple, font=\small\bfseries] at (10,2.2) {Cellule polaire};
% Pressions
\node[font=\small\bfseries, poppink] at (0.4,-0.9) {Basses P};
\node[font=\small\bfseries, popblue] at (3.8,-0.9) {Hautes P};
\node[font=\small\bfseries, poppink] at (7.8,-0.9) {Basses P};
\node[font=\small\bfseries, popblue] at (11.6,-0.9) {Hautes P};
% Vents de surface
\node[font=\scriptsize, popblue, align=center] at (2,0.75) {alizés\\ de NE};
\node[font=\scriptsize, popblue, align=center] at (6,0.75) {vents\\ d'ouest};
\node[font=\scriptsize, popblue, align=center] at (10,0.75) {vents\\ d'est};
\end{tikzpicture}
\end{center}
\legende{Schéma de la circulation générale de l'atmosphère dans l'hémisphère Nord (corrigé du Document 2). Flèches orange, vertes et violettes : mouvements verticaux et en altitude des trois cellules de convection ; flèches bleues au sol : vents de surface, dirigés des hautes pressions (Hautes P) vers les basses pressions (Basses P).}
\end{popfigure}

\courssub{Bilan de l'activité}

Cette mission a permis de mettre en évidence trois idées essentielles de la leçon :

\begin{itemize}
\item Les grands mouvements de l'atmosphère et de l'hydrosphère forment un \cle{système couplé} : le courant chaud de Guinée alimente l'atmosphère en chaleur et en humidité, et la circulation atmosphérique (alizés, mousson, ZCIT) transforme cette énergie en phénomènes météorologiques.
\item Un phénomène dangereux n'est jamais un « hasard » : les pluies diluviennes de Douala s'expliquent par une \cle{chaîne causale dynamique} précise (évaporation, convergence, ascendance, condensation), que l'on peut décrire, prévoir et expliquer à la population.
\item Comprendre l'origine d'un phénomène est la première étape de la \cle{protection des populations} : une alerte fondée sur la science permet d'organiser la prévention, la conduite à tenir et les secours, et de sauver des vies.
\end{itemize}

Tu as ainsi mobilisé les trois savoir-faire de la leçon — décrire, relier, proposer — dans une situation qui est exactement celle des météorologues et des services de protection civile du Cameroun. Retiens que la rigueur scientifique du raisonnement est ce qui rend un bulletin d'alerte crédible et un plan de protection efficace.

\newpage

\coursec{Méthodes et exercices résolus}

\courssub{Méthodes et exercices résolus}

\textbf{Introduction.}\par Cette partie te donne les clés pour réussir les exercices types du Probatoire. Pour chaque savoir-faire officiel, une fiche \textbf{Méthode} structure ta démarche, suivie d'un \textbf{Exercice résolu} corrigé pas à pas. Travaille-les activement : cache la solution, essaie, puis compare.

% ============================================================
% SAVOIR-FAIRE 1 : Décrire un mouvement atmosphérique ou océanique à partir d'un schéma ou d'une carte
% ============================================================

\begin{methode}[Décrire la circulation générale de l'atmosphère sur une carte/schéma]
\textbf{Objectif :} Produire une description rigoureuse, structurée et vocabulaire précis d'une figure représentant les cellules de convection et les vents planétaires.
\begin{enumerate}
    \item \textbf{Identifier le cadre :} Préciser s'il s'agit d'une coupe méridienne (latitude vs altitude) ou d'une planisphère (vue du dessus). Noter l'hémisphère concerné (Nord/Sud) ou la vue globale.
    \item \textbf{Repérer les grandes cellules :} Nommer les trois cellules par hémisphère : \cle{Cellule de Hadley} (équateur $\to$ 30°), \cle{Cellule de Ferrel} (30° $\to$ 60°), \cle{Cellule polaire} (60° $\to$ pôle). Indiquer le sens de rotation (horaire/anti-horaire) et les mouvements verticaux (ascendance/descendance).
    \item \textbf{Localiser les zones de convergence/divergence :} \cle{ZCIT} (Zone de Convergence Intertropicale) à l'équateur (ascendance, basses pressions, pluies) ; \cle{Subtropicaux} vers 30° (descendance, hautes pressions, déserts) ; \cle{Front polaire} vers 60° (convergence air polaire/air tempéré, ascendance, dépressions).
    \item \textbf{Nommer et orienter les vents de surface (vents planétaires) :} \cle{Alizés} (Est $\to$ Ouest, équateur $\to$ 30°, déviés par Coriolis : NE hémisphère Nord, SE hémisphère Sud) ; \cle{Vents d'ouest} (Ouest $\to$ Est, 30° $\to$ 60°) ; \cle{Vents polaires d'est} (Est $\to$ Ouest, 60° $\to$ pôle).
    \item \textbf{Mentionner la force de Coriolis :} Expliquer la déviation vers la \textbf{droite} dans l'hémisphère Nord, vers la \textbf{gauche} dans l'hémisphère Sud (cause : rotation de la Terre).
    \item \textbf{Conclure par le bilan :} Redistribution de l'énergie (chaleur) des tropiques vers les pôles.
\end{enumerate}
\textbf{Reflexes Probatoire :} Toujours citer les latitudes approximatives (0°, 30°, 60°, 90°). Utiliser les termes « ascendance », « descendance », « convergence », « divergence », « alizés », « vents d'ouest ». Ne pas confondre vents d'altitude (retour) et vents de surface.
\end{methode}

\begin{exoresolu}[Analyse d'une coupe méridienne de la circulation atmosphérique générale]
\textbf{Énoncé.} On te donne une coupe méridienne schématique de l'atmosphère terrestre (hémisphère Nord représenté). Le schéma montre :
\begin{itemize}
    \item À l'équateur : fort réchauffement, air montant, nuages convectifs, pluies abondantes.
    \item Vers 30° N : air descendant, ciel clair, zones désertiques (Sahara).
    \item Vers 60° N : rencontre d'air froid polaire et d'air tempéré, front polaire, ascendance, dépressions.
    \item Flèches au sol : vents du NE entre 0° et 30° N ; vents du SO entre 30° et 60° N ; vents du NE au-delà de 60° N.
    \item Flèches en altitude : sens inverse des vents de surface.
\end{itemize}
\textbf{Questions :}
\begin{enumerate}
    \item Nomme les trois cellules de convection représentées et donne leurs limites latitudinales approximatives.
    \item Explique l'origine physique du mouvement ascendant à l'équateur et du mouvement descendant aux subtropicaux.
    \item Nomme les vents de surface dans chaque cellule et précise leur direction. Justifie leur déviation par rapport à un trajet méridien direct.
    \item Quel est le rôle global de cette circulation pour le climat terrestre ?
\end{enumerate}

\tcblower
\textbf{Solution détaillée.}

\textbf{1. Identification des cellules et limites.}
\begin{itemize}
    \item \textbf{Cellule de Hadley} : s'étend de l'équateur (0°) à environ 30° de latitude Nord. C'est la cellule thermiquement directe la plus puissante (moteur principal).
    \item \textbf{Cellule de Ferrel} : s'étend d'environ 30° N à 60° N. C'est une cellule thermiquement indirecte (mécaniquement entraînée par les cellules adjacentes).
    \item \textbf{Cellule polaire} : s'étend de 60° N au pôle (90° N). Cellule thermiquement directe (refroidissement au pôle $\to$ descendance).
\end{itemize}

\textbf{2. Origine des mouvements verticaux.}
\begin{itemize}
    \item \textbf{Équateur (Ascendance / ZCIT) :} Fort ensoleillement zénital permanent $\to$ fort réchauffement de la surface $\to$ réchauffement de la couche d'air au contact $\to$ dilatation, diminution de la densité $\to$ \textbf{instabilité convective} $\to$ ascendance rapide de l'air chaud et humide. Refroidissement adiabatique $\to$ condensation $\to$ nuages cumulonimbus $\to$ pluies diluviennes.
    \item \textbf{Subtropicaux ~30° N (Descendance) :} L'air monté à l'équateur se déplace en altitude vers les pôles. Il se refroidit par rayonnement infrarouge vers l'espace, perd son humidité (précipitations équatoriales), devient dense. Par conservation du moment cinétique, il s'accumule vers 30° N $\to$ \textbf{subsidence} (descendance forcée). Échauffement adiabatique par compression $\to$ évaporation des nuages $\to$ ciel clair, hautes pressions au sol (anticyclones subtropicaux : Açores, Sahara).
\end{itemize}

\textbf{3. Vents de surface et force de Coriolis.}
\begin{center}
\begin{tabularx}{\linewidth}{|l|X|X|X|}\hline
\textbf{Cellule} & \textbf{Latitudes} & \textbf{Vent de surface (Nom \& Direction)} & \textbf{Justification déviation (Coriolis)} \\ \hline
Hadley & 0° -- 30° N & \textbf{Alizés de Nord-Est} (NE $\to$ SO) & Air part vers l'équateur (Sud). Déviation vers la \textbf{droite} (hémisphère Nord) $\to$ vient du NE. \\ \hline
Ferrel & 30° N -- 60° N & \textbf{Vents d'Ouest} (SO $\to$ NE) & Air part vers le pôle (Nord). Déviation vers la \textbf{droite} $\to$ vient de l'Ouest (SO). \\ \hline
Polaire & 60° N -- 90° N & \textbf{Alizés polaires / Vents d'Est polaires} (NE $\to$ SO) & Air part du pôle vers l'équateur (Sud). Déviation vers la \textbf{droite} $\to$ vient du NE. \\ \hline
\end{tabularx}
\end{center}
\emph{Rappel :} La force de Coriolis ($ \vec{F}_c = 2m(\vec{\Omega} \times \vec{v}) $) dévie tout mouvement horizontal vers la droite dans l'hémisphère Nord, vers la gauche dans l'hémisphère Sud. Elle est nulle à l'équateur, maximale aux pôles.

\textbf{4. Rôle global.}
Cette circulation assure le \textbf{transport méridien de chaleur} (énergie sensible + chaleur latente de l'eau évaporée aux tropiques) vers les hautes latitudes. Sans elle, l'équateur serait insupportablement chaud et les pôles infiniment plus froids. Elle maintient l'habitabilité planétaire.

\textbf{Conclusion.} La description complète associe géométrie des cellules (latitudes, sens rotation), physique des transferts verticaux (convection/subsidence), dynamique des vents de surface (noms, directions, Coriolis) et fonction climatique (redistribution énergie).
\end{exoresolu}

% ============================================================
% SAVOIR-FAIRE 1 (bis) : Décrire la circulation océanique
% ============================================================

\begin{methode}[Décrire les courants marins sur une carte planisphère]
\textbf{Objectif :} Analyser une carte des courants de surface (gyres, courants de bordure, upwelling) et relier leur origine aux vents planétaires et à la topographie.
\begin{enumerate}
    \item \textbf{Repérer les gyres subtropicaux :} 5 grands gyres (N. Pacifique, S. Pacifique, N. Atlantique, S. Atlantique, Indien). Sens de rotation : \textbf{horaire} hémisphère Nord, \textbf{anti-horaire} hémisphère Sud (conséquence de Coriolis sur le transport d'Ekman).
    \item \textbf{Distinguer courants de bordure ouest et est :}
    \begin{itemize}
        \item \textbf{Bordure Ouest (Courants chauds) :} Étroits, profonds, rapides (ex: \cle{Gulf Stream}, \cle{Kuroshio}, \cle{Brasil}, \cle{Est-Australie}, \cle{Agulhas}). Transportent chaleur vers les pôles.
        \item \textbf{Bordure Est (Courants froids) :} Larges, lents, peu profonds (ex: \cle{Canaries}, \cle{Benguela}, \cle{Californie}, \cle{Humboldt}, \cle{Australie-Ouest}). Ramènent eau froide vers l'équateur.
    \end{itemize}
    \item \textbf{Identifier l'upwelling (surgences) :} Côtes Est des océans (vents parallèles au littoral + Coriolis $\to$ transport d'Ekman au large $\to$ remontée eau profonde). Riches en nutriments $\to$ zones de pêche majeures (ex: Pérou/Chili, Namibie, Mauritanie/Sénégal).
    \item \textbf{Citer la circulation thermohaline (profonde) :} Moteur : différences de densité (Température $\downarrow$, Salinité $\uparrow$ $\to$ $\rho \uparrow$). Zones de formation d'eau profonde : Mer du Labrador, Mer du Groenland, Mer de Weddell (Antarctique). « Ceinture transporteuse » globale (temps de parcours $\sim$ 1000 ans).
    \item \textbf{Relier aux vents :} Courants de surface = réponse dynamique aux \textbf{vents planétaires} (alizés $\to$ courant équatorial ; vents d'ouest $\to$ dérive de surface). Frottement air-mer + Coriolis.
\end{enumerate}
\textbf{Reflexes :} Sur carte, flèches rouges = chauds, bleues = froids. Ne pas confondre courant de surface (ventilé) et courant profond (thermohaline). Préciser « courant de bordure ouest » pour Gulf Stream/Kuroshio.
\end{methode}

\begin{exoresolu}[Interprétation de la carte des courants de l'Atlantique Nord et conséquences climatiques]
\textbf{Énoncé.} On observe une carte des courants de surface de l'Atlantique Nord.
\begin{itemize}
    \item Le \textbf{Gulf Stream} quitte le golfe du Mexique, longe la côte Est des USA (Floride $\to$ Cap Hatteras), puis s'éloigne vers le NE en devenant \textbf{Dérive Nord-Atlantique}.
    \item Le \textbf{Courant des Canaries} descend le long de l'Afrique du Nord (Maroc, Mauritanie, Sénégal).
    \item Le \textbf{Courant équatorial Nord} traverse d'Est en Ouest sous les alizés.
    \item Des flèches indiquent un upwelling majeur au large de la Mauritanie/Sénégal.
\end{itemize}
\textbf{Questions :}
\begin{enumerate}
    \item Nomme le gyre représenté et son sens de rotation. Justifie ce sens par la dynamique des vents.
    \item Compare le Gulf Stream et le Courant des Canaries (nature thermique, largeur, vitesse, rôle climatique).
    \item Explique le mécanisme de l'upwelling au large de la Mauritanie et son importance écologique et économique pour le Sénégal et la Mauritanie.
    \item Quelle serait la conséquence climatique pour l'Europe de l'Ouest si le Gulf Stream s'affaiblissait significativement (scénario fonte Groenland) ?
\end{enumerate}

\tcblower
\textbf{Solution détaillée.}

\textbf{1. Gyre et sens de rotation.}
C'est le \textbf{Gyre subtropical de l'Atlantique Nord}. Son sens de rotation est \textbf{horaire} (dans le sens des aiguilles d'une montre).
\textbf{Justification :} Les \textbf{alizés} (NE) soufflent sur la partie sud du gyre (0°--30° N) $\to$ poussent l'eau vers l'Ouest (Courant équatorial Nord). Les \textbf{vents d'ouest} soufflent sur la partie nord (30°--60° N) $\to$ poussent l'eau vers l'Est (Dérive Nord-Atlantique). La force de Coriolis (hémisphère Nord) dévie le transport d'Ekman net vers la \textbf{droite} du vent, accumulant l'eau au centre du gyre (bosse de 1--2 m) $\to$ crée une pente de pression $\to$ courant géostrophique tournant horaire autour de cette bosse.

\textbf{2. Comparaison Gulf Stream vs Courant des Canaries.}
\begin{center}
\begin{tabularx}{\linewidth}{|l|X|X|}\hline
\textbf{Critère} & \textbf{Gulf Stream (Bordure Ouest)} & \textbf{Courant des Canaries (Bordure Est)} \\ \hline
\textbf{Nature thermique} & \textbf{Chaud} (eau tropicale, $\sim$ 25--28°C) & \textbf{Froid} (eau subtropicale/subpolaire, $\sim$ 15--20°C) \\ \hline
\textbf{Largeur / Profondeur} & \textbf{Étroit} (50--100 km), \textbf{Profond} (> 1000 m) & \textbf{Large} (500--1000 km), \textbf{Superficiel} (< 500 m) \\ \hline
\textbf{Vitesse} & \textbf{Rapide} (1--2 m/s, $\sim$ 3--7 km/h) & \textbf{Lent} (0,1--0,5 m/s) \\ \hline
\textbf{Transport de chaleur} & Massif vers le pôle ($\sim$ 1,3 PW = $10^{15}$ W) & Faible vers l'équateur \\ \hline
\textbf{Rôle climatique} & Réchauffe l'Europe de l'Ouest (hivers doux) & Refroidit les côtes africaines, favorise aridité (désert Sahara/Sahel côtier) \\ \hline
\end{tabularx}
\end{center}
\textit{Note :} L'intensification des courants de bordure ouest (Stommel, Munk) est due à la variation de Coriolis avec la latitude ($\beta$-effet).

\textbf{3. Upwelling mauritanien/sénégalais.}
\begin{itemize}
    \item \textbf{Mécanisme :} Les \textbf{alizés} (vents du NE) soufflent \textbf{parallèles à la côte} (orientée N-S). Force de Coriolis (hémisphère Nord) dévie le transport de surface (couche d'Ekman) vers la \textbf{droite} du vent, c'est-à-dire \textbf{au large} (vers l'Ouest).
    \item \textbf{Conséquence :} Éloignement de l'eau de surface côtière $\to$ remontée compensatoire (\textbf{upwelling}) d'eau profonde, froide, riche en sels nutritifs (nitrates, phosphates, silices) provenant de la décomposition de la matière organique en profondeur.
    \item \textbf{Importance :} Explosion de productivité primaire (phytoplancton) $\to$ base chaîne alimentaire $\to$ \textbf{zones de pêche les plus riches au monde} (sardinelles, maquereaux, céphalopodes). Ressource vitale pour la sécurité alimentaire et l'économie (exportations) du Sénégal et de la Mauritanie.
\end{itemize}

\textbf{4. Affaiblissement du Gulf Stream (Scénario AMOC).}
\begin{itemize}
    \item \textbf{Cause :} Fonte massive du Groenland $\to$ apport d'eau douce (faible salinité) en Mer du Labrador/Île de Baffin $\to$ diminution de la densité de l'eau de surface $\to$ inhibition de la convection profonde (formation d'eau profonde NADW) $\to$ ralentissement de la \cle{Circulation Méridienne de Retournement de l'Atlantique (AMOC)}.
    \item \textbf{Conséquence climatique Europe de l'Ouest :} Diminution drastique du transport de chaleur océanique vers le NE $\to$ \textbf{refroidissement relatif} (ou réchauffement moins fort que la moyenne globale) sur l'Europe de l'Ouest, hivers plus rudes, modification du régime des pluies (dépressions déviées). Paradoxe : réchauffement global $\to$ refroidissement local temporaire régional.
    \item \textbf{Autres impacts :} Élévation niveau mer accélérée côte Est USA (effet géostrophique), décalage ZCIT $\to$ sécheresses Sahel, perturbation mousson.
\end{itemize}

\textbf{Conclusion.} La description d'une carte océanique exige de nommer les structures (gyres, courants de bordure, upwelling), d'en expliquer la dynamique (vents, Coriolis, topographie) et d'en déduire les impacts climatiques et écologiques régionaux.
\end{exoresolu}

% ============================================================
% SAVOIR-FAIRE 2 : Relier un phénomène météorologique dangereux à son origine dynamique
% ============================================================

\begin{methode}[Relier un cyclone tropical à son origine dynamique et aux conditions de formation]
\textbf{Objectif :} Expliquer la genèse, la structure, l'énergie et la trajectoire d'un cyclone (ouragan/typhon) en utilisant le vocabulaire précis.
\begin{enumerate}
    \item \textbf{Rappeler les 6 conditions nécessaires (Check-list de Gray) :}
    \begin{enumerate}
        \item \textbf{Température de surface de la mer (TSM) $\ge$ 26,5°C} sur une couche d'au moins 50 m d'épaisseur (réservoir d'énergie).
        \item \textbf{Instabilité conditionnelle de l'atmosphère} (gradient de température vertical fort) $\to$ convection profonde possible.
        \item \textbf{Humidité relative élevée} dans la moyenne troposphère (500 hPa) $\to$ limite l'entrainement d'air sec qui tue la convection.
        \item \textbf{Distance à l'équateur $\ge$ 5° de latitude} (force de Coriolis $f = 2\Omega \sin\varphi$ suffisante pour initier la rotation cyclonique).
        \item \textbf{Pré-existence d'une perturbation de basse pression} (onde tropicale, creux barométrique, ZCIT active) $\to$ « grain » initial.
        \item \textbf{Faible cisaillement vertical du vent} (différence vitesse/direction vent 200 hPa -- 850 hPa $< 10$ m/s) $\to$ évite de déchirer la structure verticale (déplacement du « chaud » par rapport au centre).
    \end{enumerate}
    \item \textbf{Décrire le moteur thermique (Machine de Carnot) :} Source chaude = mer (TSM) ; Puits froid = tropopause ($\sim$ -70°C). Rendement $\eta = 1 - T_{froid}/T_{chaud}$. Énergie = \textbf{Chaleur latente de condensation} libérée dans l'œil et le mur de l'œil. Pression centrale chute $\to$ gradient de pression $\to$ vents violents.
    \item \textbf{Structure :} \cle{Œil} (centre, air descendant, clair, calme, chaud) ; \cle{Mur de l'œil} (anneau convection intense, vents max, pluies max) ; \cle{Bandes spirales} (pluies/orages) ; \cle{Hauteur} $\sim$ 15--18 km (tropopause).
    \item \textbf{Trajectoire :} Guidage par le flux de direction (vent moyen 850--300 hPa). Généralement Ouest $\to$ Nord-Ouest (alizés) puis recourbement vers Nord-Est (vents d'ouest) en latitudes moyennes. Déviation vers la droite (Hémisphère Nord) due à Coriolis ($\beta$-dérive).
    \item \textbf{Dangers :} Vent violent (échelle Saffir-Simpson), \cle{Surcote} (élévation niveau mer $\sim$ vent + basse pression), Pluies diluviennes (inondations), Tornades embarquées.
\end{enumerate}
\textbf{Reflexes :} Ne pas dire « cyclone se forme à l'équateur » (Coriolis nul). Ne pas confondre cyclone tropical (cœur chaud, symétrique, pas de front) et dépression tempérée (cœur froid, fronts, asymétrique). « Ouragan » = Atlantique/Nord-Est Pacifique ; « Typhon » = Nord-Ouest Pacifique ; « Cyclone » = Sud-Ouest Pacifique/Océan Indien.
\end{methode}

\begin{exoresolu}[Genèse et dangerosité d'un cyclone tropical : cas d'un système en Atlantique Nord]
\textbf{Énoncé.} Fin août, une onde tropicale quitte la côte africaine (Sénégal) vers 15° N, 20° W. La TSM est de 29°C. Le cisaillement vertical est faible (5 m/s). L'atmosphère est instable et humide. Le système s'organise en dépression tropicale, puis tempête tropicale (vents 65 km/h), puis ouragan de catégorie 3 (vents 185 km/h) en approchant des Petites Antilles.
\textbf{Questions :}
\begin{enumerate}
    \item Vérifie que les 6 conditions de formation (Gray) sont réunies pour ce système. Identifie la condition qui \textbf{n'est pas} remplie à l'origine (côte africaine) et explique comment elle le devient.
    \item Explique le mécanisme énergétique : d'où vient l'énergie cinétique des vents ? Pourquoi la pression centrale chute-t-elle ?
    \item Décris la structure verticale et horizontale de l'ouragan catégorie 3 au stade mature.
    \item Prédis la trajectoire probable et justifie-la. Quels sont les 3 dangers majeurs pour la population des îles touchées ?
\end{enumerate}

\tcblower
\textbf{Solution détaillée.}

\textbf{1. Check-list des conditions de formation (Gray).}
\begin{enumerate}
    \item \textbf{TSM $\ge$ 26,5°C :} \textbf{OUI} (29°C donné, couche chaude épaisse en fin d'été).
    \item \textbf{Instabilité conditionnelle :} \textbf{OUI} (air maritime tropical chaud/humide en basses couches, air plus sec/froid en altitude).
    \item \textbf{Humidité moyenne troposphère :} \textbf{OUI} (masse d'air maritime tropical, pas de couche d'air saharien sec (SAL) mentionnée comme inhibiteur ici).
    \item \textbf{Distance équateur $\ge$ 5° :} \textbf{OUI} (15° N $\to$ Coriolis $f \approx 3,7 \times 10^{-5}$ s$^{-1}$, suffisant).
    \item \textbf{Perturbation initiale :} \textbf{OUI} (Onde tropicale / \cle{Onde d'Est africaine} -- AEW -- quittant la côte).
    \item \textbf{Faible cisaillement vertical :} \textbf{OUI} (5 m/s $< 10$ m/s).
\end{enumerate}
\textbf{Point subtil :} Au large de l'Afrique (naissance onde), la condition \textbf{n°3 (humidité)} est souvent \textbf{NON remplie} à cause de la \cle{Couche d'Air Saharien (SAL)} : air très sec, chaud, poussiéreux (850--500 hPa) qui inhibe la convection. Le système doit s'éloigner vers l'Ouest (vers 40--50° W) pour sortir de l'influence de la SAL et trouver une atmosphère humide. C'est souvent le facteur limitant en début de saison.

\textbf{2. Moteur thermique et chute de pression.}
\begin{itemize}
    \item \textbf{Source d'énergie :} \textbf{Évaporation} massive à la surface de l'océan (flux de chaleur latente $L_v E$, $L_v \approx 2,5 \times 10^6$ J/kg). L'air chaud et humide converge vers le centre en basses couches (couche limite).
    \item \textbf{Ascendance dans le mur de l'œil :} L'air monte en se refroidissant adiabatiquement $\to$ \textbf{condensation} $\to$ libération chaleur latente $\to$ réchauffement diabatique de la colonne d'air centrale.
    \item \textbf{Conséquence hydrostatique :} Colonne d'air plus chaude $\to$ moins dense $\to$ poids de la colonne diminue $\to$ \textbf{Pression au sol chute} (jusqu'à 900--950 hPa cat 3, $<$ 900 hPa cat 5).
    \item \textbf{Gradient de pression :} Différence pression centre / environnement ($\Delta p \sim 50$--100 hPa sur 50--100 km) $\to$ Force de gradient de pression intense $\to$ vents violents (équilibre cyclostrophique/gradient dans le mur de l'œil : $V \approx \sqrt{\frac{1}{\rho}\frac{\partial p}{\partial r} r}$).
    \item \textbf{Rendement Carnot :} $\eta \approx \frac{T_{mer} - T_{tropopause}}{T_{mer}} \approx \frac{302 - 200}{302} \approx 33\%$. Efficacité remarquable.
\end{itemize}

\textbf{3. Structure mature (Cat 3 : vents 178--208 km/h $\approx$ 50--58 m/s).}
\begin{itemize}
    \item \textbf{Horizontal :} \textbf{Œil} (diam 20--60 km) : subsidence (air descendant $\to$ adiabatique $\to$ réchauffement $\to$ dissipation nuages), vent calme, pression mini. \textbf{Mur de l'œil} : anneau 10--20 km large, convection la plus intense, vents max (niveau 850--700 hPa), pluie extrême ($> 50$ mm/h). \textbf{Bandes spirales} (rainbands) : convergence basse, ascendance, orages, vents forts par rafales, organisation spirale logarithmique.
    \item \textbf{Vertical :} Convergence basses couches (0--3 km) $\to$ Ascendance cœur (mur œil) jusqu'à tropopause (15--18 km) $\to$ Divergence haute (épanchement cirrus/épais nuageux en enclume) $\to$ Subsidence dans l'œil et à la périphérie lointaine. \textbf{Cœur chaud} (anomalie $+10$ à $+15°C$ à 250 hPa vs environnement).
\end{itemize}

\textbf{4. Trajectoire et dangers.}
\begin{itemize}
    \item \textbf{Trajectoire :} Guidage par flux directeurs (vent moyen 850--300 hPa). Initialement porté par \textbf{alizés} (Est $\to$ Ouest) vers les Petites Antilles. Approche latitudes moyennes ($\sim$ 25--30° N) $\to$ rencontre \textbf{vents d'ouest} $\to$ \textbf{recourbement parabolique} vers le Nord puis Nord-Est. Déviation $\beta$ (Coriolis variable) pousse vers le Nord-Ouest naturellement.
    \item \textbf{3 dangers majeurs :}
    \begin{enumerate}
        \item \textbf{Surcote (Storm surge) :} Élévation niveau mer = Basse pression (effet barométrique inversé : 1 hPa $\approx$ 1 cm) + Effet vent (poussée eau vers côte) + Vague déferlante. Peut atteindre 3--6 m (cat 3--5). \textbf{Cause n°1 de mortalité.}
        \item \textbf{Vents violents :} Destruction habitat précaire, toitures, lignes électriques, projectiles. Échelle Saffir-Simpson Cat 3 = « Dévastateur » (dégâts structurels petits bâtiments, arbres déracinés).
        \item \textbf{Pluies diluviennes / Inondations :} Cumul 200--500 mm en 24h (lent déplacement $\to$ plus de pluie). Crues éclair, glissements terrain (îles volcaniques : Montagne Pelée, Soufrière).
    \end{enumerate}
\end{itemize}

\textbf{Conclusion.} Relier le cyclone à son origine dynamique, c'est enchaîner : Conditions initiales (thermodynamique + dynamique) $\to$ Moteur thermique (chaleur latente) $\to$ Structure (cœur chaud, gradient pression) $\to$ Trajectoire (flux directeurs + $\beta$-dérive) $\to$ Aléas (surcote, vent, pluie).
\end{exoresolu}

% ============================================================
% SAVOIR-FAIRE 2 (bis) : Relier une tempête tempérée (dépression des latitudes moyennes) à son origine
% ============================================================

\begin{methode}[Analyser une dépression tempérée (tempête) sur une carte météorologique de surface]
\textbf{Objectif :} Lier la structure d'une dépression des latitudes moyennes (fronts, secteurs) à la dynamique du \cle{Front polaire} et au \cle{Courant-jet}.
\begin{enumerate}
    \item \textbf{Identifier le centre de basse pression (L) :} Pression centrale (ex: 980 hPa). Isobares serrées $\to$ vent fort (gradient).
    \item \textbf{Repérer les fronts :} \textbf{Front chaud} (ronds rouges, air chaud remplace froid, pente douce 1/100--1/300, nuages étagés Cs/As/Ns $\to$ pluies continues) ; \textbf{Front froid} (triangles bleus, air froid remplace chaud, pente raide 1/50--1/100, nuages convectifs Cb $\to$ averses/orages violents, rafales) ; \textbf{Front occlus} (mélange, symbole mixte, stade mature/déclin).
    \item \textbf{Localiser les secteurs :} \textbf{Secteur chaud} (entre front chaud et front froid, air tropical maritime, doux, humide, vent fort de SO) ; \textbf{Secteur froid} (arrière front froid, air polaire maritime, instable, averses).
    \item \textbf{Relier au Courant-jet (Jet-stream) :} Creux de vague de Rossby en altitude $\to$ divergence en altitude (sortie gauche / entrée droite du jet) $\to$ aspiration air en surface $\to$ creusement dépression (cyclogenèse explosive si divergence forte + advection chaleur basse couche + advection froid haute couche).
    \item \textbf{Évolution type (Modèle Norvégien - Bjerknes) :} Ondulation front $\to$ Jeune dépression (fronts chaud/froid distincts) $\to$ Mature (occlusion commence, pression min, vents max) $\to$ Occlus (front occlus, centre comblé, remplissage).
\end{enumerate}
\textbf{Reflexes :} « Tempête » = vent moyen $\ge$ 89 km/h (Force 10 Beaufort) ou rafales $\ge$ 100 km/h. Ne pas confondre avec cyclone tropical (pas d'œil, fronts présents, cœur froid, asymétrique, énergie = énergie potentielle disponible barocline, pas chaleur latente seule).
\end{methode}

\begin{exoresolu}[Analyse d'une tempête hivernale sur l'Atlantique Nord touchant l'Europe]
\textbf{Énoncé.} Carte météo surface 12h UTC, 15 janvier. Dépression centrée 55° N, 20° W (972 hPa). Front chaud vers 50° N (Biscaye), front froid vers 45° N (Açores). Courant-jet à 300 hPa : creux marqué sur l'Atlantique, vitesse max 180 km/h (50 m/s) en sortie de creux vers l'Europe. Vent de surface au large Bretagne : SO 90 km/h, rafales 130 km/h.
\textbf{Questions :}
\begin{enumerate}
    \item Nomme le type de dépression et son stade d'évolution d'après le modèle de Bergen.
    \item Explique le rôle du courant-jet (niveau 300 hPa) dans l'approfondissement (creusement) de cette dépression.
    \item Décris le temps qu'il fait dans le secteur chaud et derrière le front froid pour un observateur à Brest (Bretagne) lors du passage des fronts.
    \item Pourquoi cette dépression est-elle classée « tempête » ? Quelle est la différence fondamentale de moteur énergétique avec un cyclone tropical ?
\end{enumerate}

\tcblower
\textbf{Solution détaillée.}

\textbf{1. Type et stade.}
C'est une \textbf{dépression tempérée (extratropicale) des latitudes moyennes}, aussi appelée \textbf{perturbation atlantique} ou \textbf{dépression du front polaire}.
Stade : \textbf{Mature (début d'occlusion)}. Le front froid rattrape le front chaud (vitesse front froid $\approx$ 2x vitesse front chaud). L'occlusion commence près du centre (point triple). Pression minimale atteinte (972 hPa). Gradient de pression maximal $\to$ vents max.

\textbf{2. Rôle du Courant-jet (Niveau 300 hPa $\approx$ 9000 m).}
\begin{itemize}
    \item Le Jet-stream n'est pas rectiligne : il ondule (\cle{Vagues de Rossby}). La dépression de surface est située sous l'\textbf{avant d'un creux} (axe creux $\to$ Ouest de la dépression).
    \item Dans la \textbf{sortie gauche du jet} (ou entrée droite), la dynamique (équilibre quasi-géostrophique) impose une \textbf{divergence en altitude} (l'air s'éloigne plus vite qu'il n'arrive).
    \item Cette divergence en altitude « aspire » l'air des basses couches vers le haut (mouvement vertical synoptique $\omega < 0$).
    \item Pour compenser la perte de masse en altitude, la pression au sol \textbf{chute} (creusement) : $\frac{\partial p_s}{\partial t} \propto - \int \nabla \cdot \vec{V} \, dp$. C'est la \textbf{cyclogenèse explosive} (bombe météorologique si $\Delta p / 24h \ge 24$ hPa corrigé latitude).
    \item \textbf{Advections couplées :} Advection d'air chaud (AAC) en basses couches (devant front chaud) + Advection d'air froid (AAF) en haute couche (arrière creux) $\to$ instabilité barocline $\to$ conversion énergie potentielle $\to$ énergie cinétique (vent).
\end{itemize}

\textbf{3. Temps à Brest (passage perturbation Ouest $\to$ Est).}
\begin{center}
\begin{tabularx}{\linewidth}{|l|X|X|X|}\hline
\textbf{Phase} & \textbf{Pression} & \textbf{Vent / Temps} & \textbf{Phénomènes significatifs} \\ \hline
\textbf{Approche Front Chaud} & Baisse régulière & Vent SE $\to$ S, se renforce. Ciel : Ci $\to$ Cs $\to$ As $\to$ Ns. & Pluies continues, modérées, prolongées (6--12h). Nuages bas (brouillard/stratus). Douceur. \\ \hline
\textbf{Passage Front Chaud} & Palier / inflexion & Tourne SO brutalement. Rafales. & Cesse-pluie temporaire. Température $\uparrow$ brutale (secteur chaud). Humidité $\uparrow$. \\ \hline
\textbf{Secteur Chaud} & Baisse lente / stable & Vent SO \textbf{Fort à Tempête} (90 km/h, rafales 130 km/h). Ciel bas (Sc/St) ou dégagé. & Temps lourd, doux, humide. Risque rafales violentes (cisaillement vent vertical). \\ \hline
\textbf{Passage Front Froid} & \textbf{Hausse brutale} (saut barométrique) & Tourne SO $\to$ O $\to$ NO. Rafales max (ligne de grain). & \textbf{Averses violentes, orages, grêle possible, rafales descendantes (downbursts).} Chute T° brutale. \\ \hline
\textbf{Arrière Front Froid (Secteur Froid)} & Hausse régulière & Vent O/NO fort, rafaleux. Ciel : Cu, Cb épars. & Averses (grains), éclaircies. Instabilité (air froid sur mer « chaude »). Tempête de neige possible si T° $\le$ 0°C. \\ \hline
\end{tabularx}
\end{center}

\textbf{4. Classification « Tempête » et différence moteur.}
\begin{itemize}
    \item \textbf{Critère tempête :} Vent moyen $\ge$ 89 km/h (Force 10 Beaufort) ou Rafales $\ge$ 100 km/h. Ici : Vent moyen 90 km/h, Rafales 130 km/h $\to$ \textbf{Tempête confirmée} (ex: type Lothar/Martin 1999, Xynthia 2010).
    \item \textbf{Moteur énergétique :}
    \begin{itemize}
        \item \textbf{Cyclone Tropical :} \textbf{Thermodynamique (Cœur Chaud)}. Énergie = Chaleur latente condensation (océan chaud). Symétrique. Pas de fronts. Nécessite Coriolis $\neq$ 0 mais faible cisaillement.
        \item \textbf{Dépression Tempérée :} \textbf{Barocline (Cœur Froid)}. Énergie = \textbf{Énergie Potentielle Disponible (EPD)} du gradient de température horizontal (air polaire froid vs air tropical chaud). Asymétrique (fronts). Nécessite fort cisaillement vertical (Jet) et fort gradient T° horizontal. \textbf{Conversion EPD $\to$ Énergie Cinétique} (instabilité barocline d'Eady/Charney).
    \end{itemize}
\end{itemize}

\textbf{Conclusion.} L'analyse d'une tempête tempérée repose sur la géométrie des fronts (modèle norvégien), la dynamique du jet-stream (divergence altitude $\to$ creusement surface) et le contraste de masses d'air (baroclinicité). Les dangers : vents violents (secteur chaud + arrière front froid), pluie/neige, surcote (marée + basse pression + vent).
\end{exoresolu}

% ============================================================
% SAVOIR-FAIRE 3 : Proposer des mesures de protection face aux phénomènes dangereux
% ============================================================

\begin{methode}[Proposer des mesures de prévention et protection adaptées (cyclones, tempêtes, submersions)]
\textbf{Objectif :} Structurer une réponse opérationnelle selon le cycle de la gestion des risques : \textbf{Prévention (amont)} -- \textbf{Préparation/Alerte (imminent)} -- \textbf{Protection/Gestion de crise (pendant)} -- \textbf{Réhabilitation (après)}.
\begin{enumerate}
    \item \textbf{Analyser l'aléa et l'enjeu :} Type de phénomène (cyclone, tempête tempérée, submersion marine, crue éclair). Zone concernée (urbain, rural, littoral, montagne). Population vulnérable (précarité, handicap, isolement).
    \item \textbf{Mesures de Prévention Structurelle (Long terme / Aménagement) :}
    \begin{itemize}
        \item \textbf{Règlementation / Urbanisme :} \cle{PPRI} (Plan de Prévention des Risques d'Inondation/Submersion), \cle{POS/PLU} interdisant construction en zone rouge (aléa fort), prescriptions constructives (surélévation, matériaux résistants, ouvertures protégées).
        \item \textbf{Ouvrages de protection :} Digues, brise-lames, rechargement de plages, restauration mangroves/récifs coralliens (solutions fondées sur la nature / \cle{SbN}), bassins de rétention.
        \item \textbf{Réseaux d'observation :} Radars météo (Doppler), bouées, satellites, stations côtières (marégraphes).
    \end{itemize}
    \item \textbf{Mesures de Préparation et Alerte (Court terme / J-3 à H-0) :}
    \begin{itemize}
        \item \textbf{Vigilance Météo (Météo-France / DGM Cameroun) :} Niveaux Vert/Jaune/Orange/Rouge. Comprendre la signification de chaque niveau.
        \item \textbf{Plans d'urgence :} \cle{PCS} (Plan Communal de Sauvegarde) obligatoire communes à risque. \cle{PIMS} (Plan d'Intervention de Mise en Sûreté) pour établissements recevant du public (écoles, hôpitaux).
        \item \textbf{Information population :} Consignes (kit survie 72h, mise en sécurité objets, évacuation bétail, bouchage ouvertures).
        \item \textbf{Évacuation préventive :} Zones basses, mobile-homes, habitats précaires, avant H-0 (vents $>$ 80 km/h rendent évacuation impossible).
    \end{itemize}
    \item \textbf{Mesures de Protection / Gestion de Crise (Pendant l'événement) :}
    \begin{itemize}
        \item \textbf{Confinement :} Rester à l'intérieur, pièce centrale sans fenêtre, couper gaz/électricité si inondation imminente.
        \item \textbf{Ne pas sortir :} Risque chutes arbres, tuiles, câbles, hydroplanage, surcote soudaine.
        \item \textbf{Suivi info officielles :} Radio (batterie), alertes SMS (FR-Alert / systèmes locaux), comptes officiels.
    \end{itemize}
    \item \textbf{Mesures Post-Crise (Résilience) :} Évaluation dégâts, déblaiement, épidémiosurveillance (eau, moustiques), soutien psychologique, reconstruction « mieux » (Build Back Better), retour d'expérience (REX).
\end{enumerate}
\textbf{Reflexes Probatoire :} Distinguer \textbf{Prévention} (réduire vulnérabilité/exposition AVANT) et \textbf{Protection} (réduire conséquences PENDANT). Citer des exemples camerounais : Plan de contingence inondations Extrême-Nord, surveillance volcan Mont Cameroun, gestion risques côtiers Kribi/Limbé/Douala. Vocabulaire : \cle{Aléa}, \cle{Vulnérabilité}, \cle{Exposition}, \cle{Risque = Aléa $\times$ Vulnérabilité $\times$ Exposition}, \cle{Résilience}.
\end{methode}

\begin{exoresolu}[Élaboration d'un plan de protection communal face au risque cyclone/surcote à Limbé (Cameroun)]
\textbf{Énoncé.} La ville de Limbé (Sud-Ouest, Cameroun), située au pied du Mont Cameroun, sur la côte atlantique, est exposée aux cyclones tropicaux (ex: tempête tropicale \textit{« ... »} hypothétique cat 2) et aux fortes pluies orographiques. Le quartier \textit{Bota Land} est en zone basse (mangrove remblayée), habitat précaire. L'hôpital de district est à 200 m de la mer. La route d'accès unique (RN3) traverse une zone de glissement de terrain.
\textbf{Questions :}
\begin{enumerate}
    \item Identifie les 3 composantes du risque (Aléa, Vulnérabilité, Exposition) pour ce scénario.
    \item Propose 3 mesures de \textbf{prévention structurelle} (aménagement/urbanisme) adaptées à ce contexte.
    \item Détaille le dispositif d'\textbf{alerte et d'évacuation} (niveaux vigilance, acteurs, moyens, lieux d'hébergement) pour une vigilance Orange « Cyclone » émise à J-2.
    \item Quelles consignes de \textbf{protection individuelle} donner à la population pendant le passage de l'œil (calme trompeur) et du mur de l'œil ?
\end{enumerate}

\tcblower
\textbf{Solution détaillée.}

\textbf{1. Analyse du Risque (R = A $\times$ V $\times$ E).}
\begin{itemize}
    \item \textbf{Aléa (A) :} \textbf{Cyclone tropical Cat 2} (vents 154--177 km/h, surcote 1,5--2,5 m, pluies $>$ 300 mm/24h) + \textbf{Glissements de terrain} (pluies orographiques Mont Cameroun, sols volcaniques instables) + \textbf{Crues éclairs} (bassins versants courts, pentus).
    \item \textbf{Vulnérabilité (V) :} \textbf{Forte}. Habitat précaire (tôles, bois, parpaings non ferraillés) $\to$ destruction vent/surcote. Réseaux (électricité, eau, assainissement) vétustes. Hôpital non surélevé, non parasismique/paracyclonique. Voie d'accès unique (RN3) = point de rupture unique (SPOF).
    \item \textbf{Exposition (E) :} \textbf{Forte}. Population dense (Bota Land $\sim$ 20 000 hab), activités économiques (port, CDC plantations, raffinerie SONARA à proximité), infrastructures critiques (hôpital, écoles) en zone inondable/surcotable.
\end{itemize}

\textbf{2. Mesures de Prévention Structurelle (Long terme).}
\begin{enumerate}
    \item \textbf{Révision du PLU / PPRI :} Classement \textbf{Zone Rouge (inconstructible)} pour Bota Land (aléa surcote + glissement). \textbf{Relocalisation progressive} des populations vers zones sûres (hauteurs, zones d'extension urbainisées avec voiries/réseaux). Indemnisation / accompagnement social indispensable.
    \item \textbf{Protection de l'hôpital (Infrastructure critique) :} \textbf{Surélévation} du plateau technique (urgences, blocs, groupes électrogènes, réserve eau) au-dessus cote de surcote centennale + 0,5 m (marge climatique). \textbf{Bardage pare-cyclonique} (baies vitrées, toiture). Accès héliportable sécurisé. Stocks autonomes 72h (eau, carburant, médicaments, sang).
    \item \textbf{Sécurisation de l'accès (RN3) :} \textbf{Ouvrages de stabilisation} (parois clouées, drains profonds, filets pare-blocs) sur les secteurs de glissement identifiés. \textbf{Création d'une voie de contournement} (piste de secours) par les hauteurs (Mabeta, Batoke) pour désenclaver si RN3 coupée. \textbf{Entretien rigoureux} fossés/caniveaux (désenvasement pré-saison).
\end{enumerate}

\textbf{3. Dispositif Alerte / Évacuation (Vigilance Orange Cyclone à J-2).}
\begin{itemize}
    \item \textbf{J-2 (Orange) :} \textbf{Maire} (chef opérations PCS) active \textbf{Cellule de Crise Communale}. Réunion \textbf{CODIS} (Comité Opérationnel Départemental Incendie Secours) / Préfet.
    \begin{itemize}
        \item Diffusion alerte : Sirènes communales, SMS géolocalisés (si système existe), Radio locale (Ocean City Radio, CRTV), crieurs publics, relais associatifs/chefs de quartiers.
        \item Message type : « Vigilance Orange Cyclone. Risque vents destructeurs, surcote, inondations. Préparer kit 72h. Évacuation zone Bota Land / bord de mer PRÉVUE demain J-1 6h. »
        \item Pré-positionnement secours : Sapeurs-pompiers, Protection Civile, Croix-Rouge, armée (BIR/BIA si réquisition) vers zones d'hébergement.
    \end{itemize}
    \item \textbf{J-1 (Matin, passage Rouge probable) :} \textbf{Ordre d'évacuation OBLIGATOIRE} zones basses (Bota Land, Down Beach, Mile 4).
    \begin{itemize}
        \item \textbf{Moyens :} Bus municipaux, camions bennes, pirogues motorisées (si routes coupées), marche à pied encadrée.
        \item \textbf{Lieux d'hébergement (Centres d'Hébergement d'Urgence - CHU) :} Écoles/lycées en zone haute (Limbé Lycée, Molyko Buea), églises, salles des fêtes. Capacité, eau, latrines, infirmerie vérifiées AVANT.
        \item \textbf{Registre :} Pointage nominatif évacués (listes chefs de quartiers) pour recherche disparus post-événement.
    \end{itemize}
    \item \textbf{H-0 (Début vents violents) :} \textbf{Confinement total}. Plus aucune évacuation. Secours en attente abris renforcés.
\end{itemize}

\textbf{4. Consignes Protection Individuelle (Pendant le passage).}
\begin{itemize}
    \item \textbf{Passage de l'ŒIL (Calme, ciel clair, 15--30 min) :} \textbf{NE PAS SORTIR.} C'est le piège mortel n°1. Le vent va reprendre \textbf{violemment en sens inverse} (arrière mur de l'œil) en quelques minutes. Profiter pour : vérifier étanchéité (bâches), vider eau entrée, écouter radio, \textbf{ne pas aller voir dehors}, ne pas commencer réparations.
    \item \textbf{Passage du Mur de l'ŒIL (Vents max, pluies diluviennes) :} Se mettre \textbf{au centre du bâtiment}, pièce sans fenêtre (couloir, salle de bain, WC), \textbf{à l'abri des débris volants} (matelas, table solide au-dessus). Couper électricité/gaz si eau monte. S'éloigner des portes/fenêtres. Protéger tête/cou. Ne pas utiliser ascenseur. Ne pas prendre voiture.
    \item \textbf{Surcote soudaine :} Si eau monte vite dans la maison $\to$ \textbf{Monter à l'étage / toit} (dernier refuge), emporter gilets de sauvetage / bidons vides / téléphone. Ne pas nager dans l'eau courante (courants, débris).
\end{itemize}

\textbf{Conclusion.} La protection efficace repose sur l'anticipation (urbanisme résilient, PCS testé), une chaîne d'alerte fiable et comprise de tous (y compris populations analphabètes/anglophones/francophones), une évacuation PRÉCOCE (avant vents forts) vers des lieux sûrs identifiés, et des consignes de survie claires pendant l'événement (piège de l'œil). La résilience post-crise (reconstruction durable, soutien psycho-social) clôt le cycle.
\end{exoresolu}

% ============================================================
% EXERCICE COMPLÉMENTAIRE : Interactions Océan-Atmosphère (El Niño / La Niña) - Savoir-faire 1 & 2
% ============================================================

\begin{methode}[Analyser les anomalies ENSO (El Niño / La Niña) et leurs téléconnexions]
\textbf{Objectif :} Relier l'état de l'océan Pacifique équatorial (TSM, alizés, thermocline) aux impacts climatiques mondiaux (sécheresses, inondations).
\begin{enumerate}
    \item \textbf{Identifier la phase :} \textbf{Neutre} (Alizés Est $\to$ Ouest forts, thermocline penchée Est bas / Ouest haut, upwelling Est, piscine chaude Ouest) ; \textbf{El Niño} (Alizés faiblissent/inversent, piscine chaude migre Est, thermocline s'aplatit, upwelling stoppé Est, convection suit chaleur $\to$ pluies Équateur Est/Pérou) ; \textbf{La Niña} (Alizés renforcés, upwelling intensifié Est, piscine chaude confinée Ouest, convection renforcée Ouest $\to$ pluies Indonésie/Australie).
    \item \textbf{Quantifier :} Indice \cle{ONI} (Oceanic Niño Index) = Moyenne mobile 3 mois anomalie TSM région Niño 3.4 (5°N-5°S, 170°-120°W). El Niño si ONI $\ge$ +0,5°C (5 trimestres) ; La Niña si ONI $\le$ -0,5°C.
    \item \textbf{Téléconnexions majeures (Impacts Probatoire) :}
    \begin{itemize}
        \item \textbf{El Niño :} Sécheresse \textbf{Australie, Indonésie, Nord Brésil, Afrique Australe, Sahel (souvent)}. Pluies excédentaires \textbf{Pérou/Équateur, Sud USA, Corne de l'Afrique (souvent)}. Cyclones Atlantique \textbf{réduits} (cisaillement fort).
        \item \textbf{La Niña :} Inondations \textbf{Australie, Indonésie, Nord-Est Brésil}. Sécheresse \textbf{Pérou/Équateur, Sud USA, Corne de l'Afrique}. Cyclones Atlantique \textbf{augmentés} (cisaillement faible). Hivers plus froids Europe/Amérique Nord.
    \end{itemize}
    \item \textbf{Contexte Cameroun :} El Niño $\to$ souvent déficit pluviométrique au Nord (Sahel) et sur la côte (Douala/Limbé) ; La Niña $\to$ souvent excédent pluviométrique (risque inondations Extrême-Nord, Logone, Bénoué).
\end{enumerate}
\textbf{Reflexes :} Ne pas inverser les signes d'anomalie TSM. « El Niño = Réchauffement Est Pacifique ». Ne pas généraliser : « souvent », « tendance », « probabilité accrue », pas certitude absolue. Distinguer impact direct (Pacifique) et téléconnexion (atmosphérique via circulation de Walker modifiée).
\end{methode}

\begin{exoresolu}[Analyse d'un épisode El Niño 2023-2024 et conséquences au Cameroun]
\textbf{Énoncé.} En décembre 2023, l'indice ONI (Niño 3.4) atteint +2,0°C (El Niño fort). Les alizés sont inversés (vents d'Ouest) sur le Pacifique occidental. La thermocline est plate. La ZCIT est déplacée vers l'équateur central Pacifique.
\textbf{Questions :}
\begin{enumerate}
    \item Décris les anomalies océaniques et atmosphériques dans le Pacifique équatorial lors de cet El Niño fort (comparé à la normale).
    \item Explique le mécanisme de la \textbf{téléconnexion} vers l'Atlantique tropical et l'Afrique de l'Ouest/Cameroun (circulation de Walker / Hadley modifiée).
    \item Prédis les impacts probables sur la campagne agricole 2024 au Cameroun (Nord vs Sud) et sur l'activité cyclonique en Atlantique Nord.
    \item Propose une mesure d'adaptation pour les agriculteurs du Nord-Cameroun (région de l'Extrême-Nord).
\end{enumerate}

\tcblower
\textbf{Solution détaillée.}

\textbf{1. Anomalies Pacifique Équatorial (El Niño fort).}
\begin{center}
\begin{tabularx}{\linewidth}{|l|X|X|}\hline
\textbf{Paramètre} & \textbf{Conditions Normales (Neutre)} & \textbf{Anomalies El Niño Fort (Déc 2023)} \\ \hline
\textbf{Alizés (Surface)} & Forts, Est $\to$ Ouest (Convergence Ouest) & \textbf{Affaiblis / Inversés (Ouest $\to$ Est)} sur Pacifique Centre/Ouest \\ \hline
\textbf{TSM (Température Surface Mer)} & Froide Est (Upwelling) / Chaude Ouest (Piscine > 29°C) & \textbf{Anomalie positive majeure (+2 à +3°C) Est/Centre} (Niño 3.4). Piscine chaude migre Est. \\ \hline
\textbf{Thermocline} & Penchée : Profonde Ouest (150m), Faible Est (30m) & \textbf{Aplanlie / Profonde partout} (Ondes de Kelvin descendantes). Upwelling Est \textbf{stoppé}. \\ \hline
\textbf{Convection / Pluies (ZCIT)} & Concentrée Ouest (Indonésie, Australie) & \textbf{Déplacée vers Centre/Est Pacifique} (Ligne de changement date). Pluies diluviennes Pérou/Équateur. \\ \hline
\textbf{Pression (SOI)} & Haute Est (Darwin) / Basse Ouest (Tahiti) & \textbf{Inversée : Basse Est / Haute Ouest} (SOI très négatif). \\ \hline
\end{tabularx}
\end{center}

\textbf{2. Mécanisme de Téléconnexion vers Atlantique / Afrique.}
\begin{itemize}
    \item \textbf{Circulation de Walker modifiée :} L'ascendance massive déplacée vers le Pacifique Central/Est crée une \textbf{cellule de Walker anormale}.
    \item \textbf{Subsidence compensatrice :} La masse d'air montant au Pacifique doit redescendre ailleurs. Elle descend préférentiellement sur l'\textbf{Atlantique tropical Ouest} (Nord-Est Brésil, Golfe de Guinée) et l'\textbf{Océan Indien Ouest}.
    \item \textbf{Conséquence dynamique :} Subsidence $\to$ \textbf{stabilisation atmosphérique} (inversion de température, air sec en moyenne troposphère) $\to$ \textbf{inhibition de la convection} (ZCIT africaine affaiblie, déplacée au Sud).
    \item \textbf{Circulation de Hadley :} Renforcement de la branche descendante Hadley sur l'Atlantique tropical $\to$ Alizés plus forts / plus secs sur le Golfe de Guinée $\to$ Moins d'humidité advectée vers le continent.
    \item \textbf{Cisaillement vertical Atlantique Nord :} Renforcement des vents d'ouest en altitude (200 hPa) via téléconnexion $\to$ \textbf{Cisaillement vertical fort} $\to$ inhibe cyclogenèse cyclones tropicaux Atlantique.
\end{itemize}

\textbf{3. Impacts Prédictifs Cameroun 2024.}
\begin{itemize}
    \item \textbf{Nord-Cameroun (Extrême-Nord, Nord, Adamaoua - Sahel/Soudanien) :} \textbf{Risque accru de DÉFICIT pluviométrique.} Début saison des pluies retardé, poches de sécheresse, fins de saison précoces. Impact : \textbf{Baisse rendements mil/sorgho/maïs/arachide}, stress hydrique pâturages (conflits éleveurs/agriculteurs), insécurité alimentaire (déjà fragilisée par insécurité Boko Haram).
    \item \textbf{Sud-Cameroun (Zone équatoriale / forêt / littoral) :} Signal moins net. Légère tendance à déficit sur le littoral (Douala, Limbé, Kribi) mais variabilité interannuelle forte. Saison sèche courte (janv-fév) possiblement plus sèche/chaude.
    \item \textbf{Activité Cyclonique Atlantique Nord 2024 :} \textbf{Activité INFÉRIEURE à la normale} (cisaillement fort, air sec saharien (SAL) favorisé par subsidence). \textit{Note réelle : Saison 2024 a été paradoxalement très active (record TSM Atlantique dominant sur cisaillement El Niño), illustrer l'incertitude.}
\end{itemize}

\textbf{4. Mesure d'Adaptation Agriculteurs Extrême-Nord.}
\textbf{Mesure :} \textbf{Promotion variétés à cycle court / tolérantes à la sécheresse (Mil Souna 3, Sorgho S35, Maïs TZEE-W) + Agroforesterie (Faidherbia albida / Acacia) + Zaï / Demi-lunes.}
\begin{itemize}
    \item \textbf{Variétés courtes (80--90 j) :} Échappent à la fin de saison précoce. Tolérance stress hydrique floraison.
    \item \textbf{Zaï / Demi-lunes :} Techniques de \textbf{récupération eaux de ruissellement} (fossés, trous amendés fumure) $\to$ concentration eau/fertilité au pied plante. Adapté sols dégradés (glacis).
    \item \textbf{Agroforesterie (Faidherbia) :} Arbre perd feuilles saison humide (pas d'ombre) $\to$ feuilles saison sèche (pâturage, fertilité azotée, ombre bétail). Améliore structure sol, rétention eau.
    \item \textbf{Appui institutionnel :} Distribution semences subventionnées (SODECOTON, MINADER, ONG) AVANT saison (avril). Information climatique (PRESASS, ACMAD) diffusée par radio locale (langues fulfulde, arabe choua) pour décision semis.
\end{itemize}

\textbf{Conclusion.} L'analyse ENSO demande de décrire l'anomalie physique Pacifique (Océan + Atmosphère couplés), d'expliquer la téléconnexion atmosphérique (Walker/Hadley modifiées $\to$ subsidence téléconnectée), et de décliner les impacts probabilistes sectoriels (agriculture, eau, santé, cyclones) pour proposer des adaptations concrètes, locales et équitables.
\end{exoresolu}

% ============================================================
% BOÎTE ATTENTION FINALE : LES 5 ERREURS QUI COÛTENT DES POINTS
% ============================================================

\begin{attention}[Les 5 erreurs qui coûtent des points au Probatoire - NCT07]
\begin{enumerate}
    \item \textbf{Confondre Cyclone Tropical et Dépression Tempérée (Tempête).}
    \textit{Erreur :} Dire qu'un cyclone a des fronts, ou qu'une tempête tempérée a un œil et un cœur chaud.
    \textit{Correct :} \textbf{Cyclone tropical} = Cœur chaud, symétrique, \textbf{sans fronts}, énergie = chaleur latente (océan $\ge$ 26,5°C), Coriolis nécessaire mais cisaillement faible. \textbf{Dépression tempérée} = Cœur froid, \textbf{fronts chaud/froid/occlus}, asymétrique, énergie = instabilité barocline (gradient T° horizontal + Jet-stream), cisaillement fort requis.
    
    \item \textbf{Oublier la Force de Coriolis ou inverser le sens de déviation.}
    \textit{Erreur :} « Les alizés soufflent du Nord au Sud » (sans déviation). Ou « Déviation à gauche dans l'hémisphère Nord ».
    \textit{Correct :} Coriolis dévie \textbf{vers la DROITE dans l'hémisphère Nord} (alizés $\to$ NE, vents d'ouest $\to$ SO), \textbf{vers la GAUCHE dans l'hémisphère Sud} (alizés $\to$ SE, vents d'ouest $\to$ NO). Nulle à l'équateur. C'est la clé de la rotation des gyres (horaire N, anti-horaire S) et des cyclones (anti-horaire N, horaire S).
    
    \item \textbf{Mélanger Upwelling (Surgences) et Downwelling, ou mal localiser l'upwelling.}
    \textit{Erreur :} Placer l'upwelling sur les côtes Ouest (bordure ouest gyres) ou dire qu'il réchauffe l'eau.
    \textit{Correct :} Upwelling majeur = \textbf{Côtes EST des océans} (Benguela, Canaries, Humboldt, Californie, Somalie/Oman). Mécanisme : Vent parallèle à côte + Coriolis $\to$ Transport d'Ekman \textbf{au large} $\to$ Remontée eau \textbf{FROIDE}, \textbf{RICHE en nutriments}. Conséquence : Productivité $\uparrow$, Pêche $\uparrow$, Climat côtier \textbf{froid et sec} (déserts côtiers : Atacama, Namib, Sahara atlantique).
    
    \item \textbf{Ne pas structurer la réponse « Mesures de protection » selon le cycle de gestion du risque.}
    \textit{Erreur :} Lister des mesures en vrac (« construire des digues, écouter la radio, évacuer ») sans distinguer Prévention (long terme), Préparation/Alerte (court terme), Protection (pendant), Réhabilitation (après).
    \textit{Correct :} Utiliser le cadre \textbf{PPRI / PCS / Vigilance / Confinement / REX}. Citer les acteurs (Maire, Préfet, Météo-France/DGM, Protection Civile, Population). Distinguer \textbf{structurel} (urbanisme, digues, SbN) et \textbf{non-structurel} (alerte, assurance, éducation, plan d'urgence). Contextualiser (Cameroun : zones inondables Logone/Bénoué, risque volcanique Mont Cameroun, risque côtier Douala/Kribi/Limbé).
    
    \item \textbf{Imprécision sur les latitudes limites des cellules et la ZCIT.}
    \textit{Erreur :} « La cellule de Hadley va de l'équateur aux pôles » ou « La ZCIT est fixe à l'équateur ».
    \textit{Correct :} Cellules : \textbf{Hadley 0°--30°}, \textbf{Ferrel 30°--60°}, \textbf{Polaire 60°--90°}. \textbf{ZCIT} = Zone de Convergence Intertropicale. \textbf{Mobile} : suit le soleil zénital (migration saisonnière $\sim$ 5°--10° N en juil, $\sim$ 5° S en janv). Au Cameroun : ZCIT au Nord (12°--14° N) en août $\to$ saison des pluies Nord ; ZCIT au Sud (2°--4° N) en janv $\to$ saison sèche Nord, petite saison sèche Sud. La position de la ZCIT commande le régime des pluies.
\end{enumerate}
\textbf{Astuce finale :} Sur un schéma, \textbf{légende TOUT} (flèches vents, isothermes, isobares, fronts, courants, upwelling, cellules). Une flèche sans nom = 0 point. Une latitude approximative (0, 30, 60, 90) = point bonus. Courage, la rigueur paie !
\end{attention}

\newpage

\coursec{Exercices}

\courssub{Je vérifie mes connaissances}

\begin{exercice}[QCM : Circulation atmosphérique générale]
Parmi les affirmations suivantes, cochez celle(s) qui est(sont) correcte(s) :
\begin{enumerate}
    \item L'air chaud, moins dense, s'élève en créant une zone de haute pression au sol.
    \item La cellule de Hadley s'étend de l'équateur jusqu'aux environs de 30° de latitude Nord et Sud.
    \item Les alizés sont des vents d'Ouest qui soufflent des latitudes moyennes vers l'équateur.
    \item L'effet de Coriolis dévie les vents vers la droite dans l'hémisphère Nord et vers la gauche dans l'hémisphère Sud.
    \item Les zones de convergence intertropicale (ZCIT) correspondent à des zones de subsidence (air descendant).
\end{enumerate}
\end{exercice}

\begin{exercice}[Vrai ou Faux justifié : Courants océaniques]
Indiquez si les affirmations sont vraies (V) ou fausses (F) et justifiez brièvement votre réponse.
\begin{enumerate}
    \item Le Gulf Stream est un courant froid qui descend de l'Arctique vers l'Europe.
    \item Le courant de Benguela, au large de l'Afrique australe, est un courant de remontée (\emph{upwelling}) riche en nutriments.
    \item Les courants marins de surface sont principalement entraînés par les différences de densité de l'eau (thermohaline).
    \item L'El Niño correspond à un réchauffement anormal des eaux de surface dans le Pacifique équatorial Est.
    \item La circulation thermohaline globale (le « tapis roulant ») se fait sur une échelle de temps de quelques mois.
\end{enumerate}
\end{exercice}

\begin{exercice}[Définitions et vocabulaire]
Donnez une définition précise et rigoureuse pour chacun des termes suivants :
\begin{enumerate}
    \item Cellule de convection atmosphérique (citer les trois cellules principales).
    \item Effet de Coriolis.
    \item Courant de dérive (ou courant de surface) vs Courant de profondeur.
    \item Cyclone tropical (conditions de formation minimales : température de l'eau, latitude, cisaillement).
    \item Surcote marine (storm surge).
\end{enumerate}
\end{exercice}

\begin{exercice}[Calcul direct : Vitesse et trajet d'un courant]
Le courant du Gulf Stream a une vitesse moyenne de \SI{1,5}{\meter\per\second} au large des côtes américaines.
\begin{enumerate}
    \item Exprimez cette vitesse en \SI{}{\kilo\meter\per\hour} et en \SI{}{\kilo\meter\per\day}.
    \item Une masse d'eau met environ 60 jours pour traverser l'Atlantique Nord (distance approximative \SI{5000}{\kilo\meter}) via le Gulf Stream et le courant du Nord-Atlantique. Calculez la vitesse moyenne effective de ce trajet en \SI{}{\centi\meter\per\second}. Comment expliquez-vous la différence avec la vitesse maximale du Gulf Stream ?
\end{enumerate}
\end{exercice}

\courssub{Je m'entraîne}

\begin{exercice}[Analyse de carte : L'Atlantique Nord]
On vous donne une carte muette de l'Atlantique Nord montrant les isothermes de surface et les principaux courants (Gulf Stream, Courant du Nord-Atlantique, Courant des Canaries, Courant équatorial Nord, Drift Nord-Atlantique).
\begin{enumerate}
    \item Repérez et nommez les quatre courants cités ci-dessus en indiquant leur sens (flèches).
    \item Classez ces courants en « courants chauds » et « courants froids » par rapport aux latitudes qu'ils traversent.
    \item Expliquez l'influence du Gulf Stream et du Drift Nord-Atlantique sur le climat de l'Europe de l'Ouest (températures, précipitations) en les comparant à des lieux situés à même latitude sur la côte Est de l'Amérique du Nord (ex. : Québec vs Nantes).
    \item Pourquoi le courant des Canaries favorise-t-il l'aridité du Sahara occidental et des îles Canaries ?
\end{enumerate}
\end{exercice}

\begin{exercice}[Mécanisme : Formation d'un cyclone tropical]
À partir du schéma ci-dessous représentant une coupe transversale d'un cyclone tropical mature, répondez aux questions.
\begin{popfigure}
\begin{tikzpicture}[scale=0.8]
% Eye
\draw[popblue, thick, fill=popblue!10] (0,0) ellipse (1.5 and 1);
\node[popblue] at (0,0) {\textbf{ŒIL}};
% Eyewall
\draw[poporange, thick, fill=poporange!20] (-3,0) .. controls (-2,4) and (2,4) .. (3,0) .. controls (2,-3) and (-2,-3) .. (-3,0);
\node[poporange] at (0,2.5) {\textbf{Mur de l'œil}};
% Spiral rainbands
\draw[popblue!50, thick] (-6,1) .. controls (-4,2) .. (-2,1.5);
\draw[popblue!50, thick] (6,-1) .. controls (4,-2) .. (2,-1.5);
\draw[popblue!50, thick] (-5,-2) .. controls (-3,-3) .. (-1,-2.5);
\draw[popblue!50, thick] (5,2) .. controls (3,3) .. (1,2.5);
% Arrows
\draw[->, thick, popgreen] (0,-3.5) -- (0,-1.5) node[midway, right] {Air montant chaud/humide};
\draw[->, thick, popred] (0,4.5) -- (0,3) node[midway, right] {Air descendant sec (œil)};
\draw[->, thick, popdark] (-4.5,0) -- (-3.2,0) node[midway, above] {Vents convergents};
\draw[->, thick, popdark] (4.5,0) -- (3.2,0) node[midway, above] {Vents convergents};
\draw[->, thick, popmuted] (-2,4.5) -- (-2,3.5) node[midway, left] {Divergence haute altitude};
\draw[->, thick, popmuted] (2,4.5) -- (2,3.5) node[midway, right] {Divergence haute altitude};
% Ocean
\draw[fill=popblue!30] (-7,-4) rectangle (7,-3);
\node at (0,-3.5) {\textbf{OCÉAN CHAUD ($\ge$ 26,5°C)}};
\end{tikzpicture}
\legende{Coupe schématique d'un cyclone tropical mature (vue méridienne).}
\end{popfigure}
\begin{enumerate}
    \item Nommez les zones A (centre calme), B (anneau de forts vents/pluies), C (bandes spiralées).
    \item Décrivez le circuit thermodynamique (cycle de Carnot) qui alimente le cyclone : évaporation en surface, condensation dans le mur de l'œil, divergence en altitude, subsidence dans l'œil.
    \item Pourquoi la rotation de la Terre (force de Coriolis) est-elle indispensable pour déclencher la rotation cyclonique ? Pourquoi pas de cyclone à l'équateur ?
    \item Que se passe-t-il lorsque le cyclone touche terre (\emph{landfall}) ou arrive sur eau froide ? Justifiez par l'arrêt du « moteur thermique ».
\end{enumerate}
\end{exercice}

\begin{exercice}[Lecture de carte météorologique : Anticyclone et Dépression]
On considère une carte d'isobares (pression au niveau de la mer) sur l'Europe de l'Ouest un jour de janvier. L'isobare 1020 hPa entoure un centre à 1035 hPa (centre A) au large du Portugal. L'isobare 1000 hPa entoure un centre à 990 hPa (centre B) sur le Golfe de Gascogne. Les isobares sont espacées de 5 hPa.
\begin{enumerate}
    \item Identifiez la nature des centres A et B (Anticyclone / Dépression). Justifiez.
    \item Tracez schématiquement le sens des vents en surface autour de A et de B dans l'hémisphère Nord (indiquez la déviation par Coriolis).
    \item Où le gradient de pression (et donc le vent) est-il le plus fort ? Quelle conséquence sur l'état de la mer ?
    \item Prévoyez le temps qu'il fait au centre A et au centre B (ciel, précipitations, vent). Expliquez par les mouvements verticaux d'air (subsidence / convergence-ascendance).
\end{enumerate}
\end{exercice}

\begin{exercice}[Étude de cas : Le phénomène El Niño / La Niña]
Le document ci-après présente des anomalies de température de surface de la mer (SST) dans le Pacifique équatorial pour deux situations.
\textbf{Situation 1 (Normal) :} Eau chaude (\SI{29}{\celsius}) à l'Ouest (Indonésie), eau plus fraîche (\SI{23}{\celsius}) à l'Est (côte péruvienne) avec \emph{upwelling}. Alizés d'Est forts.
\textbf{Situation 2 (El Niño) :} Eau chaude étendue jusqu'à l'Est (\SI{27}{\celsius} au large du Pérou). Alizés affaiblis ou inversés. \emph{Upwelling} stoppé.
\begin{enumerate}
    \item En situation normale, expliquez pourquoi l'eau est plus froide à l'Est (courant de Humboldt, \emph{upwelling}, alizés). Quelle conséquence sur la pêche au large du Pérou ?
    \item Lors d'El Niño, décrivez les modifications de la circulation atmosphérique (cellule de Walker) et océanique. Pourquoi les pluies diluviennes touchent-elles le Pérou/Équateur et la sécheresse l'Australie/Indonésie ?
    \item Citez deux conséquences climatiques mondiales (téléconnexions) d'un fort épisode El Niño (ex. : mousson indienne, cyclones Atlantique, température globale).
\end{enumerate}
\end{exercice}

\courssub{Je me prépare au Probatoire}

\begin{exercice}[Sujet type Probatoire : Cyclone Idai (Mozambique, mars 2019)]
\textbf{Document 1 :} Extrait de rapport OMS/Météo-France. « Le cyclone tropical intense Idai a touché terre près de Beira (Mozambique) le 14 mars 2019 avec des vents moyens de \SI{170}{\kilo\meter\per\hour} et des rafales à \SI{240}{\kilo\meter\per\hour}. Une surcote marine de \SI{6}{\meter} a submergé la ville. Précipitations : \SI{400}{\milli\meter} en 24h. Bilan : > 1000 morts, 3 millions de sinistrés, destruction de 90\% des infrastructures de Beira. »
\textbf{Document 2 :} Carte de trajectoire. Idai s'est formé dans le canal de Mozambique (eau \SI{29}{\celsius}), a fait une boucle, s'est intensifié rapidement avant de toucher terre, puis s'est affaibli en dépression résiduelle apportant des pluies torrentielles sur le Zimbabwe et le Malawi.
\begin{enumerate}
    \item \textbf{Origine dynamique (6 pts) :} En vous appuyant sur les documents et vos connaissances, expliquez les conditions nécessaires à la genèse et à l'intensification rapide d'Idai dans le canal de Mozambique (facteurs océaniques, atmosphériques, géographiques).
    \item \textbf{Phénomènes associés (4 pts) :} Détaillez les trois phénomènes météorologiques dangereux produits par Idai (vents violents, surcote, pluies diluviennes) en reliant chacun à la structure dynamique du cyclone (mur de l'œil, spirales, translation lente).
    \item \textbf{Prévention et protection (5 pts) :} En tant qu'expert en gestion des risques pour la ville de Beira, proposez \textbf{trois} mesures concrètes de prévision/alerte, \textbf{deux} mesures d'aménagement du territoire/urbanisme et \textbf{deux} mesures de préparation des populations qui auraient pu réduire la vulnérabilité. Justifiez chaque mesure.
    \item \textbf{Synthèse (5 pts) :} Rédigez une conclusion argumentée (15 lignes max) répondant à la problématique : « Les mouvements de l'atmosphère et de l'hydrosphère : facteurs de danger et nécessité de prévention ». Votre réponse doit intégrer les échelles spatiale et temporelle, le rôle du changement climatique sur l'intensité des cyclones, et la disparité des moyens de prévention entre pays développés et en développement.
\end{enumerate}
\end{exercice}

\begin{exercice}[Sujet type Probatoire : Synthèse - Circulation générale et climats régionaux]
\textbf{Problématique :} « Comment la circulation générale de l'atmosphère et des océans organise-t-elle les climats régionaux et génère-t-elle des risques météorologiques ? »
On vous donne une carte du monde muette (projection équirectangulaire) avec les parallèles 0°, 30°, 60°, 90°.
\begin{enumerate}
    \item \textbf{Schéma annoté (8 pts) :} Sur la carte fournie (ou reproduite sur votre copie), tracez et légendez :
    \begin{itemize}
        \item Les trois cellules de convection par hémisphère (Hadley, Ferrel, Polaire) avec le sens de la circulation (flèches montantes/descendantes, vents de surface/altitude).
        \item Les zones de haute et basse pression au sol (HP, BP) aux latitudes caractéristiques.
        \item Les principaux vents planétaires (Alizés, Vents d'Ouest, Vents d'Est polaires) avec leur déviation par Coriolis.
        \item Deux courants marins chauds majeurs (ex: Gulf Stream, Kuroshio) et deux courants froids majeurs (ex: Courant du Labrador, Courant de Benguela, Humboldt).
    \end{itemize}
    \item \textbf{Analyse climatique (6 pts) :} En vous appuyant sur votre schéma, expliquez :
    \begin{enumerate}
        \item Pourquoi les grands déserts chauds (Sahara, Atacama, Kalahari, Australie) se situent-ils vers 30° de latitude ?
        \item Pourquoi les façades ouest des continents aux latitudes moyennes (Europe de l'Ouest, Chili, Californie) bénéficient-elles d'un climat océanique doux et humide ?
        \item Pourquoi la côte du Pérou (désert d'Atacama) est-elle l'un des déserts les plus arides du monde malgré sa proximité avec l'océan ?
    \end{enumerate}
    \item \textbf{Risques et adaptation (6 pts) :}
    \begin{enumerate}
        \item Citez deux phénomènes météorologiques violents liés aux instabilités de la circulation atmosphérique (ex: tempêtes extratropicales, tornades, vagues de chaleur bloquantes). Pour l'un d'eux, décrivez le mécanisme dynamique à l'origine.
        \item Comment la surveillance spatiale (satellites météorologiques, altimétrie) et la modélisation numérique permettent-elles la prévention ? Donnez l'exemple d'un système d'alerte précoce efficace (ex: cyclones, crues, tempêtes).
    \end{enumerate}
\end{enumerate}
\end{exercice}

\begin{exercice}[Sujet type Probatoire : Étude intégrée - Le Lac Tchad et la mousson africaine]
\textbf{Contexte :} Le lac Tchad, situé au cœur du Sahel (environ 13° N), a vu sa surface passer de \SI{25000}{\kilo\meter\squared} (1960) à moins de \SI{1500}{\kilo\meter\squared} (années 1980-90) avant une légère reprise récente. Sa variabilité est intimement liée à la mousson ouest-africaine.
\begin{enumerate}
    \item \textbf{Mécanisme de la mousson (5 pts) :} Expliquez le basculement saisonnier de la ZCIT (Zone de Convergence InterTropicale) au-dessus de l'Afrique de l'Ouest. Reliez la position de la ZCIT (janvier vs août), le flux d'alizés (harmattan vs mousson), et le régime des pluies sur le bassin du Lac Tchad.
    \item \textbf{Interactions Océan-Atmosphère (5 pts) :} Comment les anomalies de température de surface de l'Atlantique tropical (gradient Nord/Sud, « Atlantic Niño ») influencent-elles la position de la ZCIT et l'intensité des pluies de mousson au Sahel ? Expliquez les mécanismes de téléconnexion.
    \item \textbf{Retroaction sol-atmosphère (4 pts) :} La diminution de la surface du lac modifie l'évapotranspiration locale. Expliquez comment cette modification de surface (albédo, humidité, rugosité) peut créer une rétroaction négative sur les précipitations locales (effet « lac-champ »).
    \item \textbf{Gestion durable (6 pts) :} Le projet de transfert d'eau du bassin de l'Oubangui (Congo) vers le Lac Tchad (projet Transaqua) est débattu. En tant que géographe/environnementaliste, analysez ce projet sous l'angle :
    \begin{enumerate}
        \item De la faisabilité hydrologique et énergétique (dénivelé, pertes).
        \item Des impacts écologiques potentiels (espèces invasives, modification du régime de crue, delta du Logone).
        \item De la gouvernance transfrontalière (Commission du Bassin du Lac Tchad - CBLT, 6 pays membres).
    \end{enumerate}
    Concluez sur la nécessité d'une gestion intégrée des ressources en eau (GIRE) à l'échelle du bassin face au changement climatique.
\end{enumerate}
\end{exercice}

\newpage

\coursec{Corrigés des exercices}

\begin{corrige}[Exercice 1 — QCM : Circulation atmosphérique générale]
\textbf{Réponses correctes : 2 et 4.}

\begin{enumerate}
    \item \textbf{FAUX.} L'air chaud, moins dense, s'élève bien, mais en s'élevant il \emph{allège} la colonne d'air au-dessus du sol : il crée donc une zone de \textbf{basse pression} (dépression) au sol, et non de haute pression. C'est l'air descendant (subsidence) qui crée les \cle{hautes pressions} (anticyclones).
    \item \textbf{VRAI.} La \cle{cellule de Hadley} s'étend de l'équateur (ascendance à la \cle{ZCIT}) jusqu'à environ 30° de latitude N et S, où l'air redescend (anticyclones subtropicaux).
    \item \textbf{FAUX.} Les \cle{alizés} sont des vents \textbf{d'Est} (NE dans l'hémisphère Nord, SE dans l'hémisphère Sud) qui soufflent des hautes pressions subtropicales (30°) vers l'équateur. Les vents d'Ouest sont ceux des latitudes moyennes (cellule de Ferrel).
    \item \textbf{VRAI.} La \cle{force de Coriolis}, due à la rotation terrestre, dévie tout mouvement vers la \textbf{droite} dans l'hémisphère Nord et vers la \textbf{gauche} dans l'hémisphère Sud. Elle est nulle à l'équateur et maximale aux pôles.
    \item \textbf{FAUX.} La ZCIT est une zone de \textbf{convergence et d'ascendance} (air chaud et humide qui monte, d'où nuages et pluies abondantes). Les zones de subsidence sont les anticyclones subtropicaux vers 30°.
\end{enumerate}
\end{corrige}

\begin{attention}[Points de vigilance — Circulation atmosphérique]
Ne pas confondre « air qui monte » avec « haute pression » : c'est l'inverse. Retenir le couple : \textbf{ascendance $\rightarrow$ basse pression $\rightarrow$ temps perturbé} ; \textbf{subsidence $\rightarrow$ haute pression $\rightarrow$ temps sec}. Attention aussi au sens des alizés : vents \textbf{d'Est} (vers l'équateur), et non d'Ouest.
\end{attention}

\begin{corrige}[Exercice 2 — Vrai ou Faux justifié : Courants océaniques]
\begin{enumerate}
    \item \textbf{FAUX.} Le \cle{Gulf Stream} est un courant \textbf{chaud} qui remonte le long des côtes Est des États-Unis puis traverse l'Atlantique vers l'Europe (il se prolonge par le \cle{courant de dérive Nord-Atlantique}). Il transporte des eaux tropicales vers les hautes latitudes.
    \item \textbf{VRAI.} Le \cle{courant de Benguela} remonte le long des côtes de Namibie et d'Angola ; les alizés écartent les eaux de surface, provoquant une remontée d'eaux profondes froides riches en sels nutritifs (nitrates, phosphates) : c'est l'\cle{upwelling}, d'où une grande richesse halieutique.
    \item \textbf{FAUX.} Les courants de \textbf{surface} (courants de dérive, 0--400 m) sont principalement entraînés par les \textbf{vents} planétaires (alizés, vents d'Ouest) : ce sont des courants de dérive. La circulation \emph{thermohaline} (différences de température et de salinité, donc de densité) concerne surtout les courants de \textbf{profondeur}.
    \item \textbf{VRAI.} \cle{El Niño} est un réchauffement anormal (anomalie positive de SST) des eaux de surface du Pacifique équatorial central et oriental, lié à l'affaiblissement des alizés et à l'arrêt de l'\emph{upwelling} péruvien.
    \item \textbf{FAUX.} Le « tapis roulant » océanique (circulation thermohaline globale) fonctionne sur une échelle de \textbf{plusieurs siècles à un millénaire} (environ 1000 à 1600 ans pour un cycle complet), et non quelques mois.
\end{enumerate}
\end{corrige}

\begin{attention}[Points de vigilance — Courants océaniques]
Distinction fondamentale : \textbf{courants de surface} = moteurs éoliens (vents), rapides (\SIrange{1}{3}{\kilo\meter\per\hour}) ; \textbf{courants de profondeur} = moteurs thermohalins (densité), lents (\SI{}{\centi\meter\per\second}), échelle de temps séculaire. Ne pas inverser les températures : Gulf Stream = chaud, Benguela/Humboldt/Labrador = froids.
\end{attention}

\begin{corrige}[Exercice 3 — Définitions et vocabulaire]
\begin{enumerate}
    \item \textbf{\cle{Cellule de convection atmosphérique} :} circuit fermé de circulation de l'air à l'échelle planétaire, comprenant une branche ascendante, un déplacement en altitude, une branche descendante et un retour en surface. Les trois cellules par hémisphère sont : la \textbf{cellule de Hadley} (0°–30°, directe), la \textbf{cellule de Ferrel} (30°–60°, \emph{indirecte}) et la \textbf{cellule polaire} (60°–90°, directe).
    \item \textbf{\cle{Effet de Coriolis} :} force apparente due à la rotation de la Terre, qui dévie la trajectoire de tout corps en mouvement (vent, courant marin) vers la \textbf{droite} dans l'hémisphère Nord et vers la \textbf{gauche} dans l'hémisphère Sud. Elle est nulle à l'équateur et maximale aux pôles.
    \item \textbf{\cle{Courant de dérive} (surface) :} courant des 100 à 400 premiers mètres, entraîné par le frottement des vents dominants (alizés, vents d'Ouest) ; vitesse de l'ordre de quelques \SI{}{\kilo\meter\per\hour}. \textbf{Courant de profondeur :} courant lent des grandes profondeurs, mis en mouvement par les différences de densité (température et salinité) : circulation \emph{thermohaline}, d'échelle séculaire.
    \item \textbf{\cle{Cyclone tropical} :} système dépressionnaire intense, organisé autour d'un œil calme, se formant sur les océans tropicaux. Conditions minimales : eau de surface $\geq$ \SI{26,5}{\celsius} sur au moins \SI{50}{\meter} d'épaisseur ; latitude $>$ 5° (force de Coriolis suffisante) ; faible cisaillement vertical du vent ; forte humidité en altitude et instabilité de l'air.
    \item \textbf{\cle{Surcote marine} (\emph{storm surge}) :} élévation anormale du niveau de la mer provoquée par le vent violent qui « pousse » l'eau vers la côte (effet de vent) et par l'effet d'aspiration de la très basse pression centrale (effet barométrique inverse : environ \SI{1}{\centi\meter} de surélévation par hPa de déficit). Elle provoque des submersions côtières catastrophiques, aggravées à marée haute.
\end{enumerate}
\end{corrige}

\begin{attention}[Points de vigilance — Vocabulaire]
Cellule de Ferrel = \emph{indirecte} (rotation inverse de Hadley, maintenue par les cellules voisines). Cyclone : l'œil est zone de \textbf{subsidence} (air descendant), le mur de l'œil zone d'\textbf{ascendance} violente. Surcote = effet vent + effet pression ; ne pas oublier l'aggravation à marée haute.
\end{attention}

\begin{corrige}[Exercice 4 — Calcul direct : Vitesse et trajet d'un courant]
\begin{enumerate}
    \item Conversion de la vitesse :
    \begin{align*}
    v &= \SI{1,5}{\meter\per\second} = 1,5 \times 3,6 = \SI{5,4}{\kilo\meter\per\hour}\\
    v &= 5,4 \times 24 = \SI{129,6}{\kilo\meter\per\day} \approx \SI{130}{\kilo\meter\per\day}
    \end{align*}
    \item Vitesse moyenne effective sur la traversée :
    \begin{align*}
    t &= 60 \text{ jours} \times 24 \times 3600 = \SI{5\,184\,000}{\second}\\
    d &= \SI{5000}{\kilo\meter} = \SI{5e8}{\centi\meter}\\
    v_{\text{moy}} &= \frac{d}{t} = \frac{\num{5e8}}{\num{5,184e6}} \approx \SI{96,5}{\centi\meter\per\second} \approx \SI{0,96}{\meter\per\second}
    \end{align*}
    \textbf{Interprétation :} la vitesse effective ($\approx$ \SI{1}{\meter\per\second}) est inférieure à la vitesse maximale du Gulf Stream (\SI{1,5}{\meter\per\second}, pouvant atteindre \SI{2,5}{\meter\per\second} dans son cœur) car : le trajet n'est pas rectiligne (méandres), le courant ralentit en s'élargissant et en se divisant (courant Nord-Atlantique plus diffus), et la vitesse décroît du centre vers les bords du courant. La valeur \SI{1,5}{\meter\per\second} correspond au cœur du courant au large des côtes américaines, non à l'ensemble du trajet.
\end{enumerate}
\end{corrige}

\begin{attention}[Points de vigilance — Calculs de vitesse]
Vérifier systématiquement les conversions : \SI{1}{\meter\per\second} = \SI{3,6}{\kilo\meter\per\hour}. Pour les \SI{}{\centi\meter\per\second}, convertir la distance en cm \emph{et} le temps en secondes avant le calcul. Attention aux puissances de 10 : \SI{5000}{\kilo\meter} = \num{5e8} cm.
\end{attention}

\begin{corrige}[Exercice 5 — Analyse de carte : L'Atlantique Nord]
\begin{enumerate}
    \item \textbf{Repérage :} le \textbf{Courant équatorial Nord} s'écoule d'Est en Ouest le long de 10–20° N (poussé par les alizés) ; le \textbf{Gulf Stream} remonte le long de la Floride puis s'écarte vers le large au niveau du cap Hatteras (direction NE) ; le \textbf{Courant de dérive Nord-Atlantique} traverse l'Atlantique vers l'Europe du Nord-Ouest ; le \textbf{Courant des Canaries} redescend le long des côtes du Maroc et des îles Canaries vers le Sud. L'ensemble forme un \cle{gyre subtropical} tournant dans le \textbf{sens horaire} dans l'hémisphère Nord.
    \item \textbf{Classement :}
    \begin{itemize}
        \item Courants \textbf{chauds} (transportant de l'eau tropicale vers les hautes latitudes) : Courant équatorial Nord, Gulf Stream, Drift Nord-Atlantique.
        \item Courant \textbf{froid} (redescendant des latitudes moyennes vers l'équateur) : Courant des Canaries.
    \end{itemize}
    \item \textbf{Influence sur le climat européen :} le Gulf Stream puis le Drift Nord-Atlantique transportent d'énormes quantités de chaleur vers les côtes d'Europe de l'Ouest. Les \cle{vents d'Ouest} dominants (cellule de Ferrel) advectent cette douceur océanique sur le continent : hivers doux, étés tempérés, précipitations régulières. À latitude comparable, \textbf{Nantes} (47° N) connaît des températures moyennes de janvier autour de \SI{6}{\celsius}, tandis que \textbf{Québec} (47° N, côte Est canadienne, sous influence du courant froid du Labrador et des vents continentaux) descend vers \SI{-12}{\celsius}. L'écart (15 à \SI{20}{\celsius} en hiver) illustre le rôle thermorégulateur des courants chauds et de l'advection océanique.
    \item \textbf{Courant des Canaries et aridité :} ce courant froid refroidit les basses couches de l'atmosphère au-dessus de lui ; l'air refroidi devient \textbf{stable} (peu de convection), l'évaporation est réduite et l'humidité ne peut se condenser en pluies. De plus, la subsidence de la cellule de Hadley (anticyclone subtropical vers 30° N) renforce la sécheresse. Résultat : aridité du Sahara occidental et des îles Canaries malgré la proximité de l'océan.
\end{enumerate}
\end{corrige}

\begin{attention}[Points de vigilance — Analyse de carte]
Gyres : sens horaire (NH), antihoraire (SH) — cohérent avec Coriolis. Courants de bord ouest (Gulf Stream, Kuroshio) = étroits, rapides, profonds, chauds. Courants de bord est (Canaries, Californie, Benguela) = larges, lents, froids, upwelling. Climat façade ouest (Europe) = océanique doux ; façade est (Afrique) = désertique.
\end{attention}

\begin{corrige}[Exercice 6 — Mécanisme : Formation d'un cyclone tropical]
\begin{enumerate}
    \item \textbf{Zones :} A = l'\textbf{œil} (centre calme, air descendant, ciel dégagé, pression minimale) ; B = le \textbf{mur de l'œil} (anneau de convection violente : vents maximaux, pluies torrentielles) ; C = les \textbf{bandes spiralées} (bandes pluvio-orageuses externes en spirale).
    \item \textbf{Circuit thermodynamique (machine de Carnot) :}
    \begin{enumerate}
        \item Au-dessus de l'océan chaud ($\geq$ \SI{26,5}{\celsius}), forte \textbf{évaporation} : l'air de surface se charge en chaleur sensible et en chaleur latente (vapeur d'eau).
        \item Cet air chaud et humide converge vers le centre et \textbf{s'élève} dans le mur de l'œil ; en montant, il se refroidit, la vapeur d'eau \textbf{condense} en libérant la chaleur latente : c'est cette chaleur qui réchauffe la colonne d'air, entretient l'ascendance et fait chuter la pression centrale.
        \item En altitude (tropopause), l'air \textbf{diverge} vers l'extérieur, évacuant la masse d'air.
        \item Une partie de l'air \textbf{redescend dans l'œil} (subsidence), s'y réchauffe par compression adiabatique : d'où le calme et le ciel clair au centre.
    \end{enumerate}
    Le cyclone fonctionne donc comme une \textbf{machine thermique} : source chaude = océan tropical, source froide = haute troposphère.
    \item \textbf{Rôle de Coriolis :} la convergence de l'air vers la dépression ne peut créer de rotation organisée que si la force de Coriolis dévie les trajectoires (vers la droite dans l'hémisphère Nord) : l'air « rate » le centre et se met à tourner (sens \textbf{antihoraire} dans l'hémisphère Nord, \textbf{horaire} dans l'hémisphère Sud). À l'équateur, la force de Coriolis est \textbf{nulle} : l'air converge directement vers le centre sans rotation ; aucun cyclone ne peut donc s'y former ni y persister (d'où une bande « interdite » d'environ 5° de part et d'autre de l'équateur).
    \item \textbf{Landfall ou eau froide :} le « moteur thermique » du cyclone est l'apport de chaleur et d'humidité océaniques. Sur terre, cet apport cesse (et le frottement augmente) ; sur eau froide ($<$ \SI{26}{\celsius}), l'évaporation s'effondre. Sans chaleur latente, la convection s'arrête, la dépression se comble : le cyclone \textbf{s'atténue rapidement} en tempête puis en dépression résiduelle (qui peut néanmoins encore produire des pluies diluviennes, comme Idai sur le Zimbabwe).
\end{enumerate}
\end{corrige}

\begin{attention}[Points de vigilance — Cyclone tropical]
Ne pas écrire que l'œil est une zone d'ascendance : c'est une zone de \textbf{subsidence}. Ne pas dire que Coriolis « crée » le cyclone : il organise la rotation, mais l'énergie vient de l'océan chaud (chaleur latente). Seuil \SI{26,5}{\celsius} sur \SI{50}{\meter} et latitude $>$ 5° sont des conditions \textbf{nécessaires} (pas suffisantes seules).
\end{attention}

\begin{corrige}[Exercice 7 — Lecture de carte météorologique : Anticyclone et Dépression]
\begin{enumerate}
    \item \textbf{Centre A} (\SI{1035}{\hecto\pascal}, pression croissante vers le centre) : \textbf{anticyclone}. \textbf{Centre B} (\SI{990}{\hecto\pascal}, pression décroissante vers le centre) : \textbf{dépression}. Justification : un anticyclone est un maximum de pression, une dépression un minimum.
    \item \textbf{Sens des vents (hémisphère Nord) :} le vent est d'abord dirigé du centre des hautes pressions vers les basses (force du gradient de pression), puis dévié vers la \textbf{droite} par Coriolis jusqu'à circuler presque parallèlement aux isobares. Autour de \textbf{A} : rotation \textbf{horaire} (divergente, vers l'extérieur). Autour de \textbf{B} : rotation \textbf{antihoraire} (convergente, vers l'intérieur).
    \item \textbf{Gradient de pression :} il est maximal là où les isobares sont \textbf{les plus resserrées} ; ici autour de B (\SI{990}{\hecto\pascal} entouré de \SI{1000}{\hecto\pascal}, soit 10 hPa sur une courte distance) le gradient est plus fort qu'autour de A. Conséquence : \textbf{vents forts} près de B, donc mer agitée à forte (forte houle), dangereuse pour la navigation dans le Golfe de Gascogne ; mer calme près de A.
    \item \textbf{Temps prévu :}
    \begin{itemize}
        \item \textbf{Centre A (anticyclone) :} air \textbf{descendant} (subsidence) qui se réchauffe par compression adiabatique et évapore les nuages $\rightarrow$ ciel dégagé, temps sec, vent faible (gel possible la nuit en janvier).
        \item \textbf{Centre B (dépression) :} air \textbf{convergent puis ascendant}, qui se refroidit et condense $\rightarrow$ ciel couvert, nuages épais, \textbf{précipitations}, vent fort.
    \end{itemize}
\end{enumerate}
\end{corrige}

\begin{attention}[Points de vigilance — Carte météo]
Vent $\parallel$ isobares (équilibre gradient/Coriolis). Sens : A (HP) = horaire/divergent ; B (BP) = antihoraire/convergent. Gradient fort = isobares serrées = vent fort. Subsidence (HP) = temps sec ; Ascendance (BP) = temps perturbé.
\end{attention}

\begin{corrige}[Exercice 8 — Étude de cas : El Niño / La Niña]
\begin{enumerate}
    \item \textbf{Situation normale (La Niña / neutre) :} les alizés d'Est poussent les eaux chaudes de surface vers l'Ouest (accumulation près de l'Indonésie, « réservoir d'eau chaude », niveau de la mer plus haut à l'Ouest). Le long des côtes du Pérou, le courant froid de \textbf{Humboldt} remonte vers le Nord et l'évacuation des eaux de surface par les alizés provoque un \textbf{\emph{upwelling}} : remontée d'eaux profondes froides (\SI{23}{\celsius}) riches en nutriments. Conséquence : plancton abondant $\rightarrow$ pêche extrêmement riche (anchois, sardines), première ressource halieutique du Pérou. La \cle{cellule de Walker} (ascendance à l'Ouest, subsidence à l'Est) est bien établie.
    \item \textbf{El Niño :} les alizés s'affaiblissent ou s'inversent (vents d'Ouest anormaux) ; l'eau chaude « reflue » vers l'Est (\SI{27}{\celsius} au large du Pérou), l'\emph{upwelling} s'arrête. La \textbf{cellule de Walker} bascule : l'ascendance se déplace vers le Pacifique central et oriental $\rightarrow$ \textbf{pluies diluviennes et inondations} sur les côtes désertiques du Pérou et de l'Équateur ; la subsidence s'installe sur l'Indonésie et l'Australie $\rightarrow$ \textbf{sécheresses et incendies}. Sans nutriments, la pêche péruvienne s'effondre.
    \item \textbf{Téléconnexions mondiales (exemples) :}
    \begin{itemize}
        \item Affaiblissement de la \textbf{mousson indienne} (déficit pluviométrique en Inde) ;
        \item Réduction de l'activité \textbf{cyclonique dans l'Atlantique Nord} (cisaillement vertical accru en altitude) mais augmentation dans le Pacifique central/oriental ;
        \item Hausse temporaire de la \textbf{température moyenne globale} (libération de chaleur océanique stockée) ;
        \item Perturbation des pluies en Afrique de l'Est et australe (sécheresses au Sahel oriental, pluies anormales en Afrique de l'Est).
    \end{itemize}
\end{enumerate}
\end{corrige}

\begin{attention}[Points de vigilance — El Niño / La Niña]
El Niño = réchauffement Pacifique Est + affaiblissement alizés + arrêt upwelling + bascule Walker. La Niña = situation normale renforcée (eaux plus froides à l'Est, alizés plus forts). Téléconnexions : penser « atmosphère globale couplée » — ce qui se passe dans le Pacifique modifie les cellules de Hadley/Walker et les jets streams planétaires.
\end{attention}

\begin{corrige}[Exercice 9 — Sujet type Probatoire : Cyclone Idai (Mozambique, 2019)]
\textbf{Barème indicatif : 20 points.}

\begin{enumerate}
    \item \textbf{Origine dynamique (6 pts).} Conditions réunies dans le canal de Mozambique :
    \begin{itemize}
        \item \emph{Facteur océanique (2 pts) :} eau très chaude (\SI{29}{\celsius}, $>$ \SI{26,5}{\celsius}) sur une couche épaisse $\rightarrow$ évaporation intense, apport massif de chaleur latente ; faible contenu thermique limitant le refroidissement par brassage.
        \item \emph{Facteurs atmosphériques (2 pts) :} dépression préexistante, convergence des alizés, faible cisaillement vertical du vent (les courants ascendants restent verticaux), forte humidité en moyenne troposphère, bonne divergence en altitude.
        \item \emph{Facteur géographique (2 pts) :} latitude d'environ 15–20° S $\rightarrow$ force de Coriolis suffisante pour organiser la rotation (sens \textbf{horaire} dans l'hémisphère Sud) ; canal de Mozambique = couloir maritime chaud et confiné favorisant l'intensification rapide avant le \emph{landfall}.
    \end{itemize}
    \item \textbf{Phénomènes associés (4 pts).}
    \begin{itemize}
        \item \emph{Vents violents (170–\SI{240}{\kilo\meter\per\hour}) :} concentrés dans le \textbf{mur de l'œil}, là où le gradient de pression est maximal (1,5 pt) ;
        \item \emph{Surcote de \SI{6}{\meter} :} vent poussant l'eau vers la côte + aspiration par la très basse pression centrale, amplifiée par la faible profondeur du plateau continental devant Beira (ville en estuaire, quasi au niveau de la mer) (1,5 pt) ;
        \item \emph{Pluies diluviennes (\SI{400}{\milli\meter}/24 h) :} convection extrême des bandes spiralées + \textbf{translation lente} du système après le \emph{landfall} (dépression résiduelle quasi stationnaire sur le Zimbabwe et le Malawi) $\rightarrow$ crues et glissements de terrain (1 pt).
    \end{itemize}
    \item \textbf{Prévention et protection (5 pts).}
    \begin{itemize}
        \item \emph{Prévision/alerte (3 mesures, 1,5 pt) :} surveillance satellitaire et radars couplée à la modélisation numérique de trajectoire ; système d'alerte précoce multicanal (radio, SMS, sirènes) en langues locales ; exercices d'évacuation et plans d'urgence communaux testés chaque saison cyclonique.
        \item \emph{Aménagement (2 mesures, 2 pts) :} digues et surélévation des défenses côtières + restauration des \textbf{mangroves} (écran naturel dissipant la houle) ; interdiction de construire dans les zones inondables et normes parasismiques/paracycloniques (toits ancrés, bâtiments-refuges surélevés).
        \item \emph{Préparation des populations (2 mesures, 1,5 pt) :} éducation aux risques (réflexes, itinéraires d'évacuation, kits d'urgence) ; organisation communautaire de secours (stocks d'eau potable, vivres, postes de santé de fortune).
    \end{itemize}
    \item \textbf{Synthèse (5 pts) — éléments attendus :} les mouvements de l'atmosphère et de l'hydrosphère redistribuent l'énergie solaire à l'échelle planétaire (cellules, courants) mais concentrent localement une énergie considérable (cyclones, surcotes) à des échelles spatiales (centaines de km) et temporelles (quelques jours) qui en font des catastrophes. Le changement climatique, en réchauffant les océans, augmente l'\emph{intensité} potentielle des cyclones et la quantité de pluie (air plus chaud = plus de vapeur d'eau), ainsi que la surcote via l'élévation du niveau marin. Or les moyens de prévention (satellites, alerte, infrastructures) sont très inégalement répartis : les pays en développement comme le Mozambique subissent les pertes humaines les plus lourdes pour des phénomènes d'intensité comparable. La prévention (prévision, aménagement, éducation, solidarité internationale) est donc une nécessité vitale et un enjeu de justice climatique.
\end{enumerate}
\end{corrige}

\begin{attention}[Points de vigilance — Sujet Probatoire Cyclone]
Exigences du correcteur : citer explicitement le seuil de \SI{26,5}{\celsius}, le rôle de Coriolis (latitude $>$ 5°), le faible cisaillement ; relier chaque dégât à une structure du cyclone (mur de l'œil, bandes, surcote) ; en synthèse, mobiliser les deux échelles (spatiale et temporelle) et la dimension Nord/Sud (inégalité face au risque).
\end{attention}

\begin{corrige}[Exercice 10 — Sujet type Probatoire : Circulation générale et climats régionaux]
\textbf{Barème indicatif : 20 points.}

\begin{enumerate}
    \item \textbf{Schéma annoté (8 pts).} Éléments attendus sur la carte :
    \begin{itemize}
        \item \emph{Cellules (3 pts) :} Hadley (0° $\leftrightarrow$ 30° : ascendance à l'équateur, subsidence à 30°) ; Ferrel (30° $\leftrightarrow$ 60°, cellule \emph{indirecte}) ; polaire (60° $\leftrightarrow$ 90°), dans chaque hémisphère, avec flèches verticales et horizontales.
        \item \emph{Pressions (2 pts) :} BP équatoriales (ZCIT), HP subtropicales ($\approx$ 30°), BP subpolaires ($\approx$ 60°), HP polaires.
        \item \emph{Vents (2 pts) :} alizés de NE (hém. N) et de SE (hém. S) entre 30° et l'équateur ; vents d'Ouest entre 30° et 60° ; vents d'Est polaires ; déviations cohérentes avec Coriolis (droite NH, gauche SH).
        \item \emph{Courants (1 pt) :} ex. Gulf Stream et Kuroshio (chauds, flèches rouges) ; Labrador et Benguela/Humboldt (froids, flèches bleues).
    \end{itemize}
    \item \textbf{Analyse climatique (6 pts).}
    \begin{enumerate}
        \item \emph{Déserts vers 30° (2 pts) :} sous les HP subtropicales, l'air de la cellule de Hadley \textbf{descend} (subsidence), se réchauffe par compression adiabatique : son humidité relative chute, les nuages se dissipent $\rightarrow$ climats arides (Sahara, Kalahari, désert australien).
        \item \emph{Façades ouest des continents aux latitudes moyennes (2 pts) :} exposées aux \textbf{vents d'Ouest} dominants qui advectent de l'air maritime doux et humide ; présence de courants chauds (Drift Nord-Atlantique, Kuroshio) et passage fréquent des perturbations du front polaire $\rightarrow$ climat océanique doux et pluvieux toute l'année.
        \item \emph{Désert d'Atacama (2 pts) :} triple cause : subsidence de l'anticyclone subtropical du Pacifique Sud ; courant froid de \textbf{Humboldt} qui stabilise l'air et bloque la convection ; effet d'\textbf{abri orographique} (effet de foehn inverse) de la cordillère des Andes qui bloque les flux humides d'Est. D'où une aridité extrême au bord de l'océan.
    \end{enumerate}
    \item \textbf{Risques et adaptation (6 pts).}
    \begin{enumerate}
        \item \emph{Phénomènes (3 pts) :} tempêtes extratropicales (cyclogenèse sur le front polaire : contraste thermique entre air polaire et air subtropical, creusement dépressionnaire rapide, vents violents en hiver sur l'Europe) ; tornades (convection violente + fort cisaillement dans les supercellules des plaines américaines) ; dômes de chaleur (blocage anticyclonique persistant). Mécanisme détaillé exigé pour l'un d'eux (ex. tempête : creusement explosif par divergence en altitude + advection d'air chaud/humide).
        \item \emph{Surveillance et prévention (3 pts) :} satellites géostationnaires (imagerie continue) et défilants, altimétrie satellitaire (niveau et température de la mer), bouées et radars $\rightarrow$ données assimilées par les modèles numériques de prévision. Exemples : le système d'alerte aux cyclones de l'OMM/centres régionaux (RSMC La Réunion pour l'océan Indien Sud-Ouest), les alertes tempête de Météo-France (vigilance), les systèmes d'alerte aux crues : anticipation de 3 à 5 jours permettant évacuations et mise en sécurité.
    \end{enumerate}
\end{enumerate}
\end{corrige}

\begin{attention}[Points de vigilance — Circulation générale \& Climats]
Sur le schéma : la cellule de Ferrel est \emph{indirecte} (sens de rotation inverse de Hadley) ; ne pas faire souffler les alizés d'Ouest ; vérifier la cohérence des déviations de Coriolis dans chaque hémisphère (droite NH, gauche SH). Pour Atacama : citer les 3 facteurs (subsidence, courant froid, abri orographique).
\end{attention}

\begin{corrige}[Exercice 11 — Sujet type Probatoire : Le Lac Tchad et la mousson africaine]
\textbf{Barème indicatif : 20 points.}

\begin{enumerate}
    \item \textbf{Mécanisme de la mousson (5 pts).} La \cle{ZCIT} suit (avec retard) le zénith apparent du Soleil, attirée par les zones de réchauffement maximal. En \textbf{janvier}, la ZCIT est au Sud (environ 5° S) : le Sahel est sous l'influence de l'\textbf{harmattan}, alizé de NE chaud et sec venu du Sahara $\rightarrow$ saison sèche sur le bassin du Lac Tchad. En \textbf{août}, la ZCIT remonte jusqu'à 18–20° N : le flux de \textbf{mousson} (alizé de SE de l'hémisphère Sud dévié en vent de SO en traversant l'équateur, chargé d'humidité atlantique) pénètre loin au Nord $\rightarrow$ saison des pluies sur le bassin tchadien. Le niveau du lac dépend donc directement de l'amplitude et de la durée de cette remontée de la ZCIT.
    \item \textbf{Interactions Océan-Atmosphère (5 pts).} La position de la ZCIT est contrôlée par le \textbf{gradient de température de surface (SST) entre Atlantique tropical Nord et Sud}. Si l'Atlantique Nord tropical est anormalement chaud (ou le Sud anormalement froid), le contraste thermique attire la ZCIT plus au Nord $\rightarrow$ mousson renforcée, pluies abondantes au Sahel. À l'inverse, un « \textbf{Atlantic Niño} » (eaux chaudes anormales dans le golfe de Guinée, au Sud) maintient la ZCIT basse $\rightarrow$ déficit pluviométrique sahélien (mécanisme impliqué dans les grandes sécheresses des années 1970–80). Téléconnexions : El Niño Pacifique perturbe aussi la circulation de Walker à l'échelle globale et tend à réduire la mousson ouest-africaine.
    \item \textbf{Rétroaction sol-atmosphère (4 pts).} Le retrait du lac remplace une surface d'eau (albédo faible, forte évaporation) par des sols nus ou des sables (albédo plus élevé, évapotranspiration réduite). Conséquences : moins de vapeur d'eau locale recyclée dans l'atmosphère $\rightarrow$ moins de convection et de pluies locales ; échauffement diurne plus fort mais sans humidité disponible. Cette rétroaction \textbf{positive} (au sens d'amplification du phénomène initial : la sécheresse) — appelée ici effet « lac-champ » — aggrave localement le déficit pluviométrique : moins d'eau $\rightarrow$ moins d'évaporation $\rightarrow$ moins de pluie $\rightarrow$ encore moins d'eau. (Le terme « rétroaction négative sur les précipitations » signifie : effet défavorable aux pluies.)
    \item \textbf{Gestion durable (6 pts).}
    \begin{enumerate}
        \item \emph{Faisabilité (2 pts) :} le transfert Oubangui $\rightarrow$ Chari suppose de franchir un seuil topographique (dénivelé de plusieurs centaines de mètres) nécessitant pompage ou ouvrages gravitaires complexes, sur plus de \SI{2000}{\kilo\meter}, avec des pertes par évaporation et infiltration importantes sous climat sahélien ; coût énergétique et financier considérable.
        \item \emph{Impacts écologiques (2 pts) :} risque de transfert d'espèces invasives entre bassins (Congo $\neq$ Tchad), modification du régime naturel des crues du Chari et du delta du Logone (inondations saisonnières vitales pour les pâturages et la pêche), prélèvement pouvant affecter les écosystèmes du bassin du Congo.
        \item \emph{Gouvernance (2 pts) :} le bassin est partagé entre 6 États (CBLT : Cameroun, Tchad, Niger, Nigéria, Centrafrique, Libye) $\rightarrow$ nécessité d'accords de répartition, de financement et de suivi environnemental communs ; risques de tensions si les bénéfices sont perçus comme inégaux.
    \end{enumerate}
    \textbf{Conclusion :} face au changement climatique et à la variabilité de la mousson, seule une \textbf{GIRE} (gestion intégrée des ressources en eau) à l'échelle du bassin — irrigation raisonnée, reboisement, recharge des nappes, concertation transfrontalière — peut durablement concilier sauvegarde du lac, sécurité alimentaire et paix entre les usagers.
\end{enumerate}
\end{corrige}

\begin{attention}[Points de vigilance — Lac Tchad \& Mousson]
Bien distinguer \cle{harmattan} (sec, NE, hiver boréal) et \cle{mousson} (humide, SO, été boréal) ; ne pas attribuer le retrait du lac uniquement au climat : les prélèvements agricoles (irrigation) y contribuent fortement ; pour la rétroaction, expliciter la chaîne causale complète (surface $\rightarrow$ albédo/évaporation $\rightarrow$ humidité atmosphérique $\rightarrow$ convection $\rightarrow$ pluie). GIRE = échelle bassin, concertation, adaptation.
\end{attention}

\newpage

\coursec{Fiche bilan}

\courssub{Carte mentale de synthèse}
\begin{popfigure}
\begin{tikzpicture}[
    mindmap,
    grow cyclic,
    every node/.style={concept, execute at begin node=\hskip0pt},
    concept color=popblue,
    level 1/.append style={level distance=4.5cm, sibling angle=60},
    level 2/.append style={level distance=3.2cm, sibling angle=45},
    text=white,
    font=\small\bfseries
]
\node[concept, scale=1.1, align=center]{Grands mouvements\\ atmosphère \& hydrosphère}
    child[concept color=poporange] {node[concept] {1. Moteurs}
        child {node[concept, scale=0.7] {Rayonnement solaire inégal}}
        child {node[concept, scale=0.7] {Rotation Terre (Coriolis)}}
        child {node[concept, scale=0.7] {Différences de pression}}}
    child[concept color=popgreen] {node[concept] {2. Circulation atmosphérique}
        child {node[concept, scale=0.7] {Cellules : Hadley, Ferrel, Polaire}}
        child {node[concept, scale=0.7] {Vents planétaires : Alizés, Ouest, Polaires}}
        child {node[concept, scale=0.7] {ZCIT : convergence intertropicale}}}
    child[concept color=poppurple] {node[concept] {3. Circulation océanique}
        child {node[concept, scale=0.7] {Courants de surface (vents)}}
        child {node[concept, scale=0.7] {Courants profonds (thermohaline)}}
        child {node[concept, scale=0.7] {Gyres subtropicaux}}}
    child[concept color=poppink] {node[concept] {4. Couplage Océan-Atmosphère}
        child {node[concept, scale=0.7] {Échanges chaleur/humidité}}
        child {node[concept, scale=0.7] {El Niño / La Niña (ENSO)}}
        child {node[concept, scale=0.7] {Impact pluviométrie (Sahel, Cameroun)}}}
    child[concept color=popdark] {node[concept] {5. Phénomènes dangereux}
        child {node[concept, scale=0.7] {Cyclones tropicaux (structure, échelle Saffir-Simpson)}}
        child {node[concept, scale=0.7] {Tempêtes extratropicales}}
        child {node[concept, scale=0.7] {Raz-de-marée, houle cyclonique}}}
    child[concept color=popgold] {node[concept] {6. Prévention \& Protection}
        child {node[concept, scale=0.7] {Surveillance météo \& satellite}}
        child {node[concept, scale=0.7] {Systèmes d'alerte précoce}}
        child {node[concept, scale=0.7] {Aménagement littoral, éducation}}};
\end{tikzpicture}
\legende{Carte mentale récapitulant les 6 piliers de la leçon NCT07 : moteurs énergétiques, circulation atmosphérique (cellules/vents), circulation océanique (courants/gyres), interactions ENSO, phénomènes violents (cyclones/tempêtes) et stratégies de réduction des risques.}
\end{popfigure}

\courssub{Bilan des savoirs et savoir-faire essentiels}
\begin{aretenir}[Bilan des savoirs et savoir-faire essentiels]

\textbf{Définitions clés à connaître par cœur}
\begin{itemize}
    \item \cle{Cellule de convection} : Mouvement circulaire d'une masse fluide (air/eau) entraîné par des différences de densité (chaud $\uparrow$ / froid $\downarrow$). Trois cellules par hémisphère : \textbf{Hadley} (0--30°), \textbf{Ferrel} (30--60°), \textbf{Polaire} (60--90°).
    \item \cle{Force de Coriolis} : Force apparente due à la rotation de la Terre, déviant les fluides vers la \textbf{droite} dans l'hémisphère Nord, vers la \textbf{gauche} dans l'hémisphère Sud. Nulle à l'équateur, maximale aux pôles.
    \item \cle{Alizés} : Vents réguliers soufflant des subtropicales (hautes pressions) vers l'équateur (basses pressions), déviés par Coriolis (NE au Nord, SE au Sud). Convergent à la \cle{ZCIT} (Zone de Convergence Intertropicale).
    \item \cle{Courant de surface} : Entraîné par les vents planétaires (ex: Gulf Stream, Courant des Canaries, Courant de Benguela, Courant de Guinée). \cle{Courant profond (thermohaline)} : Entraîné par différences de densité (température + salinité) : eau froide/salée $\downarrow$ ($NADL$, $AABW$).
    \item \cle{Gyre subtropical} : Boucle de courants fermée (ex: Gyre Nord-Atlantique) : courant chaud occidental (rapide, étroit) + courant froid oriental (lent, large).
    \item \cle{ENSO (El Niño / La Niña)} : Oscillation couplée océan-atmosphère dans le Pacifique équatorial. \textbf{El Niño} : affaiblissement alizés, réchauffement Est Pacifique $\rightarrow$ sécheresses Australie/Indonésie, pluies Am. Sud. \textbf{La Niña} : renforcement alizés, refroidissement Est Pacifique $\rightarrow$ effets inverses.
    \item \cle{Cyclone tropical} : Dépression organisée sur mer chaude ($\ge$ 26,5°C), structure symétrique (œil, mur de l'œil, bandes spiralées), vents $>$ 118 km/h. Échelle \textbf{Saffir-Simpson} (cat. 1 à 5).
\end{itemize}

\textbf{Mécanismes fondamentaux (formules/relations)}
\begin{center}
\begin{tabularx}{\linewidth}{|l|X|}\hline
\textbf{Concept} & \textbf{Relation / Règle à retenir} \\\hline
Moteur thermique & $\Delta T$ (Équateur-Pôles) $\rightarrow$ Gradient de pression $\rightarrow$ Vent (air va de HP vers BP) \\\hline
Force de Coriolis ($F_c$) & $F_c = 2 \Omega v \sin\varphi$ \quad ($\Omega$ vitesse angulaire, $v$ vitesse vent, $\varphi$ latitude). Déviation : \textbf{Droite (N), Gauche (S)}. \\\hline
Équilibre géostrophique & Vent $\parallel$ aux isobares (HP à droite dans hémisphère N). Vitesse $\propto$ gradient de pression. \\\hline
Courants de surface & Direction $\approx$ vent dominant + déviation Coriolis (45° théorique, \textbf{Ekman}). Transport d'Ekman $\perp$ au vent (droite N). \\\hline
Circulation thermohaline & Formation eau profonde : \textbf{$NADL$} (Atlantique Nord, mer du Labrador/Groenland) et \textbf{$AABW$} (Antarctique, mer de Weddell/Ross). \\\hline
Formation cyclone & Condition nécessaire : $T_{\text{mer}} \ge 26,5^{\circ}\text{C}$ sur 50m + cisaillement vertical faible + vorticité initiale (ZCIT) + distance $\ge 5^{\circ}$ équateur (Coriolis $\ne 0$). \\\hline
\end{tabularx}
\end{center}

\textbf{Méthodes express (Savoir-faire Probatoire)}
\begin{enumerate}
    \item \textbf{Lire une carte de circulation générale} : Identifier cellules (ascendance ZCIT, subsidence 30°), nommer vents (Alizés NE/SE, Vents d'Ouest), indiquer déviation Coriolis.
    \item \textbf{Analyser une carte de courants marins} : Différencier courant chaud (flèche rouge, vers pôles, bord ouest) et froid (flèche bleue, vers équateur, bord est). Repérer gyres (sens horaire N, antihoraire S).
    \item \textbf{Expliquer la genèse d'un cyclone} : Étape 1: Évaporation intense mer chaude $\rightarrow$ air chaud/humide $\uparrow$. Étape 2: Condensation $\rightarrow$ chaleur latente libérée $\rightarrow$ renforce ascendance (retroaction +). Étape 3: Coriolis organise rotation $\rightarrow$ œil (subsidence centrale) + mur de l'œil (vents max, pluies diluviennes).
    \item \textbf{Relier ENSO aux impacts climatiques} : El Niño $\rightarrow$ ZCIT déplace vers Est $\rightarrow$ sécheresse Indonésie/Australie, inondations Pérou/Equateur, \textbf{sécheresse au Sahel/Cameroun Nord}. La Niña $\rightarrow$ effets opposés (excès pluies Sahel).
    \item \textbf{Proposer mesures protection} : \textit{Prévision} : satellites, bouées, modèles numériques. \textit{Alerte} : vigilance météo (jaune/orange/rouge), sirènes, SMS. \textit{Aménagement} : digues, mangroves (brise-lames naturel), zones inconstructibles, plans d'évacuation. \textit{Éducation} : exercices simulation, kit survie (eau, radio, papiers).
\end{enumerate}

\end{aretenir}

\courssub{Checklist avant le Probatoire}
\begin{aretenir}[Checklist NCT07 -- Grands mouvements atmosphère/hydrosphère]
$\square$ Je sais schématiser et nommer les \textbf{3 cellules de convection} par hémisphère (Hadley, Ferrel, Polaire) et indiquer le sens des vents aux surfaces (Alizés, Vents d'Ouest, Vents polaires).\\
$\square$ J'explique l'origine de la \textbf{Force de Coriolis} et j'applique la règle de déviation : \textbf{Droite (Nord), Gauche (Sud)} sur les vents et courants.\\
$\square$ Je localise la \textbf{ZCIT} (ascendance, pluies, calmes équatoriaux) et les \textbf{hautes pressions subtropicales} (subsidence, déserts, Alizés).\\
$\square$ Je distingue \textbf{courants de surface} (vent + Coriolis, ex: Gulf Stream, Canaries, Benguela, Guinée) et \textbf{courants profonds} (thermohaline, $NADL$, $AABW$, « tapis roulant » global).\\
$\square$ Je décris un \textbf{gyre subtropical} : courant occidental chaud/rapide/étroit (ex: Gulf Stream) + courant oriental froid/lent/large (ex: Canaries).\\
$\square$ Je caractérise \textbf{El Niño} (affaiblissement alizés, réchauffement Est Pacifique, sécheresses Ouest Pacifique/Sahel) et \textbf{La Niña} (renforcement alizés, refroidissement Est, pluies Sahel/Cameroun).\\
$\square$ J'énonce les \textbf{conditions nécessaires à la cyclogenèse} : $T_{\text{mer}} \ge 26,5^{\circ}\text{C}$, cisaillement faible, vorticité (ZCIT), Coriolis $\ne 0$ ($>5^{\circ}$ lat.).\\
$\square$ Je décris la \textbf{structure d'un cyclone mature} : Œil (calme, subsidence, BP min), Mur de l'œil (vents max, convection intense), Bandes spiralées (orages).\\
$\square$ Je classe un cyclone sur l'\textbf{échelle de Saffir-Simpson} (Cat 1 à 5) à partir du vent moyen maximal (10 min ou 1 min soutenu).\\
$\square$ Je relie \textbf{tempête extratropicale} (front polaire, gradient T air, asymétrique) et \textbf{cyclone tropical} (symétrique, cœur chaud, mer chaude).\\
$\square$ Je cite \textbf{3 mesures de prévention} (surveillance satellite/radar, alerte population, aménagement littoral/mangroves) et \textbf{3 mesures de protection individuelle} (kit survie, respect consignes évacuation, mise en sécurité habitat).
\end{aretenir}

\findecours

\end{document}
